\chapter{微积分初步}
\label{chap:calculus}

第二章研究了函数的单调性、连续性及函数值的范围。现在进一步问：函数在某个位置变化得有多快？能否用简单的函数近似它？由此进入对函数变化的研究。

一辆车在一段时间内的平均速度，是位置变化量与时间的比。考察越来越短的时间内的运动，研究这个比是否趋于确定的值，就能描述瞬时速度。曲线上两点逐渐靠近时，连线的方向也可能趋于确定，给出一点处的切线方向。这两个问题将在导数的定义下获得共同的表达。

确定局部变化率之后，可以用一条直线近似曲线的一小段，研究函数的增减、极值及图像的形态。若需要更精确的近似，就进一步研究变化率本身的变化，用更高次的多项式建立局部模型。近似的精度，则由实际函数与模型之间的误差决定。

反过来，若知道各时刻的速度，怎样得到一段时间内的位移？可以将时间分成小段，用速度乘时间近似各段位移，再把它们相加。曲边图形的面积也可以用许多小矩形近似。研究分割不断加细时这些和的变化，就从有限求和进入连续的整体累积。

局部变化与整体累积可以分别研究。把一个量累积到不同终点，所得函数的变化率与原来的量有什么关系？一个函数的总变化，又能否从各点的变化率中恢复？本章将分别定义导数与定积分，再由微积分基本定理证明它们之间的联系，使这些认识成为计算和推理的依据。

近似也可以从一点附近推进到整个区间。若用一个函数代替另一个函数，怎样同时控制各处的误差？近似函数的积分、导数是否也接近原来的结果？进一步把函数作为整体，把求导与积分看作作用在函数上的运算，就能研究这些问题，并重新理解运算的线性性质、原函数中的常数及基本定理。

在此基础上，还可以反过来提出问题：若一个种群的增长速度由当前数量决定，能否确定数量随时间的变化？这里给定的是变化规律，待求的是符合规律的函数。微分方程将把导数、积分、初始状态与递推近似放在同一个问题中，使我们从研究一个给定函数，进一步走向由规律寻找函数。

第一章的极限、第二章的连续性和第三章的线性思想，将在这些问题中继续发挥作用。差商、有限求和与递推提供可以计算的起点，极限与误差控制使近似获得确定的结论。计算公式和解题方法，将随着对象的建立及其相互联系自然产生。

\FloatBarrier
\newpage
\section{导数}
\label{sec:calculus-derivatives}

\begingroup
\setlength{\parskip}{5pt}

\subsection{平均变化率}

设自变量由\(x_0\)变到\(x_0+h\)。相应的自变量增量与函数增量分别为
\begin{equation*}
 \Delta x=h,\qquad \Delta y=f(x_0+h)-f(x_0).
\end{equation*}
当\(h\ne0\)时，称
\begin{equation}\label{eq:calculus-average-rate}
 \frac{\Delta y}{\Delta x}=\frac{f(x_0+h)-f(x_0)}h
\end{equation}
为这一段上的{\heiti{}平均变化率}。在函数图像上，它是两点\((x_0,f(x_0))\)、\((x_0+h,f(x_0+h))\)所确定割线的斜率；当\(f\)表示位置、\(x\)表示时间时，它就是平均速度。

例如，\(f(x)=x^2\)时，
\begin{equation*}
 \frac{(x_0+h)^2-x_0^2}{h}=2x_0+h.
\end{equation*}
固定\(x_0\)，缩短所考察的范围，平均变化率就趋于\(2x_0\)；改变\(x_0\)，这个数也随之改变。前者是在一个位置附近确定变化率，后者则是比较不同位置处的变化。

\begin{figure}[!ht]
\centering
\includegraphics[width=.43\linewidth]{calculus-secants}
\caption{固定点\(P\)，另一点沿曲线趋近它；割线斜率趋于切线\(\ell\)的斜率。}
\label{fig:calculus-secants}
\end{figure}

\begin{example}
某质点的位置为\(s(t)=t^2+2t\)。分别计算从\(t=1\)到\(t=1+h\)及从\(t=2\)到\(t=2+h\)的平均速度。令\(h\)分别取正数和负数，解释所比较的时间范围。
\end{example}

\subsection{一点处的导数}

\begin{definition}
设\(f\)在\(x_0\)附近有定义。若有限极限
\begin{equation}\label{eq:calculus-derivative-definition}
 f'(x_0)=\lim_{h\to0}\frac{f(x_0+h)-f(x_0)}h
\end{equation}
存在，则称\(f\)在\(x_0\)处{\heiti{}可导}，这个极限称为\(f\)在该点的{\heiti{}导数}。
\end{definition}

令\(x=x_0+h\)，定义也可以写为
\begin{equation*}
 f'(x_0)=\lim_{x\to x_0}\frac{f(x)-f(x_0)}{x-x_0}.
\end{equation*}
这里始终固定\(x_0\)，让另一个点向它靠近。瞬时速度由这样的极限给出；它的确定并不需要令时间增量等于零。

只取\(h\to0^-\)或\(h\to0^+\)，分别得到左导数与右导数。函数在内点可导，当且仅当两者都存在且相等。区间端点处需要导数时，按相应的单侧导数理解。

由定义直接得到，常值函数的导数为零，\(f(x)=kx+b\)在每一点的导数均为\(k\)，而\(f(x)=x^2\)在\(x_0\)处的导数为\(2x_0\)。对于\(f(x)=|x|\)，原点处的差商为\(|h|/h\)，左右极限分别为\(-1\)与\(1\)，所以它在原点不可导。

\begin{example}
用定义求\(f(x)=x^3\)在\(x_0\)处的导数。再设\(g(x)=\sqrt{x}\)，用有理化求\(g'(x_0)\)，其中\(x_0>0\)。
\end{example}

\begin{exercise}
\mbox{}\par\nobreak

（1）用定义求\(f(x)=1/x\)在\(x_0\ne0\)处的导数。

（2）设\(f(x)=x|x|\)。判断它在原点是否可导，并与\(|x|\)比较。

（3）设\(f(0)=0\)，\(x\ne0\)时\(f(x)=x\sin(1/x)\)。判断\(f\)在原点是否连续、可导。
\end{exercise}

\subsection{可导函数的连续性}

可导性研究的是函数增量与自变量增量的比。若这个比趋于有限的数，那么自变量增量趋零时，函数增量本身也必趋零。

\begin{theorem}\label{thm:calculus-differentiable-continuous}
若\(f\)在\(x_0\)处可导，则它在该点连续。
\end{theorem}

\begin{proof}
当\(h\ne0\)时，
\begin{equation*}
 f(x_0+h)-f(x_0)=\frac{f(x_0+h)-f(x_0)}h\,h\longrightarrow f'(x_0)\cdot0=0.
\end{equation*}
因此\(f(x_0+h)\to f(x_0)\)。
\end{proof}

反过来，连续只要求函数增量趋零，并未规定它与\(h\)之间的比例。\(|x|\)在原点连续而不可导，就是一个例子。

比较\(|x|\)、\(x\sin(1/x)\)与\(x^2\sin(1/x)\)，后两者在原点补定义为零。三者都在原点连续，但差商分别为\(|h|/h\)、\(\sin(1/h)\)、\(h\sin(1/h)\)，只有最后一个趋于零。

\subsection{导函数}

\begin{definition}
将\(f\)在各可导点处的导数与该点对应，得到的函数称为\(f\)的{\heiti{}导函数}，记为\(f'\)。其定义域是\(f\)所有可导点组成的集合。
\end{definition}

例如，\(f(x)=x^2\)的导函数为\(f'(x)=2x\)。原函数的图像是一条抛物线，导函数的图像则是一条直线；后者的纵坐标记录前者在相应横坐标处的变化率。\(f'(1)=2\)是一个数，\(f'\)则是一个新的函数。

若记\(y=f(x)\)，也常写
\begin{equation*}
 y'=f'(x)=\frac{dy}{dx},\qquad \frac{d}{dx}f(x)=f'(x).
\end{equation*}
其中\(d/dx\)表示对\(x\)求导。正的导数与函数上升、负的导数与函数下降之间的联系，可以从图像观察；要把这一观察推进为区间上的判据，还需要后面的中值定理。

\begin{example}
求\(f(x)=|x|\)的导函数及其定义域。已知某函数在\((-1,0)\)上导数为负、在\((0,1)\)上导数为正，试画出可能的原函数图像，并考虑\(x=0\)处的取值是否已被这些信息确定。
\end{example}

\subsection{切线}

图像上经过\(P=(x_0,f(x_0))\)与\(Q=(x_0+h,f(x_0+h))\)的割线，可以取方向向量
\begin{equation*}
 \left(1,\frac{f(x_0+h)-f(x_0)}h\right).
\end{equation*}
若\(f\)在\(x_0\)处可导，割线方向趋于\((1,f'(x_0))\)，所以切线方程为
\begin{equation}\label{eq:calculus-tangent}
 y-f(x_0)=f'(x_0)(x-x_0).
\end{equation}
法线的方向向量可取\((-f'(x_0),1)\)，因而其方程可写为
\begin{equation*}
 (x-x_0)+f'(x_0)(y-f(x_0))=0.
\end{equation*}
切线由局部方向确定，不要求图像始终位于直线同一侧。例如，\(y=x^3\)在原点的切线为\(y=0\)，图像却穿过这条切线。

\begin{example}
用导数求抛物线\(y=x^2\)在\((1,1)\)处的切线和法线，并用第四章的方法核对。再求\(y=x^3\)在原点的切线。
\end{example}

\begin{example}[2012年广东卷理科第12题]
\label{ex:calculus-gaokao-polynomial-tangent}
求曲线\(y=x^3-x+3\)在点\((1,3)\)处的切线方程。用导数定义确定切线斜率。
\end{example}

\begin{exercise}
\mbox{}\par\nobreak

（1）（2013年大纲卷文科第10题改编）曲线\(y=x^4+ax^2+1\)在点\((-1,a+2)\)处切线的斜率为\(8\)，求\(a\)。用差商极限完成计算。

（2）（2012年安徽卷文科第17题第（2）问）设\(f(x)=ax+1/(ax)+b\)，\(a>0\)。若曲线\(y=f(x)\)在点\((1,f(1))\)处的切线方程为\(y=3x/2\)，求\(a,b\)。
\end{exercise}

\subsection{微分}
\label{subsec:calculus-differential}

平方函数的增量满足
\begin{equation*}
 (x_0+h)^2-x_0^2=2x_0h+h^2.
\end{equation*}
其中\(2x_0h\)关于\(h\)线性，而余下的\(h^2\)与\(h\)之比趋零。一般可导函数也有这种结构。

\begin{proposition}\label{prop:calculus-first-order}
函数\(f\)在\(x_0\)处可导，当且仅当存在常数\(a\)，使
\begin{equation}\label{eq:calculus-first-order}
 f(x_0+h)=f(x_0)+ah+r(h),\qquad \lim_{h\to0}\frac{r(h)}h=0.
\end{equation}
此时\(a=f'(x_0)\)。
\end{proposition}

\begin{proof}
若可导，取\(a=f'(x_0)\)，将定义中的差商减去\(a\)，就得到\(r(h)/h\to0\)。反过来，将式\eqref{eq:calculus-first-order}减去\(f(x_0)\)并除以\(h\)，差商便趋于\(a\)。
\end{proof}

\begin{definition}
对可导函数\(y=f(x)\)，称\(dx=\Delta x\)为自变量的微分，称函数增量的线性部分
\begin{equation*}
 dy=f'(x)\,dx
\end{equation*}
为函数在\(x\)处的{\heiti{}微分}。
\end{definition}

因此，一次近似为\(f(x+h)\approx f(x)+f'(x)h\)。\(\Delta y\)是实际增量，\(dy\)是线性近似，两者之差与\(\Delta x\)之比趋零。导数记号\(dy/dx\)也就与这个线性关系相容。

\begin{figure}[!ht]
\centering
\includegraphics[width=.47\linewidth]{calculus-differential}
\caption{同一自变量增量下，\(\Delta y\)由曲线确定，\(dy\)由切线确定，两者之差为余项。}
\label{fig:calculus-differential}
\end{figure}

例如，\(x=4\)时，\(\sqrt{x}\)的导数为\(1/4\)，于是
\begin{equation*}
 \sqrt{4+h}=2+\frac h4+r(h),\qquad \frac{r(h)}h\to0.
\end{equation*}
还可以由有理化直接检验误差：
\begin{equation*}
 \sqrt{4+h}-2-\frac h4=-\frac{h^2}{4(\sqrt{4+h}+2)^2}.
\end{equation*}
近似值的意义与误差的控制，在这里来自同一个增量表达式。

一次近似也可以估计测量误差。正方形边长由\(a\)变为\(a+h\)，面积的实际增量为\(2ah+h^2\)，微分为\(2ah\)。当\(a>0\)固定、\(|h|\)很小时，面积的相对变化约为\(2h/a\)，是边长相对变化的两倍。

若\(f'(x_0)=0\)，一次近似成为\(f(x_0+h)\approx f(x_0)\)。这表示增量比\(h\)更快趋零，并不表示函数在附近恒定。\(x^2\)与\(x^3\)在原点的线性部分都为零，实际增量却分别为\(h^2\)、\(h^3\)。要继续区分这些变化，后面需要更高阶的近似。

\begin{exercise}
\mbox{}\par\nobreak

（1）用微分近似计算\(\sqrt{4.04}\)，再用上式给出误差上界。

（2）对\(f(x)=1/x\)，写出在\(x=1\)处的增量、微分及两者之差。

（3）设\(f(0)=0\)，\(x\ne0\)时\(f(x)=x^2\sin(1/x)\)。证明\(f\)在原点可导，并写出其在原点的一次近似。

（4）设\(f(a)=2\)、\(f'(a)=3\)。分别写出\(f(a+2h)\)与\(f(a-h)\)在\(h\to0\)时的一次近似，并求\(\displaystyle\lim_{h\to0}\frac{f(a+2h)-f(a-h)}h\)。

（5）圆的半径为\(r>0\)，测量误差为\(h\)。写出面积的实际增量与微分；若半径的相对误差不超过\(1\%\)，估计面积的相对误差，并给出由精确增量得到的上界。

（6）设\(f\)在\(a\)处可导，证明\(\bigl(f(a+h)+f(a-h)-2f(a)\bigr)/h\to0\)。这一结论能否推出\(f\)在\(a\)附近为常值？

（7）设\(x<0\)时\(f(x)=ax+b\)，\(x\ge0\)时\(f(x)=x^2\)。分别求\(f\)在原点连续、可导的条件，并用增量的线性部分解释所得结果。
\end{exercise}

\endgroup
\FloatBarrier
\newpage
\subsection{向量值函数的微分}\label{subsec:calculus-vector-differential}

回到第\ref{subsec:vector-valued-functions}小节的动点。若位置由 \(\vec r(t)\) 给出，从时刻 \(t\) 到 \(t+h\) 的位移就是 \(\vec r(t+h)-\vec r(t)\)。取第三章的抛物线 \(\vec r(t)=(t,t^2)\)，直接相减得
\[
 \vec r(t+h)-\vec r(t)=(h,2th+h^2)=h(1,2t)+(0,h^2).
\]
先看速度。将位移除以 \(h\)，得到平均速度 \((1,2t+h)\)；令 \(h\to0\)，就得到时刻 \(t\) 的瞬时速度 \((1,2t)\)。其中 \(2t\) 只表示纵坐标的变化率；速度是整个向量，它的模 \(\sqrt{1+4t^2}\) 是速率。

一般地，设 \(\vec r(t)=(x(t),y(t))\)。若下面的向量极限存在，就称它为 \(\vec r\) 在 \(t\) 处的{\heiti{}导数}：
\[
 \vec r\,'(t)=\lim_{h\to0}\frac{\vec r(t+h)-\vec r(t)}h.
\]
差商的两个分量正是 \(x,y\) 各自的差商。由第三章的向量极限判据，\(\vec r\) 可导当且仅当两个分量都可导，此时 \(\vec r\,'(t)=(x'(t),y'(t))\)。当导向量非零时，它给出轨迹在该处的切向。

再看位移。上例中 \(h(1,2t)\) 随时间增量 \(h\) 线性变化，余下的 \((0,h^2)\) 的模除以 \(|h|\) 趋于零。一般地，从导数的定义即得
\[
 \vec r(t+h)-\vec r(t)=h\vec r\,'(t)+\vec R(h),\qquad
 \frac{|\vec R(h)|}{|h|}\longrightarrow0.
\]
取 \(dt=h\)，称位移的线性部分 \(d\vec r=\vec r\,'(t)\,dt\) 为{\heiti{}微分}。在动点的情形，它就是按当前速度运动 \(h\) 个时间单位所产生的位移。

第三章的例\ref{ex:vector-position-segment}中，两个点沿同一条线段从 \(A\) 走向 \(B\)。记 \(\vec u=\overrightarrow{AB}\)，它们的位置分别为 \(\vec r(t)=\vec a+t\vec u\)、\(\vec s(t)=\vec a+t^2\vec u\)，其中 \(0\leq t\leq1\)。直接比较一段时间内的位移：
\[
 \vec r(t+h)-\vec r(t)=h\vec u,\qquad
 \vec s(t+h)-\vec s(t)=2th\vec u+h^2\vec u.
\]
第一个点的速度恒为 \(\vec u\)，微分恰好等于实际位移；第二个点的速度为 \(2t\vec u\)，用微分估计位移时，相差 \(h^2\vec u\)。两点先后经过中点时，前者的速度是 \(\vec u\)，后者的速度是 \(\sqrt2\,\vec u\)：同一位置、同一方向，速率却不同。

\begin{exercise}
\mbox{}\par\nobreak

（1）对 \(\vec r(t)=(t,t^2)\)，求 \(t_0=1\) 处的位置向量的一次近似，直接检验余项条件。

（2）对参数曲线 \(\vec r(t)=(t^2,t^3)\)，分别求 \(t_0=1\) 与 \(t_0=0\) 处位移的线性部分。零点处的导向量能否给出切线方向？再由割线方向判断该处的切线。
\end{exercise}

\subsection{局部线性模型*}
\label{subsec:calculus-linear-model}

前两小节中，自变量的增量都是实数 \(h\)，函数的增量却可以是实数，也可以是向量。为什么两种微分都有同样的形式？设 \(L\) 是关于 \(h\) 的线性映射，由数乘的运算律，\(L(h)=hL(1)\)。所以只要知道单位增量的像 \(L(1)\)，整个线性映射便已确定；它的值是数还是向量，都不影响这个结论。

可导性要求的正是这样一个映射 \(L\)：用 \(L(h)\) 近似实际增量，误差的大小与 \(|h|\) 之比须趋于零。将实际增量除以 \(h\)，所得差商便趋于 \(L(1)\)，因而 \(L(1)\) 就是导数，线性主部也由此唯一确定。在基点 \(t_0\)，实值函数的微分把 \(h\) 送到 \(hf'(t_0)\)，位置函数的微分则把同一个 \(h\) 送到 \(h\vec r\,'(t_0)\)。

切线对应的映射为
\begin{equation*}
 x\longmapsto f(x_0)+L(x-x_0).
\end{equation*}
它在基点和平移后的增量之间建立联系，具有仿射形式。线性主部描述变化，基点决定这段近似放在什么位置。

例如，将函数的图像写成 \(\vec r(t)=(t,f(t))\)，则向量的一次近似给出
\[
 (x,y)=(t_0,f(t_0))+h(1,f'(t_0)).
\]
由第一分量可知 \(h=x-t_0\)，代入第二分量，恰得前面的切线方程。斜率与切向量在这里描述的是同一条直线。

\begin{exercise}
\mbox{}\par\nobreak

（1）求\(f(x)=x^3\)在\(x_0=1\)处的线性主部，并区分它与切线所给的仿射映射。

（2）回看例\ref{ex:calculus-gaokao-polynomial-tangent}，写出函数在\(x=1\)处的线性主部，并求原函数与切线在\(x=1+h\)处的高度差。当\(|h|\le1/10\)时，给出这个误差关于\(h^2\)的上界。
\end{exercise}

\FloatBarrier
\newpage
\section{导数的计算}
\label{sec:calculus-differentiation}

\begingroup
\setlength{\parskip}{3.5pt}
\setlength{\abovedisplayskip}{10pt}
\setlength{\belowdisplayskip}{10pt}

按定义求导，关键是研究函数增量与自变量增量的比。若函数由若干已知函数组合而成，可以先研究这些组合如何改变增量，再把所得法则反复使用。

\subsection{四则运算法则}

\begin{theorem}\label{thm:calculus-derivative-rules}
若\(f,g\)在所考察的点可导，则
\begin{equation*}
 \begin{aligned}
 (af+bg)'&=af'+bg',\\
 (fg)'&=f'g+fg',\\
 \left(\frac fg\right)'&=\frac{f'g-fg'}{g^2}\quad(g\ne0),
 \end{aligned}
\end{equation*}
其中\(a,b\)为常数。
\end{theorem}

线性性质由极限运算法则直接得到。设两个因子的增量为\(\Delta f,\Delta g\)，乘积的增量为
\begin{equation*}
 \Delta(fg)=g\Delta f+f\Delta g+\Delta f\Delta g.
\end{equation*}
除以\(h\)，前两项分别趋于\(gf'\)、\(fg'\)，而\((\Delta f/h)\Delta g\to0\)，便得到乘积法则。这里的\(f,g\)均在固定基点取值。

非零处的倒数满足
\begin{equation*}
 \frac{1/g(x+h)-1/g(x)}h
 =-\frac{g(x+h)-g(x)}{h\,g(x+h)g(x)}
 \longrightarrow-\frac{g'(x)}{g(x)^2}.
\end{equation*}
再对\(f\cdot(1/g)\)使用乘积法则，就得到商法则。

由\((x)'=1\)及乘积法则，归纳可得
\begin{equation}\label{eq:calculus-integer-power}
 (x^n)'=nx^{n-1}\qquad(n\in\mathbb N^*).
\end{equation}
负整数幂由倒数法则得到同样形式。因此，多项式与有理函数的求导都归结到这些法则。

乘积法则可以逐次使用。例如，对三个可导函数，先把前两个看作一个整体，得到
\begin{equation*}
 (fgh)'=(fg)'h+fgh'=f'gh+fg'h+fgh'.
\end{equation*}
每一项恰有一个因子被求导，其余因子保留。若两个因子相同，则\((f^2)'=2ff'\)。例如\(\bigl((x^2+1)^2\bigr)'=4x(x^2+1)\)，也可先展开为多项式再求导。

\begin{example}
求\(f(x)=(x^2+1)/(x-1)\)的导函数。比较直接用商法则与先将分式写成\(x+1+2/(x-1)\)两种计算。
\end{example}

\newpage
\subsection{复合函数的导数}
\label{subsec:calculus-chain-rule}

设\(u=g(x)\)、\(y=f(u)\)。自变量增量\(h\)先引起\(\Delta u\)，再由\(\Delta u\)引起\(\Delta y\)。两步的一次近似应当依次复合。

\begin{theorem}[链式法则]\label{thm:calculus-chain-rule}
若\(g\)在\(x_0\)处可导，\(f\)在\(u_0=g(x_0)\)处可导，则复合函数在\(x_0\)处可导，且
\begin{equation*}
 (f\circ g)'(x_0)=f'(u_0)g'(x_0).
\end{equation*}
\end{theorem}

\begin{proof}
由\(f\)可导，可写
\begin{equation*}
 f(u_0+k)-f(u_0)=\bigl(f'(u_0)+\varepsilon(k)\bigr)k,
 \qquad\varepsilon(k)\to0,
\end{equation*}
并规定\(\varepsilon(0)=0\)。令\(k=\Delta u=g(x_0+h)-g(x_0)\)。由\(g\)连续，\(\Delta u\to0\)，故
\begin{equation*}
 \frac{\Delta y}h=\bigl(f'(u_0)+\varepsilon(\Delta u)\bigr)\frac{\Delta u}h
 \longrightarrow f'(u_0)g'(x_0).
\end{equation*}
这个表达式在\(\Delta u=0\)时也成立。
\end{proof}

由此，\(y=f(u)\)、\(u=g(x)\)满足
\begin{equation}\label{eq:calculus-differential-invariance}
 dy=f'(u)g'(x)\,dx=f'(u)\,du.
\end{equation}
无论\(u\)是独立变量还是中间变量，微分都保持\(dy=f'(u)\,du\)的形式，称为{\heiti{}一阶微分形式不变性}。当\(u\)是中间变量时，\(du=g'(x)\,dx\)，它表示的是\(u\)的微分。

例如，\(y=(1+x^2)^3\)中，令\(u=1+x^2\)，便有
\begin{equation*}
 dy=3u^2\,du=6x(1+x^2)^2\,dx.
\end{equation*}
识别中间变量，可以使复合结构和每一步变化率的取值位置同时明确。

对更多层的复合，可以逐层使用链式法则。设\(y=f(u)\)、\(u=g(v)\)、\(v=p(x)\)，则
\begin{equation*}
 \frac{dy}{dx}=f'(u)g'(v)p'(x)
 =f'(g(p(x)))g'(p(x))p'(x).
\end{equation*}
每个导数都在相应的中间值处取值。例如，对\(y=\bigl(1+(2x-1)^2\bigr)^3\)，依次令\(v=2x-1\)、\(u=1+v^2\)，便有\(y'=3u^2\cdot2v\cdot2=12(2x-1)\bigl(1+(2x-1)^2\bigr)^2\)。微分记号中的逐层替换与复合函数的计算次序一致。

\begin{example}
分别求\((1+x^2)^3\)、\(1+(x^2)^3\)及\((1+x^3)^2\)的导数。再设\(s=f(t)\)、\(t=2\tau+1\)，写出位置对\(\tau\)的变化率，解释系数\(2\)的作用。
\end{example}

\newpage
\subsection{反函数的导数}

第二章已经证明，区间上的连续严格单调函数有连续反函数。现在比较对应位置附近的变化率。

\begin{theorem}\label{thm:calculus-inverse-derivative}
设\(f\)在区间上连续且严格单调，在内点\(x_0\)处可导，且\(f'(x_0)\ne0\)。令\(y_0=f(x_0)\)，则反函数在\(y_0\)处可导，且
\begin{equation*}
 (f^{-1})'(y_0)=\frac1{f'(x_0)}.
\end{equation*}
\end{theorem}

\begin{proof}
令\(x=f^{-1}(y)\)。由反函数连续性，\(y\to y_0\)时\(x\to x_0\)，且\(y\ne y_0\)时\(x\ne x_0\)。所以
\begin{equation*}
 \frac{f^{-1}(y)-f^{-1}(y_0)}{y-y_0}
 =\frac{x-x_0}{f(x)-f(x_0)}\longrightarrow\frac1{f'(x_0)}.\qedhere
\end{equation*}
\end{proof}

原函数与反函数的导数在对应点互为倒数。得到可导性之后，也可以对恒等式\(f(f^{-1}(y))=y\)求导，整理出同一关系。

\begin{figure}[!ht]
\centering
\includegraphics[width=.45\linewidth]{calculus-inverse}
\caption{\(x>0\)时的平方函数及其反函数：图像和对应切线均关于\(y=x\)对称，切线斜率互为倒数。}
\label{fig:calculus-inverse}
\end{figure}

\begin{example}
设\(f(x)=x^3+x\)，其反函数为\(g\)。不求\(g\)的显式表达式，求\(g'(2)\)，并写出\(y=g(x)\)在\(x=2\)处的切线。
\end{example}

\subsection{三角函数的导数}

三角函数的求导公式以弧度制为基础。对\(0<h<\pi/2\)，比较单位圆中内接三角形、扇形及外切三角形的面积，得到
\begin{equation*}
 \sin h<h<\tan h,
 \qquad \cos h<\frac{\sin h}{h}<1.
\end{equation*}
由夹逼及奇偶性，
\begin{equation}\label{eq:calculus-sine-limit}
 \lim_{h\to0}\frac{\sin h}{h}=1.
\end{equation}
又由\(\cos h-1=-2\sin^2(h/2)\)，有\((\cos h-1)/h\to0\)。将这两个极限用于和角公式，便有
\begin{equation*}
 \frac{\sin(x+h)-\sin x}{h}
 =\sin x\frac{\cos h-1}h+\cos x\frac{\sin h}h\longrightarrow\cos x.
\end{equation*}
对余弦作相同展开，得到\((\cos x)'=-\sin x\)。其余由商法则和倒数法则直接推出：
\begin{equation*}
 \begin{aligned}
 (\sin x)'&=\cos x,&(\cos x)'&=-\sin x,\\
 (\tan x)'&=\frac1{\cos^2x},&(\cot x)'&=-\frac1{\sin^2x},\\
 (\sec x)'&=\sec x\tan x,&(\csc x)'&=-\csc x\cot x.
 \end{aligned}
\end{equation*}
例如，\((\tan x)'=(\cos^2x+\sin^2x)/\cos^2x\)，化简即得。

\begin{example}
求\(y=\sin(2x+\pi/3)\)、\(y=x\cos x\)及\(y=\tan^2x\)的导数。
\end{example}

\subsection{指数函数与对数函数的导数}

第一章以数列\(e_n=(1+1/n)^n\)的极限定义了\(e\)。求导时需要的是连续变化的增量，因此先把这一数列极限转为函数极限。

当\(h>0\)充分小时，选取正整数\(n\)，使\(1/(n+1)<h\le1/n\)。由正底数幂的单调性，
\begin{equation*}
 \left(1+\frac1{n+1}\right)^n
 \le(1+h)^{1/h}\le
 \left(1+\frac1n\right)^{n+1}.
\end{equation*}
左端为\(e_{n+1}/(1+1/(n+1))\)，右端为\(e_n(1+1/n)\)，两者均趋于\(e\)，故中间项也趋于\(e\)。再由对数连续性，
\begin{equation*}
 \lim_{h\to0^+}\frac{\ln(1+h)}h=1.
\end{equation*}
对于\(-1<h<0\)，令\(t=-h/(1+h)>0\)，则
\begin{equation*}
 \frac{\ln(1+h)}h=(1+t)\frac{\ln(1+t)}t\longrightarrow1.
\end{equation*}
因此这个极限两侧均成立，也得到\((1+h)^{1/h}\to e\)。

现在对\(x>0\)，令\(t=h/x\)，便有
\begin{equation*}
 \frac{\ln(x+h)-\ln x}h
 =\frac1x\frac{\ln(1+t)}t\longrightarrow\frac1x.
\end{equation*}
所以\((\ln x)'=1/x\)。自然指数是自然对数的反函数，在对应点\(y=e^x\)处使用反函数公式，得到\((e^x)'=e^x\)。一般底数由换底与复合关系整理：
\begin{equation}\label{eq:calculus-exponential-logarithmic}
 (a^x)'=a^x\ln a,\qquad
 (\log_a x)'=\frac1{x\ln a}\quad(a>0,\ a\ne1).
\end{equation}

\begin{example}
求\(y=e^{-x^2}\)、\(y=\ln(1+x^2)\)及\(y=\ln\frac{1+x}{1-x}\)的导数，分别说明选择的中间变量。
\end{example}

\begin{example}
根据第二章的指数表示，求\(\sinh x\)、\(\cosh x\)、\(\tanh x\)、\(\coth x\)的导数，并用双曲函数表示结果。
\end{example}

\begin{example}[2014年广东卷理科第10题]
求曲线\(y=e^{-5x}+2\)在点\((0,3)\)处的切线方程。
\end{example}

\begin{example}[2014年重庆卷理科第20题第（1）问]
\label{ex:calculus-chongqing-even-derivative}
设\(f(x)=ae^{2x}-be^{-2x}-cx\)，其中\(a,b,c\in\mathbb R\)。已知\(f'\)为偶函数，且曲线在\((0,f(0))\)处的切线斜率为\(4-c\)，求\(a,b\)。
\end{example}

\subsection{幂函数的导数}

对实数\(\alpha\)，在\(x>0\)上有\(x^\alpha=e^{\alpha\ln x}\)，由链式法则，
\begin{equation*}
 (x^\alpha)'=e^{\alpha\ln x}\frac\alpha x=\alpha x^{\alpha-1}.
\end{equation*}
它与前面整数幂公式一致。整数幂在负半轴、原点的可用范围，由原来的函数定义与式\eqref{eq:calculus-integer-power}确定。

\(x^x\)的指数随\(x\)改变，不能把它当作固定指数的幂函数。但在\(x>0\)时，\(x^x=e^{x\ln x}\)，求导得到
\begin{equation*}
 (x^x)'=x^x(1+\ln x).
\end{equation*}
这个计算体现了表达式外形与函数组成方式之间的区别。

\subsection{反三角函数的导数}

设\(y=\arcsin x\)，则\(x=\sin y\)、\(-\pi/2\le y\le\pi/2\)。在\(-1<x<1\)时，\(\cos y>0\)，故
\begin{equation*}
 (\arcsin x)'=\frac1{\cos y}=\frac1{\sqrt{1-x^2}}.
\end{equation*}
这里根号的正号由选定的反函数分支确定。同理由反余弦、反正切的对应关系，得到
\begin{equation*}
 (\arccos x)'=-\frac1{\sqrt{1-x^2}},\qquad
 (\arctan x)'=\frac1{1+x^2}.
\end{equation*}
前两式用于\((-1,1)\)，后一式用于\(\mathbb R\)。

\subsection{隐函数求导}

曲线方程未必直接写成\(y=f(x)\)，但若在所考察的位置附近，它已经确定了一个可导函数，就可以把\(y\)看成\(x\)的函数，对等式两侧求导。

例如，圆\(x^2+y^2=R^2\)在\(y\ne0\)的点附近，可由\(y=\pm\sqrt{R^2-x^2}\)确定相应的可导分支。沿这个分支求导，
\begin{equation*}
 2x+2yy'=0,\qquad y'=-\frac xy.
\end{equation*}
求导时，\((y^2)'=2yy'\)，已经使用了链式法则。

同样，椭圆\(x^2/a^2+y^2/b^2=1\)在\(y\ne0\)处有
\begin{equation*}
 \frac x{a^2}+\frac y{b^2}y'=0.
\end{equation*}
代入切线方程并利用切点在椭圆上，就还原第四章的\(x_0x/a^2+y_0y/b^2=1\)。这里先由显式分支确认局部可导性，再利用原方程简化计算。

\begin{example}
在双曲线\(x^2/4-y^2=1\)的点\((2\sqrt2,1)\)附近求\(dy/dx\)，并确定该点处的切线。再用\(y\)作为自变量，说明圆在\((R,0)\)处的切线方向。
\end{example}

\subsection{参数方程求导}

设\(x=\varphi(t)\)、\(y=\psi(t)\)，两个函数在\(t_0\)附近可导，且\(\varphi\)有连续的局部反函数。当\(\varphi'(t_0)\ne0\)时，链式法则与反函数公式给出
\begin{equation*}
 \left.\frac{dy}{dx}\right|_{t=t_0}=\frac{\psi'(t_0)}{\varphi'(t_0)}.
\end{equation*}
若坐标导数连续且\(\varphi'(t_0)\ne0\)，则它在附近同号，后面的单调性判据保证所需的局部反函数存在。

坐标的一次近似表明，非零向量\((\varphi'(t_0),\psi'(t_0))\)给出切线方向，也适用于斜率不存在的位置。例如，椭圆\(x=a\cos t\)、\(y=b\sin t\)的切向量为\((-a\sin t,b\cos t)\)。

\begin{example}
用参数求导求上述椭圆在\(t=\pi/4\)和\(t=0\)处的切线。再将参数换为\(t=2s\)，比较对应位置的切向量与切线。
\end{example}

\subsection{求导方法整理}

求导先识别函数结构：四则运算用相应法则，复合用链式法则，反函数在对应点使用倒数关系。隐式关系和参数表示的计算也以这些法则为依据。

同一个函数的不同表达式，可能对应不同的计算次序。例如，\(\ln\sqrt{1+x^2}\)可以逐层使用链式法则，也可以先利用对数性质化简。选择表达式时，应让各层运算更清楚；计算完成后，还可以用另一种表示核对结果。

设\(u(x)>0\)，对函数\(y=u(x)^{v(x)}\)取对数可得\(\ln y=v\ln u\)，因此
\begin{equation*}
 y'=u^v\left(v'\ln u+v\frac{u'}u\right).
\end{equation*}
这种先取对数再求导的方法，称为{\heiti{}对数求导法}。它也适合整理若干正因子的乘积。

\begin{exercise}[本节习题]
\mbox{}\par\nobreak

（1）求导：\(x^2e^{-x}\)，\(\sin(x^2)\)，\(\sin^2x\)，\(\ln(\ln x)\)，\(\arctan\sqrt x\)。

（2）设\(y=(1+x^2)^{\sin x}\)，用指数表示和对数求导两种方法求导。

（3）设\(f\)可导，求\(f(x^2+1)\)及\([f(x)]^2\)的导数。两式中的\(f'\)分别在哪里取值？

（4）设\(f(1)=2\)、\(f'(1)=3\)，且\(f\)在\(1\)附近有可导反函数\(g\)。求\(g'(2)\)及\(\left.\frac d{dx}g(x^2+1)\right|_{x=1}\)。

（5）参数曲线\(x=t^2\)、\(y=t^3\)在\(t=1\)处的切线是什么？在\(t=0\)处，坐标导数同时为零，此时能否仍由割线极限判断切线？

（6）设\(f(0)=0\)，\(x\ne0\)时\(f(x)=x^2\sin(1/x)\)。求\(f'\)，并判断\(f'\)在原点是否连续。

（7）（2013年广东卷理科第10题）若曲线\(y=kx+\ln x\)在点\((1,k)\)处的切线平行于\(x\)轴，求\(k\)。

（8）（2024年全国甲卷文科第7题改编）设\(f(x)=(e^x+2\sin x)/(1+x^2)\)。求曲线\(y=f(x)\)在点\((0,1)\)处的切线与两坐标轴所围成的三角形面积。

（9）（2024年新课标I卷第13题）若曲线\(y=e^x+x\)在点\((0,1)\)处的切线也是曲线\(y=\ln(x+1)+a\)的切线，求\(a\)。

（10）设\(y=x^2\sqrt{1+x^2}/(1+x)^3\)，\(x>0\)。用对数求导法求\(y'/y\)，再求\(y'(1)\)。比较这种方法与直接使用乘积、商法则的计算。

（11）由双曲函数的指数表达式求\((\tanh x)'\)、\((\coth x)'\)。再对恒等式\(\cosh^2x-\sinh^2x=1\)求导，核对\(\sinh x\)、\(\cosh x\)的导数公式。

（12）（2012年新课标卷文科第13题）求曲线\(y=x(3\ln x+1)\)在点\((1,1)\)处的切线方程。

（13）（2013年福建卷理科第17题第（1）问）设\(f(x)=x-2\ln x\)，求曲线在\((1,f(1))\)处的切线方程。

（14）（2013年新课标I卷理科第21题第（1）问）设\(f(x)=x^2+ax+b\)、\(g(x)=e^x(cx+d)\)。两曲线都过\(P=(0,2)\)，且在该点有相同的切线\(y=4x+2\)，求\(a,b,c,d\)。

（15）由\(x^2+xy+y^2=3\)确定经过\((1,1)\)的局部可导曲线，求该点处的切线与法线。再用隐式求导计算椭圆\(x^2/a^2+y^2/b^2=1\)在\(y\ne0\)处的斜率，与参数求导核对。

（16）设可导函数\(f\)满足\(f(2x)=2f(x)\)，对等式求导。再设\(f\)在零点可导，利用\(f(x)=2^nf(x/2^n)\)证明\(f(x)=f'(0)x\)。
\end{exercise}
\endgroup

\FloatBarrier
\newpage
\subsection{线性主部的运算*}

设\(g\)在\(x_0\)处的线性主部为\(L_g(h)=g'(x_0)h\)，\(f\)在\(u_0=g(x_0)\)处的线性主部为\(L_f(k)=f'(u_0)k\)。链式法则说明
\begin{equation*}
 L_{f\circ g}=L_f\circ L_g.
\end{equation*}
导数的乘积，是两个一维线性映射复合后的系数。用微分写\(dy=f'(u)\,du\)，正是在中间增量上先作用\(L_g\)，再作用\(L_f\)；一阶微分形式不变性也就获得线性映射复合的解释。

若\(f'(x_0)\ne0\)，则\(L_f\)可逆。反函数求导公式相应地写成
\begin{equation*}
 L_{f^{-1}}=(L_f)^{-1},
\end{equation*}
两边都在对应基点处理解。求导法则把函数的局部运算送到了线性映射的运算中，而函数本身仍可以是非线性的。

这个结构还说明，求导时引入中间变量不会改变最终的线性主部，但各主部必须在正确的基点复合。

\begin{exercise}
\mbox{}\par\nobreak

（1）设\(f(u)=u^2\)、\(g(x)=1+x^3\)。分别求\(g\)在\(1\)处、\(f\)在\(g(1)\)处的线性主部，再复合得到\(f\circ g\)在\(1\)处的线性主部。

（2）设平面可逆线性变换\(A\)作用在参数曲线\(r(t)\)上。由一次近似证明，像曲线的切向量为\(Ar'(t)\)，并与第四章仿射变换保持切线的结论联系起来。

（3）设\(f(a)=a\)、\(f'(a)=\lambda\)，记\(F_n\)为\(f\)与自身复合\(n\)次所得的函数。求\(F_n\)在\(a\)处的线性主部，并说明不动点条件在复合中的作用。

（4）设\(f(a)=a\)、\(f'(a)\ne0\)，且\(f\)在\(a\)附近有可导反函数\(g\)。用线性主部的加法与求逆证明\(|(f+g)'(a)|\ge2\)，并确定等号成立的条件。
\end{exercise}

\FloatBarrier
\newpage
\section{微分中值定理}
\label{sec:calculus-mean-value}

\begingroup
\setlength{\parskip}{5pt}

导数描述一点附近的变化。若要根据各点的导数判断一个区间上的函数变化，就需要把局部信息与两端的函数值联系起来。中值定理建立的正是这种联系。

\subsection{极值的必要条件}

\begin{definition}
若存在\(x_0\)的某个邻域，使其中所有\(x\)都满足\(f(x)\le f(x_0)\)，则称\(f\)在\(x_0\)处取得{\heiti{}极大值}；将不等号反向，得到{\heiti{}极小值}，统称{\heiti{}极值}。满足\(f'(x_0)=0\)的点称为{\heiti{}驻点}。
\end{definition}

极值只在一点附近比较。第二章所说的最大值、最小值，则在给定的整个集合上比较。以下极值均在所考察区间的内点讨论，端点作为区间最值的候选位置另行处理。

\begin{theorem}[Fermat定理]\label{thm:calculus-fermat}
若\(f\)在内点\(x_0\)处取得极值，且在该点可导，则\(f'(x_0)=0\)。
\end{theorem}

\begin{proof}
设\(x_0\)为极大值点。充分小的\(h>0\)使差商不大于零，充分小的\(h<0\)使差商不小于零。两侧极限相等，故只能为零。极小值同理。
\end{proof}

因此，可导函数的内点极值一定出现在驻点。但\(x^3\)在原点导数为零，附近却有比\(0\)大和比\(0\)小的函数值，故驻点未必是极值点。

\begin{example}
比较\(x^2\)、\(x^3\)与\(|x|\)在原点的极值和可导性。再考察\(f(x)=x\)在\([0,1]\)上的最值，说明驻点条件为什么不能替代端点比较。
\end{example}

\subsection{Rolle定理}

\begin{theorem}[Rolle定理]\label{thm:calculus-rolle}
设\(f\)在\([a,b]\)上连续，在\((a,b)\)内可导，且\(f(a)=f(b)\)。则存在\(\xi\in(a,b)\)，使\(f'(\xi)=0\)。
\end{theorem}

\begin{proof}
由闭区间连续函数的最值定理，\(f\)取得最大值和最小值。若\(f\)为常值，结论立即成立；否则，因两个端点的值相等，至少有一个最值在内部取得。对此内点使用Fermat定理，便得结论。
\end{proof}

两个端点同高，区间内的图像若发生变化，就必在内部达到比端点更高或更低的位置。最值定理保证这个高度可以取到，可导性则使相应切线水平。

\begin{example}
对\(f(x)=x^3-x\)，在\([-1,1]\)上求Rolle定理所保证的点。再考察\(|x|\)在\([-1,1]\)上的情形，说明可导条件的作用。
\end{example}

Rolle定理也能比较函数零点与导函数零点：任意两个不同零点之间，至少有一个导函数零点。因此，若\(f'\)在某区间内至多有\(m\)个零点，\(f\)在该区间内至多有\(m+1\)个不同零点。

这个结论可以逐次使用。若一个多项式有\(n+1\)个不同零点，在相邻零点之间应用Rolle定理，其导函数就至少有\(n\)个不同零点。每求导一次，保证得到的零点数只减少一个。次数为\(n\)的非零多项式连续求导\(n\)次后得到非零常数，所以不可能有\(n+1\)个不同实根。这也从函数变化的角度解释了多项式次数对实根数的限制。

\subsection{Lagrange中值定理}

当两端的函数值不相同时，可以先减去两端连线，使问题回到Rolle定理。

\begin{theorem}[Lagrange中值定理]\label{thm:calculus-lagrange}
若\(f\)在\([a,b]\)上连续、内部可导，则有\(\xi\in(a,b)\)，使
\begin{equation}\label{eq:calculus-lagrange}
 f(b)-f(a)=f'(\xi)(b-a).
\end{equation}
\end{theorem}

\begin{proof}
令
\begin{equation*}
 F(x)=f(x)-f(a)-\frac{f(b)-f(a)}{b-a}(x-a).
\end{equation*}
则\(F(a)=F(b)=0\)，且满足Rolle定理的条件。由\(F'(\xi)=0\)，便得所需等式。
\end{proof}

式\eqref{eq:calculus-lagrange}表明，整段区间的平均变化率，在某个内部位置由瞬时变化率实现。对于位置函数，它意味着某一时刻的速度等于整段过程的平均速度；在图像上，则有一条切线平行于端点连线。

\begin{figure}[!ht]
\centering
\includegraphics[width=.83\linewidth]{calculus-mean-value}
\caption{左图为Rolle定理的水平切线；右图为Lagrange中值定理中与端点连线平行的切线。}
\label{fig:calculus-mean-value}
\end{figure}

这里的\(\xi\)位于两点之间，通常随两点改变。中值定理给出它的存在性，并不要求先求出它的位置，也不要求导函数连续。

\begin{example}
不求\(\sqrt{101}\)的精确值，利用\(\sqrt x\)在\([100,101]\)上的中值定理，给出\(\sqrt{101}-10\)的上下界。与在\(100\)处作一次近似比较。
\end{example}

\subsection{函数增量的估计}

若区间内\(m\le f'(x)\le M\)，对任意\(u<v\)使用中值定理，便有
\begin{equation}\label{eq:calculus-increment-bounds}
 m(v-u)\le f(v)-f(u)\le M(v-u).
\end{equation}
特别地，若\(|f'(x)|\le K\)，则
\begin{equation}\label{eq:calculus-derivative-bound}
 |f(v)-f(u)|\le K|v-u|.
\end{equation}
\begin{samepage}
导数的统一上界，转化为了两点函数值之差的控制。例如，\(|\cos x|\le1\)给出\(|\sin v-\sin u|\le|v-u|\)。这也为第二章的一致连续提供了一类可检验的充分条件。
\par
\end{samepage}

\begin{corollary}\label{cor:calculus-zero-derivative}
区间上导数恒为零的函数是常值函数。因此，若\(f'=g'\)，则\(f-g\)为常数。
\end{corollary}

\begin{proof}
对区间内任意两点使用中值定理，函数值之差均为零。将此结论用于\(f-g\)，即得第二项。
\end{proof}

导数确定各点的变化率，却不能确定函数整体的高度；加上一个函数值，就能消除这个常数的不确定性。

两点间的增量估计还可以反过来限制区间长度。若\(|f'|\le K\)、\(K>0\)，而两个位置的函数值相差\(D\)，则它们的距离至少为\(D/K\)。例如，一辆车的速度绝对值始终不超过\(V\)，要使位置改变\(S\)，所用时间不可能小于\(|S|/V\)。中值定理在这里同时说明了速度上界与完成位移所需时间的关系。

\begin{example}
对\(x>0\)，证明
\begin{equation*}
 \frac{x}{1+x}<\ln(1+x)<x.
\end{equation*}
再设\(f'(x)=2x\)、\(f(1)=3\)，确定\(f\)。
\end{example}

\begin{example}[2024年北京卷第20题第（2）问改编]
设\(f(x)=x+k\ln(1+x)\)，\(k\ne0\)。证明对任意\(t>0\)，曲线在\((t,f(t))\)处的切线都不经过原点。尝试用\(\ln(1+x)\)的中值定理解释这一结论。
\end{example}

\begin{example}
设\(f'=f\)，证明\(e^{-x}f(x)\)为常数；若再给定\(f(0)=1\)，确定\(f\)。
\end{example}

\newpage
\setlength{\parskip}{3.5pt}
\setlength{\abovedisplayskip}{10pt}
\setlength{\belowdisplayskip}{10pt}
\begin{exercise}[本节习题]
\mbox{}\par\nobreak

（1）证明\(|\arctan x-\arctan y|\le|x-y|\)。

（2）设\(f\)在\([0,1]\)连续、内部可导，且\(f(0)=0\)、\(1\le f'(x)\le2\)。估计\(f(1)\)及\(f(3/4)-f(1/4)\)。

（3）证明\(e^x=1+x\)只有一个实根。可先用Rolle定理控制根数，再判断函数值。

（4）若\(f\)在\(\mathbb R\)上可导且有周期\(T>0\)，证明任意长度为\(T\)的开区间内至少有一个\(f'\)的零点。

（5）设\(f\)在\(\mathbb R\)上可导。证明：\(f'\)为偶函数，当且仅当\(f(x)-f(0)\)为奇函数；\(f'\)为奇函数，当且仅当\(f\)为偶函数。结合例\ref{ex:calculus-chongqing-even-derivative}，说明函数值\(f(0)\)的作用。

（6）对\(0<a<b\)，证明\(\dfrac{b-a}{b}<\ln\dfrac ba<\dfrac{b-a}{a}\)，并据此估计\(\ln\dfrac{101}{100}\)。

（7）设\(f\)在\([a,b]\)上连续、内部可导，且\(|f'|\le K\)。若\(f(a)=f(b)=0\)，证明
\begin{equation*}
 |f(x)|\le K\min\{x-a,b-x\}\le\frac{K(b-a)}2\qquad(x\in[a,b]).
\end{equation*}

（8）设区间\(I\)上\(f'\ge m>0\)，且\(f(\alpha)=0\)。证明\(|x-\alpha|\le|f(x)|/m\)（\(x\in I\)），并解释它对估计近似根误差的作用。

（9）（2013年北京卷理科第18题第（2）问）曲线\(y=\ln x/x\)在\((1,0)\)处的切线为\(y=x-1\)。证明除切点外，曲线都位于这条直线的下方。可先用中值定理比较\(\ln x\)与\(x-1\)。

（10）设\(f\)在\([a,b]\)上连续、内部可导，且有三个不同零点。证明\(f'\)至少有两个不同零点。若\(f'\)也可导，证明它的导函数至少有一个零点。

（11）设\(f\)在\([0,1]\)上连续、内部可导，且\(f(0)=f(1)=0\)、\(f(1/2)=1\)。证明存在\(u\in(0,1/2)\)、\(v\in(1/2,1)\)，使\(f'(u)-f'(v)=4\)。

（12）设\(f\)在\([0,1]\)上连续、内部可导，且\(f(0)=0\)、\(f(1)=1\)。证明存在\(\xi\in(0,1)\)，使\(f'(\xi)=2\xi\)。说明应先从\(f\)中减去哪个函数，再使用Rolle定理。

（13）设\(f\)将\([0,1]\)映入自身，在该区间连续、内部可导，且\(|f'|\le q<1\)。证明方程\(f(x)=x\)有且只有一个解，分别指出连续性与导数上界的作用。

（14）设\(f\)在\([0,1]\)上连续、内部可导，两端的函数值为零。证明存在\(\xi\in(0,1)\)，使\(f'(\xi)=f(\xi)\)。

（15）若区间上\(f'\ne0\)，证明\(f\)至多有一个零点。
\end{exercise}

\endgroup
\FloatBarrier
\newpage
\begingroup
\setlength{\parskip}{2.5pt}
\setlength{\abovedisplayskip}{7pt}
\setlength{\belowdisplayskip}{7pt}
\subsection{导函数的性质*}

函数可导，并不保证导函数连续。另一方面，导函数来自局部差商，它的取值也并非可以任意变化。中值定理先给出导函数极限与导数之间的联系。

\begin{theorem}[导数极限定理]\label{thm:calculus-derivative-limit}
设\(f\)在\(a\)处连续，在\(a\)的某个去心邻域内可导。若\(\lim_{x\to a}f'(x)=L\in\mathbb R\)，则\(f\)在\(a\)处可导，且\(f'(a)=L\)。
\end{theorem}

\begin{proof}
对充分小的非零\(h\)，在\(a\)与\(a+h\)之间使用Lagrange中值定理，得到
\begin{equation*}
 \frac{f(a+h)-f(a)}h=f'(\xi_h),\qquad |\xi_h-a|<|h|.
\end{equation*}
当\(h\to0\)时，\(\xi_h\to a\)，故差商趋于\(L\)，结论成立。
\end{proof}

只从一侧趋近时，同一证明得到相应的单侧导数。因此，对连接处连续的分段函数，若两侧导函数分别趋于同一个有限值，就可以据此确定连接处的导数。

反过来，可导并不要求附近的导函数具有极限。前面研究的\(f(0)=0\)、\(f(x)=x^2\sin(1/x)\)（\(x\ne0\)）满足\(f'(0)=0\)，而\(x\ne0\)时\(f'(x)=2x\sin(1/x)-\cos(1/x)\)，在零点附近不断振荡。判断一点处是否可导，最终依据仍是差商。

\begin{theorem}[Darboux定理]\label{thm:calculus-derivative-intermediate-value}
设\(f\)在区间\(I\)上可导，\(a<b\)为其中两点。若\(\lambda\)严格介于\(f'(a)\)与\(f'(b)\)之间，则存在\(\xi\in(a,b)\)，使\(f'(\xi)=\lambda\)。
\end{theorem}

\begin{proof}
先设\(f'(a)<\lambda<f'(b)\)，令\(F(x)=f(x)-\lambda x\)，则\(F'(a)<0\)、\(F'(b)>0\)。由导数定义，在\(a\)的右侧有\(F(x)<F(a)\)，在\(b\)的左侧有\(F(x)<F(b)\)。所以\(F\)在\([a,b]\)上的最小值不能在端点取得。

由最值定理，最小值在某个\(\xi\in(a,b)\)处取得。Fermat定理给出\(F'(\xi)=0\)，即\(f'(\xi)=\lambda\)。若两个端点导数的大小顺序相反，则改取\(F\)的最大值，同样得到结论。
\end{proof}

连续函数的介值性在第二章已经建立，这里则没有假定\(f'\)连续。导函数虽可振荡而不连续，仍不能像阶跃函数那样跳过一段值。特别地，若一个区间上的导数既取正值又取负值，则两者之间必有导数为零的位置。

这个证明也给出一种构造方式：要让导数等于\(\lambda\)，就从函数中减去斜率为\(\lambda\)的直线，再寻找内部极值。导数的取值问题由此转成了原函数的极值问题。

\begin{example}
一个函数在零点左侧恒为\(-1\)，右侧恒为\(1\)。无论怎样规定它在零点的值，它能否成为某个可导函数在整个区间上的导函数？与\(x^2\sin(1/x)\)的导函数比较。
\end{example}

\begin{exercise}
\mbox{}\par\nobreak

（1）设\(f(x)=e^x-1\)（\(x<0\)），\(f(x)=ax\)（\(x\ge0\)）。分别用差商与导数极限定理，求\(f\)在零点可导的条件。

（2）设区间上可导函数\(f\)的导数只取\(-1\)、\(1\)中的值。证明\(f\)必为\(x+C\)或\(-x+C\)，其中\(C\)为常数。

（3）若\(f\)在开区间内可导，且\(f'\)单调，证明\(f'\)连续。可以先在任意内点的两侧建立单调函数的单侧极限，再使用导数极限定理。
\end{exercise}

\subsection{Cauchy中值定理*}

Lagrange中值定理比较函数增量与自变量增量。若将自变量函数\(x\)换成另一个函数，就得到两种变化之间的联系。

\begin{theorem}[Cauchy中值定理]\label{thm:calculus-cauchy-mean-value}
若\(f,g\)在\([a,b]\)上连续、内部可导，则有\(\xi\in(a,b)\)，使
\begin{equation}\label{eq:calculus-cauchy-mean-value}
 (f(b)-f(a))g'(\xi)=(g(b)-g(a))f'(\xi).
\end{equation}
\end{theorem}

\begin{proof}
为消去端点增量，取
\begin{equation*}
 H(x)=(f(b)-f(a))(g(x)-g(a))-(g(b)-g(a))(f(x)-f(a)).
\end{equation*}
则\(H(a)=H(b)=0\)。由Rolle定理有\(H'(\xi)=0\)，展开就是式\eqref{eq:calculus-cauchy-mean-value}。
\end{proof}

若\(g'\)在内部处处不为零，Rolle定理还保证\(g(b)\ne g(a)\)，因而可写为
\begin{equation*}
 \frac{f(b)-f(a)}{g(b)-g(a)}=\frac{f'(\xi)}{g'(\xi)}.
\end{equation*}
右边的两个导数在同一个点取值。分别对\(f\)、\(g\)使用Lagrange中值定理，一般不能保证得到同一个点，故不能代替这一结论。

\begin{exercise}
对\(0<a<b\)，取\(f(x)=\ln x\)、\(g(x)=x^2\)，证明存在\(\xi\in(a,b)\)，使\(\ln(b/a)/(b^2-a^2)=1/(2\xi^2)\)，并据此给出该比值的上下界。
\end{exercise}

\subsection{洛必达法则*}

在\(f(x)\)与\(g(x)\)都趋零时，比值的极限不能由两者极限相除得到。若把两者看成从共同零点出发的增量，Cauchy中值定理就能把原比值与导数比联系起来。

\begin{theorem}[洛必达法则的\(0/0\)型]\label{thm:calculus-lhopital}
设\(x\to c^+\)时\(f(x)\to0\)、\(g(x)\to0\)，且两函数在\(c\)右侧某邻域内可导，\(g'(x)\ne0\)。若
\begin{equation*}
 \lim_{x\to c^+}\frac{f'(x)}{g'(x)}=L\in\mathbb R,
\end{equation*}
则\(\lim_{x\to c^+}f(x)/g(x)=L\)。左侧版本同理；两侧适用并得到同一极限时，得到双侧版本。
\end{theorem}

\begin{proof}
补定义\(f(c)=g(c)=0\)，使端点连续。对\([c,x]\)用Cauchy中值定理，得到
\begin{equation*}
 \frac{f(x)}{g(x)}=\frac{f'(\xi)}{g'(\xi)},\qquad c<\xi<x.
\end{equation*}
当\(x\to c^+\)时，\(\xi\to c^+\)，右端趋于\(L\)。
\end{proof}
\paragraph{无穷远与无穷量之比} \mbox{}

当\(x\to+\infty\)而\(f,g\to0\)时，可令\(t=1/x\)，转为有限点处的版本；两个复合导数中的\(-1/t^2\)相消，导数比仍为\(f'/g'\)。

对于\(|f(x)|,|g(x)|\to\infty\)的情形，在相同的可导和\(g'\ne0\)条件下，若\(f'/g'\to L\in\mathbb R\)，仍有\(f/g\to L\)。证明比较固定点与变化点的增量，不将无穷大当作端点函数值。

具体地，给定\(\varepsilon>0\)，在足够靠近所考察端部的范围内，有\(|f'/g'-L|<\varepsilon\)。在这一范围内固定\(x_1\)，对\(x_1\)与\(x\)之间的区间使用Cauchy中值定理，就有
\begin{equation*}
 \left|\frac{f(x)-f(x_1)}{g(x)-g(x_1)}-L\right|<\varepsilon.
\end{equation*}
再写成
\begin{equation*}
 \frac{f(x)}{g(x)}=
 \frac{f(x)-f(x_1)}{g(x)-g(x_1)}
 \left(1-\frac{g(x_1)}{g(x)}\right)+\frac{f(x_1)}{g(x)}.
\end{equation*}
由于\(|g(x)|\to\infty\)，后两个固定点产生的误差趋零，原比值最终与\(L\)任意接近。这个论证对有限端点和无穷远均适用。若导数比趋于\(+\infty\)或\(-\infty\)，将上述误差控制换成任意大的单侧界，同样得到相应结论。

\begin{example}
求\(\lim_{x\to0}(e^x-1-x)/x^2\)及\(\lim_{x\to+\infty}\ln x/x\)。再比较\(\lim_{x\to0}(\sqrt{1+x}-1)/x\)使用有理化与洛必达法则的计算。
\end{example}

使用法则之前，要先确认原式属于所讨论的类型。\(\infty-\infty\)、\(0\cdot\infty\)等表达式，可通过通分或改写成商后判断；每次继续求导，也都须重新满足条件。求导后的比值没有极限时，法则不给出判断，原比值仍可能有极限。

\begin{exercise}
\mbox{}\par\nobreak

（1）求\(\lim_{x\to0}(x-\sin x)/x^3\)及\(\lim_{x\to0}(\ln(1+x)-x)/x^2\)。

（2）求\(\lim_{x\to+\infty}x e^{-x}\)与\(\lim_{x\to0^+}x\ln x\)。

（3）计算\(\lim_{x\to0}(1+x)^{1/x}\)。说明若用对数求导和洛必达法则，这个计算为什么只是核对已有结论，不能反过来承担本章指数、对数求导公式的最初建立。

（4）令\(f(x)=x+\sin x\)、\(g(x)=x\)。比较\(x\to+\infty\)时原比值与导数比的极限。
\end{exercise}
\endgroup

\FloatBarrier
\newpage
\section{导数的应用}
\label{sec:calculus-applications}

中值定理把导数与函数增量联系起来，因而可以由导数研究区间上的单调性，再确定极值、最值和方程的根。含参数的问题也遵循这个次序，只是需要继续判断这些结构在什么参数处发生改变。

\subsection{单调性}

\begin{theorem}\label{thm:calculus-monotonicity}
设\(f\)在区间\(I\)上可导。若\(f'\ge0\)，则\(f\)单调不减；若\(f'>0\)，则\(f\)严格递增。将不等号反向，得到相应的单调不增和严格递减结论。
\end{theorem}

\begin{proof}
任取\(u<v\)，由中值定理，\(f(v)-f(u)=f'(\xi)(v-u)\)。导数的符号即给出函数值之差的符号。
\end{proof}

连续端点可以一并纳入单调区间。反过来，若可导函数单调不减，则由差商定义得\(f'\ge0\)。不过，严格递增并不要求处处导数为正，例如\(x^3\)在原点导数为零。

实际判断中，若\(f'\ge0\)，且只有有限个零点，将区间在这些点处分开，再用连续性衔接，就能得到严格递增。这样既保留了函数经过水平切线的可能，也不影响不同位置之间的严格大小关系。

\begin{figure}[!ht]
\centering
\includegraphics[width=.83\linewidth]{calculus-monotonicity}
\caption{由右图导函数的符号，确定左图函数的增减；驻点两侧是否变号，决定它是否为极值点。}
\label{fig:calculus-monotonicity}
\end{figure}

\begin{example}
求\(f(x)=x^3-3x\)的单调区间。再判断\(g(x)=x^3+x\)的单调性，并与第二章不使用导数的判断比较。
\end{example}

\begin{exercise}
\mbox{}\par\nobreak

（1）求\(x e^{-x}\)与\(\ln x/x\)的单调区间。

（2）证明\(f(x)=x-\sin x\)在\(\mathbb R\)上严格递增。导数在无穷多个点为零是否妨碍结论？可在任意两点之间选取导数为正的一小段进行比较。

（3）求实数\(a\)的范围，使\(x^3-3ax\)在\(\mathbb R\)上单调不减。

（4）（2024年新课标I卷第6题改编）设
\begin{equation*}
 f(x)=\begin{cases}-x^2-2ax-a,&x<0,\\ e^x+\ln(x+1),&x\ge0.\end{cases}
\end{equation*}
若\(f\)在\(\mathbb R\)上严格递增，求\(a\)的取值范围。分段内的导数符号与连接处的取值关系，应分别考察。
\end{exercise}

\subsection{极值与最值}

\begingroup
\setlength{\parskip}{3pt}
\setlength{\abovedisplayskip}{10pt}
\setlength{\belowdisplayskip}{10pt}

\begin{proposition}\label{prop:calculus-first-derivative-test}
设\(f\)在\(x_0\)附近连续，且在\(x_0\)两侧可导。若导数在左侧为正、右侧为负，则\(x_0\)为极大值点；若左负右正，则为极小值点。若两侧同为正或同为负，则该点无极值。
\end{proposition}

这是单调性判据的直接结果：左增右减时，两侧函数值都小于中间值；其他情形类似。它给出极值的判别，而Fermat定理只给出了可导极值点必须满足的条件。

闭区间连续函数的最大值与最小值一定存在。若内部处处可导，内部最值点必为驻点，所以只须比较端点和全部驻点的函数值。驻点是候选位置，比较后才能确定最值。

以\(f(x)=x^3-3x\)在\([-3,3]\)上为例，导数在\((-3,-1)\)、\((1,3)\)上为正，在\((-1,1)\)上为负。因此\(-1\)为极大值点，\(1\)为极小值点；但
\begin{equation*}
 f(-3)=-18,\quad f(-1)=2,\quad f(1)=-2,\quad f(3)=18.
\end{equation*}
最大值为\(18\)，最小值为\(-18\)，都在端点取得。局部极值与整个区间的最值，由同一个变化过程得到，却比较不同的范围。

在开区间或无界区间上，还应研究端部行为。例如\(x+1/x\)在\(x=1\)处取得最小值\(2\)，在\((0,1)\)上递减、\((1,+\infty)\)上递增；两端均趋于正无穷，故没有最大值。

\begin{example}
求\(f(x)=x^3-3x\)在\([0,2]\)与\([0,3]\)上的最值。再求\(g(x)=\ln x-x\)在\((0,+\infty)\)上的极值和最值，比较两种范围下的研究步骤。
\end{example}

\begin{example}[2024年新课标II卷第16题]
已知\(f(x)=e^x-ax-a^3\)。
（1）当\(a=1\)时，求曲线在\((1,f(1))\)处的切线方程；
（2）若\(f\)有极小值，且极小值小于\(0\)，求\(a\)的取值范围。
\end{example}

\begin{exercise}
\mbox{}\par\nobreak

（1）求\(f(x)=x^4-2x^2\)的极值及其在\([-2,1]\)上的最值。

（2）求\(x e^{-x}\)在\([0,+\infty)\)上的最值，并判断哪些值可以取到。

（3）设\(f(x)=x^{1/3}\)。它在原点没有有限导数，是否在原点取得极值？结合\(|x|\)比较。

\begin{samepage}
（4）（2024年新课标I卷第10题）设\(f(x)=(x-1)^2(x-4)\)，判断下列结论。

（A）\(x=3\)是\(f\)的极小值点；（B）当\(0<x<1\)时，\(f(x)<f(x^2)\)；（C）当\(1<x<2\)时，\(-4<f(2x-1)<0\)；（D）当\(-1<x<0\)时，\(f(2-x)>f(x)\)。
\par
\end{samepage}
\end{exercise}
\endgroup

\begingroup
\setlength{\parskip}{3.5pt}
\setlength{\abovedisplayskip}{10pt}
\setlength{\belowdisplayskip}{10pt}
\subsection{方程的根}

研究\(f(x)=k\)，就是比较固定图像与不同高度的水平线。连续性给出中间高度的存在性，单调性控制每一段内的根数，极值则确定水平线穿过不同单调段的临界高度。

一般地，可以先将定义域分成单调区间，分别确定各段的值域，再统计\(k\)被多少段取得。公共端点只计算一次。若只需根数上界，也可以使用Rolle定理；若要确定实际根数，还应补上存在性判断。

若连续函数先严格递减、后严格递增，两端均趋于正无穷，最小值为\(m\)，则水平线\(y=k\)在\(k<m\)时不与图像相交，\(k=m\)时交于一点，\(k>m\)时左右各交于一点。单调性控制各段的交点个数，值域保证交点存在。

两端异号保证有根，两端同号却不能排除有根，如\(x^2-1\)在\([-2,2]\)上。把区间在极值点处分开后，才便于分别使用零点存在定理和单调性。

存在性证明中的取点，可以由希望得到的不等式倒推。例如，对\(a>0\)，要在\(f(x)=x+a\ln x\)中找到一个负值，可以先令\(x=e^{-t}\)、\(t>0\)，于是\(f(e^{-t})=e^{-t}-at\)。第一项小于\(1\)，所以只须使\(at\ge1\)。取\(t=1/a\)，就得到\(f(e^{-1/a})<0\)，再与\(f(1)=1\)配对即可确定一个有根区间。

这里先找容易验证的充分条件，再选择取样点。更一般地，要使\(u(x)-v(x)>0\)，可以寻找便于处理的中间数\(c\)，使\(u(x)>c>v(x)\)。这个\(c\)和取样点都由所需的符号关系决定；无需先猜到一个特别巧的数。

\begin{example}
设\(a>0\)，证明\(e^x-ax^2\)在负半轴上有且只有一个零点，并构造一个含该零点的有限区间。可以先令\(x=-t\)，利用\(e^{-t}<1\)选择使函数值为负的\(t\)。
\end{example}

\begin{example}
讨论实数\(t\)取不同值时，方程\(x^3-3x=t\)的不同实根个数，并指出根数改变对应的函数图像位置。
\end{example}

两个图像的交点问题可以写为\(f-g=0\)。有时直接研究差更自然，有时除以一个处处为正的函数，能把参数从变化的表达式中分离出来。

\begin{example}
讨论方程\(\ln x=kx\)的正实根个数。比较研究\(\ln x-kx\)与研究\(\ln x/x\)的两种组织方式。
\end{example}

\begin{exercise}
\mbox{}\par\nobreak

（1）讨论方程\(e^x=ax\)的实根个数。

（2）证明\(x^3+x=\cos x\)有且只有一个实根，并确定一个长度不超过\(1\)的有根区间。

（3）设\(f\)在开区间\(I\)上可导且\(f'\)严格递增。证明\(f\)在\(I\)内至多有两个不同零点。

（4）（2014年新课标II卷文科第21题）设\(f(x)=x^3-3x^2+ax+2\)，曲线在点\((0,2)\)处的切线与\(x\)轴交点的横坐标为\(-2\)。求\(a\)，并证明当\(k<1\)时，曲线与直线\(y=kx-2\)只有一个交点。
\end{exercise}

\subsection{不等式的参数范围}

第二章已经把恒成立与存在性条件转为最值或值域的判断。导数现在为确定这些范围提供了工具。

要求\(f(x)\ge a\)对所有\(x\in I\)成立，就是要求\(a\)不超过全部函数值。最小值\(m\)存在时，这等价于\(a\le m\)；最小值不存在时，按实际值域判断。存在性问题只要求值域中至少有一个数满足所给比较。

例如，\(x>0\)时\(\ln x/x\)的最大值为\(1/e\)，且在\(x=e\)取得。因此
\begin{equation*}
 \ln x\le ax\ \text{对所有 }x>0\text{成立}
 \quad\Longleftrightarrow\quad a\ge\frac1e.
\end{equation*}
若改为存在\(x>0\)使\(\ln x>ax\)，则条件变为\(a<1/e\)。两个问题使用同一值域，但量词决定了比较方向。

\begin{example}
求实数\(a\)的范围，使\(x^2-2ax+2\ge0\)对所有\(x\in[0,2]\)成立。比较直接研究含参数函数的最值与把参数分离后研究一个固定函数的两种方法。
\end{example}

分离参数时，乘除因子的符号和定义域决定变形是否等价。若某个位置使参数的系数为零，可在该位置先判断原条件，再在其余范围讨论。

\paragraph{先由局部条件确定必要范围} \mbox{}

若要求\(F_a(x)\ge0\)在某区间上恒成立，而\(F_a(c)=0\)，则\(c\)附近的增量也必须满足符号要求。\(c\)为内点且函数可导时，Fermat定理给出\(F_a'(c)=0\)；若只研究\(x\ge c\)，右导数则必须非负。例如，\(e^x\ge1+ax\)对所有实数成立，先在零点使用极值必要条件就得到\(a=1\)，再由\(e^x\)的单调变化或凸性验证充分性。

局部条件只是整段条件的一部分。若它尚不能确定参数，还要回到整个研究区间；后面的Taylor展开可以在低阶项抵消时进一步辨认局部增量的符号。

\paragraph{分别控制两个部分} \mbox{}

比较\(U(x)\)与\(V(x)\)，可以分别给出前者的下界、后者的上界，再比较这两个界。特别地，若存在常数\(c\)，使整个区间上\(U(x)\ge c\)、\(V(x)\le c\)，就有\(U\ge V\)。常数分离是更强的要求，一般的逐点大小关系并不保证\(\min U\ge\max V\)。

例如，在\(x>0\)时，由\(e^x\ge1+x\)、\(\ln x\le x-1\)，相减就得到\(e^x-\ln x\ge2\)。选择中间表达式\(x+1\)、\(x-1\)，使共同的\(x\)消去，比直接寻找\(e^x-\ln x\)的最小值更方便。所选估计的方向与消去后留下的量，需要一同考虑。

当条件中出现两个自变量时，还应先确定它们怎样取值。设\(F,G\)分别在闭区间\(I,J\)上连续，记各自的最大值为\(M_F,M_G\)，\(G\)的最小值为\(m_G\)。则
\begin{equation*}
 \begin{aligned}
 &\text{对每个 }x\in I\text{，存在 }y\in J\text{ 使 }F(x)\le G(y)
 &&\Longleftrightarrow M_F\le M_G,\\
 &\text{对每个 }x\in I\text{ 和每个 }y\in J\text{ 都有 }F(x)\le G(y)
 &&\Longleftrightarrow M_F\le m_G.
 \end{aligned}
\end{equation*}
前者允许选择\(G\)的最大值与全部\(F(x)\)比较，后者则连\(G\)的最小值也必须满足要求。导数负责求出这些最值，量词决定应当比较哪两个最值。

\begin{exercise}
\mbox{}\par\nobreak

（1）求\(a\)的范围，使\(\ln x\le a(x-1)\)对所有\(x>0\)成立。

（2）分别求\(a\)的范围，使\(e^x-ax\ge0\)对所有实数\(x\)成立，以及使存在实数\(x\)满足\(e^x-ax<0\)。

（3）求\(t\)的范围，使对每个\(x\in[1,e]\)，都存在\(y\in[0,1]\)满足\(\ln x-t\le y-y^2\)。
\end{exercise}

\subsection{含参数函数的讨论}

含参数函数的分类，应依据单调结构、候选极值点或根数发生改变的位置。常见的临界情形包括导数零点合并、驻点进入研究区间、极值与指定高度相等。

例如，对\(f_a(x)=x^3-3ax\)，导数为\(3(x^2-a)\)。当\(a\le0\)时，函数严格递增；当\(a>0\)时，出现两个驻点\(\pm\sqrt a\)，对应的极大值、极小值分别为\(2a^{3/2}\)、\(-2a^{3/2}\)。因此讨论方程\(f_a(x)=t\)，可以先按\(a\)判断图像结构，再比较\(t\)与这两个临界高度。

若只研究\([0,1]\)，还需要比较驻点与端点的位置。参数\(a=1\)的意义来自\(\sqrt a=1\)，而不是因为某个代数表达式恰好可以按这个数分类。

% 题号按本地试卷PDF核对；转录汇编的条目编号与试卷有一题偏移。
\begin{example}[2023年新课标II卷第11题]
若函数\(f(x)=a\ln x+b/x+c/x^2\)（\(a\ne0\)）既有极大值也有极小值，则下列哪些结论正确？
\begin{equation*}
 \text{A. }bc>0\qquad \text{B. }ab>0\qquad
 \text{C. }b^2+8ac>0\qquad \text{D. }ac<0.
\end{equation*}
\end{example}

这类问题的导数往往可以写成一个正因子乘较简单的表达式。先确认符号由哪一部分决定，再研究其零点位置和变号情况，可以把参数条件与函数结构对应起来。

导数方程的实根还要放回原定义域。例如在\(x>0\)上，负根不会产生驻点；两个正根合并时，导数也未必变号。因此判别式、根的位置和两侧符号分别回答不同的问题，不能仅凭导数方程有两个实根就断定原函数有两个极值。

\begin{exercise}\label{exs:calculus-parameter-gaokao}
\mbox{}\par\nobreak

（1）讨论\(f_a(x)=x^3-3ax\)在\([0,1]\)上的最值，并指出需要分类的参数位置。

（2）设\(f_a(x)=\ln x-ax\)。讨论其单调性，并判断不同\(a\)下\(f_a(x)=0\)的根数。

（3）（2024年新课标I卷第18题）已知\(f(x)=\ln\frac{x}{2-x}+ax+b(x-1)^3\)。

\exercisesubitem{1}若\(b=0\)且\(f'(x)\ge0\)，求\(a\)的最小值；

\exercisesubitem{2}证明曲线\(y=f(x)\)是中心对称图形；

\exercisesubitem{3}若\(f(x)>-2\)当且仅当\(1<x<2\)，求\(b\)的取值范围。
\end{exercise}

\subsection{切线问题}

给定切点时，直接使用式\eqref{eq:calculus-tangent}。若给定的是切线经过的点，就应把切点作为未知量。

设切点横坐标为\(t\)，切线经过\((u,v)\)，则
\begin{equation}\label{eq:calculus-tangent-through-point}
 v=f(t)+f'(t)(u-t).
\end{equation}
原问题由此变为关于\(t\)的方程。\((u,v)\)未必是切点，不能直接用\(f'(u)\)代替切线斜率。

\begin{example}
求过点\((0,-1)\)且与\(y=x^2\)相切的所有直线。再求过\((0,b)\)且与\(y=\ln x\)相切的直线，并讨论其条数。
\end{example}

设两条曲线为\(y=f(x)\)、\(y=g(x)\)，切点横坐标分别为\(s,t\)。令两条切线的斜率与截距分别相同，得到公共切线的条件：
\begin{equation*}
 f'(s)=g'(t),\qquad f(s)-sf'(s)=g(t)-tg'(t).
\end{equation*}
两切点可以不同。若要求在公共点相切，还须\(s=t\)且\(f(s)=g(t)\)。几何约束可再联用第四章的曲线方程与向量方法。

\begin{exercise}
\mbox{}\par\nobreak

（1）求\(y=e^x\)上平行于\(y=2x\)的切线。

（2）求\(y=x^2\)与\(y=-(x-2)^2\)的公共切线。

（3）求实数\(a\)，使\(y=\ln x\)与\(y=ax^2\)有公共点，且在该公共点处切线相同。
\end{exercise}

\subsection{最优化问题}

导数不能替代建模。先确定变化的量、约束及允许范围，再把目标量表示成单变量函数，才能使用最值方法。最后还应把候选点还原到原问题，核对它是否满足约束。

例如，用长为\(L\)的材料围成矩形，一边为\(x\)，另一边为\(L/2-x\)，面积为\(A(x)=x(L/2-x)\)，其中\(0<x<L/2\)。由\(A'(x)=L/2-2x\)，面积先增后减，在\(x=L/4\)处取得最大值。这个结论也可以用配方或均值不等式得到；导数把同一种研究推进到更一般的目标函数。

\begin{example}
在半径为\(R\)的圆内作内接矩形，求面积的最大值。若改为在半轴长为\(a,b\)的椭圆内作边平行于坐标轴的内接矩形，怎样结合伸缩变换与导数研究？
\end{example}

\begin{example}
从边长为\(L\)的正方形纸板四角各剪去一个边长为\(x\)的小正方形，再折成无盖长方体盒。建立体积关于\(x\)的函数，求体积最大时的\(x\)。
\end{example}

约束使目标量同时受到几个相反方向的影响：盒子变高时，底面也在缩小。把它们写进同一个函数，才能比较总体变化。结果还应符合原量纲：所有长度放大为原来的\(k\)倍时，面积与体积分别放大为\(k^2\)、\(k^3\)倍。

\begin{exercise}
\mbox{}\par\nobreak

（1）制作一个体积为\(V>0\)的无盖圆柱容器，侧面与底面的单位面积成本相同。求用料最少时的底面半径与高的关系。

（2）在曲线\(y=1/x\)的第一象限部分上取点\(P\)，使\(P\)到原点的距离最小。求\(P\)，并解释为什么可以改为研究距离的平方。

（3）设\(a>0\)，点\(A=(0,a)\)、\(B=(1,0)\)。在\(x\)轴上取\(P=(t,0)\)，\(0\le t\le1\)，求\(|AP|+\tfrac12|PB|\)的最小值。驻点不在研究区间内时，应怎样完成判断？
\end{exercise}
\endgroup

\FloatBarrier
\newpage
\section{高阶导数}
\label{sec:calculus-higher-derivatives}

\begingroup
\setlength{\parskip}{4pt}
\setlength{\abovedisplayskip}{10pt}
\setlength{\belowdisplayskip}{10pt}

一阶导数描述函数的变化率。继续研究导函数的变化，可以判断图像怎样弯曲，也可以把一次近似推进为更高次的局部模型。下面先研究高阶导数和凸性，再在选讲部分建立Taylor公式。

\subsection{高阶导数}

\begin{definition}
若导函数\(f'\)在某点可导，称其导数为\(f\)在该点的{\heiti{}二阶导数}，记为\(f''\)或\(d^2y/dx^2\)。依次求导得到\(n\)阶导数，记为\(f^{(n)}\)。并约定\(f^{(0)}=f\)。
\end{definition}

若\(s(t)\)表示位置，则\(s'(t)\)表示速度，\(s''(t)\)表示速度的变化率，即加速度。加速度为正，说明速度随时间增大，并不单独说明位置正在增大；位置增减由速度符号决定。

对多项式，反复求导最终得到常数与零；对指数函数，各阶导数保持同一形式；正弦、余弦则按四步循环变化。例如
\begin{equation*}
 (e^x)^{(n)}=e^x,\qquad
 \frac{d^n}{dx^n}\sin x=\sin\left(x+\frac{n\pi}2\right).
\end{equation*}
这些结构将使高次近似中的系数能够统一计算。

\begin{theorem}[Leibniz公式]\label{thm:calculus-leibniz}
若\(f,g\)具有直到\(n\)阶的导数，则
\begin{equation}\label{eq:calculus-leibniz}
 (fg)^{(n)}=\sum_{k=0}^n\binom nk f^{(k)}g^{(n-k)}.
\end{equation}
\end{theorem}

\begin{proof}
\(n=1\)时就是乘积法则。若公式对\(n\)成立，将各项分别求导。同一个\(f^{(k)}g^{(n+1-k)}\)由相邻两项产生，系数相加为\(\binom n{k-1}+\binom nk=\binom{n+1}k\)，两端项的系数仍为\(1\)，故公式对\(n+1\)成立。由归纳法得证。
\end{proof}

例如，\((fg)''=f''g+2f'g'+fg''\)。反复求导时，每次可以选择对哪个因子求导；\(n\)次中有\(k\)次落在第一个因子上的选择数，正是\(\binom nk\)。若一个因子为多项式，它的高阶导数最终为零，和式便只留下有限个非零项。

\begin{example}
用Leibniz公式求\(x^2e^x\)与\(x\sin x\)的\(n\)阶导数，并与前两阶的直接求导核对。
\end{example}

\begin{example}
求\(f(x)=x^4-2x^2\)的各阶导数。再求\(g(x)=\ln x\)的\(n\)阶导数，观察系数与阶乘的联系。
\end{example}

\begin{example}[2014年四川卷理科第21题第（1）问]
\label{ex:calculus-sichuan-derivatives}
设\(f(x)=e^x-ax^2-bx-1\)，\(a,b\in\mathbb R\)，\(g=f'\)。求\(g\)在\([0,1]\)上的最小值。
\end{example}

研究\(f'\)的单调性，可以继续考察\(f''\)的符号；再由\(f'\)的变化判断其零点与正负，回到\(f\)的单调性和极值问题。

\subsection{凸函数}
\label{subsec:calculus-convexity}

单调性比较不同位置的高低，凸性则比较图像与弦的位置。

\begin{definition}
设\(I\)为区间。若对任意\(u,v\in I\)及\(0\le\lambda\le1\)，都有
\begin{equation}\label{eq:calculus-convex-definition}
 f((1-\lambda)u+\lambda v)
 \le(1-\lambda)f(u)+\lambda f(v),
\end{equation}
称\(f\)为\(I\)上的{\heiti{}凸函数}；若不等号反向，则称为{\heiti{}凹函数}。
\end{definition}

右端是两点连线上相应位置的高度，所以本章所说的凸，指图像不高于弦。凹函数的图像则不低于弦。这个约定也明确了后面二阶导数判据的方向。由三角不等式，\(|x|\)满足式\eqref{eq:calculus-convex-definition}，虽在原点不可导，仍是凸函数。

对\(u<t<v\)，凸性可以等价地写为
\begin{equation}\label{eq:calculus-convex-slopes}
 \frac{f(t)-f(u)}{t-u}\le\frac{f(v)-f(t)}{v-t}.
\end{equation}
因为两边乘以正的分母，整理后正好是\(t\)处图像不高于弦的条件。

\begin{theorem}\label{thm:calculus-convex-derivative}
区间内可导函数\(f\)为凸函数，当且仅当\(f'\)单调不减。若\(f\)二阶可导，这又等价于\(f''\ge0\)。
\end{theorem}

\begin{proof}
若\(f\)凸，利用弦上中间点的比较，再分别让中间点趋于\(u\)、\(v\)，得到
\begin{equation*}
 f'(u)\le\frac{f(v)-f(u)}{v-u}\le f'(v).
\end{equation*}
故\(f'\)单调不减。

反过来，若\(f'\)单调不减，对\([u,t]\)、\([t,v]\)分别使用中值定理，两个割线斜率为\(f'(\xi)\)、\(f'(\eta)\)，其中\(\xi<t<\eta\)。由\(f'(\xi)\le f'(\eta)\)，得到式\eqref{eq:calculus-convex-slopes}，故\(f\)凸。最后一个等价关系，是把单调性判据应用于\(f'\)。
\end{proof}

由上述割线斜率与导数的比较，还得到凸函数的{\heiti{}切线不等式}：
\begin{equation}\label{eq:calculus-convex-tangent}
 f(x)\ge f(a)+f'(a)(x-a).
\end{equation}
凸函数的图像不低于任意一条切线；凹函数则反向。弦与切线由此给出两个方向的估计。

例如，\((e^x)''>0\)，在\(a=0\)处使用切线不等式，得到\(e^x\ge1+x\)；\((\ln x)''<0\)，在\(a=1\)处则得到\(\ln x\le x-1\)。这些不等式也可以通过研究函数之差证明，凸性把它们组织为同一种几何关系。

\begin{figure}[!ht]
\centering
\includegraphics[width=.43\linewidth]{calculus-convexity}
\caption{凸函数的图像不低于切线\(\ell\)，在弦的两个端点之间又不高于弦\(s\)。}
\label{fig:calculus-convexity}
\end{figure}

若连续图像在一点两侧的凸、凹性质发生改变，称这个点为{\heiti{}拐点}。对二阶导数存在的函数，可用其在两侧的符号判断。\(x^3\)在原点有拐点，而\(x^4\)虽然在原点二阶导数为零，两侧却同为凸，因而没有拐点。

\begin{example}
研究\(f(x)=x^4-2x^2\)的凸性和拐点，结合前面的单调性与极值画出图像。再用凸性证明，对任意实数\(u,v\)，有
\begin{equation*}
 e^{(u+v)/2}\le\frac{e^u+e^v}{2}.
\end{equation*}
\end{example}

\begin{example}[2023年新课标II卷第22题第（1）问]
证明当\(0<x<1\)时，\(x-x^2<\sin x<x\)。分别研究相应函数的导数或凸性，比较两种证明所控制的量。
\end{example}

\FloatBarrier
\setlength{\parskip}{6pt}
\begin{exercise}
\mbox{}\par\nobreak

（1）求\(e^{2x}\)、\(1/x\)和\(\cos x\)的\(n\)阶导数。

（2）设\(s(t)=t^3-3t^2\)，\(t\ge0\)。分别讨论位置、速度的增减，判断何时运动方向改变。

（3）求\(f(x)=x e^x\)的\(n\)阶导数，并用归纳法验证所得表达式。

（4）判断\(x^2\)、\(\sqrt x\)和\(1/x\)在正半轴上的凸性，并写出\(x=1\)处的切线不等式。

（5）利用\(\ln x\)的凹性证明\(\sqrt{uv}\le(u+v)/2\)，其中\(u,v>0\)。

（6）求\(f(x)=x e^{-x}\)的拐点，比较极值点与拐点的位置。

（7）设\(f''\ge0\)。证明任意两点之间图像不高于弦，并说明只知道某一点的\(f''\)为正，为什么不足以判断整个区间的凸性。

（8）（2014年四川卷理科第21题第（2）问）对例\ref{ex:calculus-sichuan-derivatives}中的\(f\)，若\(f(1)=0\)，且\(f\)在\((0,1)\)内有零点，求\(a\)的取值范围。

（9）（2024年北京卷第20题第（3）问）设\(f(x)=x+\ln(1+x)\)，\(A=(t,f(t))\)、\(C=(0,f(t))\)，其中\(t>0\)。曲线在\(A\)处的切线与\(y\)轴交于\(B\)，\(O\)为原点。是否存在点\(A\)，使\(2S_{\triangle ACO}=15S_{\triangle ABO}\)？若存在，这样的点有几个？

参考数据：\(1.09<\ln3<1.10\)，\(1.60<\ln5<1.61\)，\(1.94<\ln7<1.95\)。

（10）（2014年新课标I卷理科第21题）设\(f(x)=ae^x\ln x+be^{x-1}/x\)，曲线在\((1,f(1))\)处的切线为\(y=e(x-1)+2\)。

（i）求\(a,b\)；（ii）证明\(f(x)>1\)。

（11）（2016年全国I卷理科第21题）\(f(x)=(x-2)e^x+a(x-1)^2\)有两个零点。

（i）求\(a\)的取值范围；（ii）设两个零点为\(x_1,x_2\)，证明\(x_1+x_2<2\)。

（12）（2016年全国III卷理科第21题）设\(f(x)=\alpha\cos2x+(\alpha-1)(\cos x+1)\)，其中\(\alpha>0\)，记\(|f(x)|\)的最大值为\(A\)。

（i）求\(f'(x)\)；（ii）求\(A\)；（iii）证明\(|f'(x)|\le2A\)。

（13）设\(f\)在\([a,b]\)上二阶可导，且\(f''\ge m>0\)。对\(g(x)=f(x)-mx^2/2\)使用凸性，证明\(\dfrac{f(a)+f(b)}2-f\!\left(\dfrac{a+b}2\right)\ge\dfrac{m(b-a)^2}8\)。

（14）设\(f\)在区间内二阶可导，且\(f''>0\)。证明每条切线与该区间内的图像只有切点一个公共点。

（15）求\(x^3e^{2x}\)的\(n\)阶导数。若\(P\)为\(m\)次多项式，写出\((Pe^{ax})^{(n)}\)的一般表达式，说明为什么至多需要\(m+1\)项。

（16）对恒等式\((1+x)^p(1+x)^q=(1+x)^{p+q}\)求\(n\)阶导数，并在零点取值，证明\(\sum_{k=0}^n\binom pk\binom q{n-k}=\binom{p+q}n\)，其中\(p,q\)为非负整数，超过上标的二项式系数约定为零。
\end{exercise}
\endgroup

\FloatBarrier
\newpage
\begingroup
\setlength{\parskip}{4pt}
\setlength{\abovedisplayskip}{10pt}
\setlength{\belowdisplayskip}{10pt}
\subsection{高阶微分*}

一阶微分将函数增量中关于输入增量的线性部分写成\(dy=f'(x)\,dx\)。固定自变量增量\(dx\)，再对这个表达式中的系数求微分，可以依次得到高阶微分。

\begin{definition}
设\(x\)为独立变量，\(y=f(x)\)具有\(n\)阶导数。在连续求微分的过程中保持\(dx\)不变，定义
\begin{equation*}
 d^2y=d(dy)=f''(x)(dx)^2,\qquad
 d^ny=f^{(n)}(x)(dx)^n.
\end{equation*}
\end{definition}

高阶导数给出系数，高阶微分还包括输入增量的相应次幂。符号\(d^2y\)表示对\(y\)连续求两次微分，并非\((dy)^2\)；\(d^2y/dx^2\)则与二阶导数的记号一致。后面的Taylor公式将说明，这些系数怎样共同组成更高次的局部近似。

一阶微分具有形式不变性，但在高阶情形下，中间变量本身的变化也需要保留。设\(y=f(u)\)、\(u=u(x)\)，则
\begin{equation*}
 \begin{aligned}
 d^2y&=d\bigl(f'(u)\,du\bigr)
       =f''(u)(du)^2+f'(u)d^2u,\\
 du&=u'(x)\,dx,\qquad d^2u=u''(x)(dx)^2.
 \end{aligned}
\end{equation*}
这就是二阶链式法则
\begin{equation*}
 \frac{d^2y}{dx^2}=f''(u)\bigl(u'(x)\bigr)^2+f'(u)u''(x)
\end{equation*}
的微分形式。若把\(u\)直接作为独立变量，保持的是\(du\)，所得到的二阶微分为\(f''(u)(du)^2\)；若\(u\)随\(x\)非线性变化，则还要计入\(f'(u)d^2u\)。独立变量的选择在二阶运算中已经产生了实质影响。

\begin{example}
设\(y=u^2\)、\(u=x^2\)。先将\(y\)写成\(x\)的函数求\(d^2y\)，再用二阶链式法则计算，并分别指出两项的来源。
\end{example}

\begin{exercise}
\mbox{}\par\nobreak

（1）以\(x\)为独立变量，求\(y=\ln x\)的前三阶微分，并比较\(d^2y\)与\((dy)^2\)。

（2）设\(y=\sin u\)、\(u=x^2+1\)。分别计算\(f''(u)(du)^2\)与\(d^2y\)，指出只保留前者漏掉了什么。

（3）若\(u=ax+b\)，说明二阶微分公式中的附加项为什么消失，并由\(d^ny=f^{(n)}(u)a^n(dx)^n\)核对线性换元下的高阶求导。

（4）设\(y=e^u\)、\(u=\ln x\)，以\(x\)为独立变量，计算二阶微分公式中的两项，解释它们为何恰好抵消，并将\(y\)直接表示为\(x\)的函数进行核对。
\end{exercise}

\subsection{Taylor多项式*}

一次近似匹配函数在基点的值与一阶变化。若继续匹配二阶、三阶直至更高阶导数，是否可以得到更精确的局部模型？先在多项式中确定应当选择的系数。

设
\begin{equation*}
 P_n(x)=c_0+c_1(x-a)+\cdots+c_n(x-a)^n.
\end{equation*}
则\(P_n^{(k)}(a)=k!c_k\)。因此，使\(P_n^{(k)}(a)=f^{(k)}(a)\)对\(0\le k\le n\)都成立的多项式是唯一的。

\begin{definition}
若\(f\)在\(a\)处具有直到\(n\)阶的导数，称
\begin{equation}\label{eq:calculus-taylor-polynomial}
 T_n(x)=\sum_{k=0}^n\frac{f^{(k)}(a)}{k!}(x-a)^k
\end{equation}
为\(f\)在\(a\)处的\(n\)阶{\heiti{}Taylor多项式}。取\(a=0\)时，也称Maclaurin多项式。
\end{definition}

例如，\(e^x\)在\(0\)处的各阶导数均为\(1\)，故
\begin{equation*}
 T_n(x)=1+x+\frac{x^2}{2!}+\cdots+\frac{x^n}{n!}.
\end{equation*}
\(\sin x\)的四阶Taylor多项式为\(x-x^3/6\)，因为常数项、二次项和四次项的系数均为零。所谓阶数，并不要求最高次项系数非零。

到这里，已经确定了匹配各阶导数的多项式，但尚未证明它的误差有多小。令\(R_n(x)=f(x)-T_n(x)\)，接下来的任务就是控制这个余项。

\begin{figure}[!ht]
\centering
\includegraphics[width=.46\linewidth]{calculus-taylor}
\caption{\(e^x\)在零点的前三阶Taylor多项式。各阶局部近似的精度，仍须由余项估计确定。}
\label{fig:calculus-taylor}
\end{figure}

\begin{exercise}
\mbox{}\par\nobreak

（1）写出\(\cos x\)、\(\ln(1+x)\)及\((1+x)^\alpha\)在\(0\)处的三阶Taylor多项式。

（2）求\(x^4\)在\(a=1\)处的四阶Taylor多项式，直接展开验证它恰为原多项式。

（3）若两个次数不超过\(n\)的多项式在\(a\)处直到\(n\)阶的导数都相同，证明它们相等。
\end{exercise}

\subsection{Taylor公式*}

\begin{theorem}[带Lagrange余项的Taylor公式]\label{thm:calculus-taylor-lagrange}
设\(f\)在包含\(a\)与\(x\)的区间上具有直到\(n+1\)阶的连续导数，\(x\ne a\)。则存在位于\(a\)与\(x\)之间的\(\xi\)，使
\begin{equation}\label{eq:calculus-taylor-lagrange}
 f(x)=\sum_{k=0}^n\frac{f^{(k)}(a)}{k!}(x-a)^k
 +\frac{f^{(n+1)}(\xi)}{(n+1)!}(x-a)^{n+1}.
\end{equation}
\end{theorem}

为了得到余项，需要将函数与Taylor多项式的差，再减去一个系数待定的高次项，使所得函数在两个位置都为零。Taylor多项式已经匹配的各阶导数，则为反复使用Rolle定理提供了条件。

\begin{proof}
固定\(x\)，选取常数\(K\)，使
\begin{equation*}
 H(t)=f(t)-T_n(t)-K(t-a)^{n+1}
\end{equation*}
满足\(H(x)=0\)。又有\(H(a)=H'(a)=\cdots=H^{(n)}(a)=0\)。由\(H(a)=H(x)=0\)，得到两者之间的一个\(H'\)零点；将它与\(H'(a)=0\)配对，再得到更靠近\(a\)的\(H''\)零点。如此反复，共用\(n+1\)次Rolle定理，得到\(H^{(n+1)}(\xi)=0\)。

因此\(f^{(n+1)}(\xi)-(n+1)!K=0\)。代回\(H(x)=0\)，即得公式。
\end{proof}

若\(|f^{(n+1)}|\le M\)，则
\begin{equation}\label{eq:calculus-taylor-bound}
 |R_n(x)|\le\frac{M}{(n+1)!}|x-a|^{n+1}.
\end{equation}
这个界不需要知道\(\xi\)的具体位置，可以直接用于数值误差估计。

\paragraph{Peano余项} \mbox{}

有时只需要知道余项相对于某次幂趋零，而不要求给出明确的误差常数。记
\begin{equation*}
 r(h)=o(h^n)\quad(h\to0)
 \qquad\Longleftrightarrow\qquad
 \frac{r(h)}{h^n}\to0.
\end{equation*}
这称为相对于\(h^n\)的高阶无穷小，\(o\)在这里描述余项满足的极限关系。

\begin{theorem}[带Peano余项的Taylor公式]\label{thm:calculus-taylor-peano}
设\(f\)在\(a\)附近有直到\(n-1\)阶的导数，且\(f^{(n-1)}\)在\(a\)处可导，其中\(n\ge1\)。则
\begin{equation}\label{eq:calculus-taylor-peano}
 f(a+h)=\sum_{k=0}^n\frac{f^{(k)}(a)}{k!}h^k+o(h^n).
\end{equation}
\end{theorem}

\begin{proof}
\(n=1\)就是可导函数的一次近似。设结论对\(n-1\)成立，注意到
\begin{equation*}
 T_n'(a+h)=\sum_{j=0}^{n-1}\frac{f^{(j+1)}(a)}{j!}h^j
\end{equation*}
正好是\(f'\)在\(a\)处的\(n-1\)阶Taylor多项式。将归纳假设用于\(f'\)，便知余项\(R_n(h)=f(a+h)-T_n(a+h)\)满足\(R_n'(h)=o(h^{n-1})\)。

给定\(\varepsilon>0\)，充分小的非零\(t\)均满足\(|R_n'(t)|\le\varepsilon|t|^{n-1}\)。又\(R_n(0)=0\)，故由中值定理，存在\(\eta\)介于\(0\)与\(h\)之间，使
\begin{equation*}
 |R_n(h)|=|hR_n'(\eta)|\le\varepsilon|h||\eta|^{n-1}\le\varepsilon|h|^n.
\end{equation*}
所以\(R_n(h)=o(h^n)\)，归纳完成。
\end{proof}

Lagrange余项利用一段区间内更高阶导数的控制，给出可计算的误差；Peano余项只在基点附近比较阶数，所需条件也较弱。应根据问题需要选择版本，不能把一种版本的条件直接用于另一种余项表达。

\begin{example}
用\(e^x\)在\(0\)处的二阶Taylor公式估计\(e^{0.1}\)，并给出一个小于\(1/2000\)的误差上界。再写出\(\ln(1+h)\)的三阶Peano展开。
\end{example}

\begin{exercise}
\mbox{}\par\nobreak

（1）用\(\cos x\)在\(0\)处的Taylor公式证明\(|\cos x-1+x^2/2|\le x^4/24\)。

（2）设\(f(a)=f'(a)=0\)、\(f''(a)=2\)。写出\(f(a+h)\)的二阶Peano展开，并判断充分靠近\(a\)的非零\(h\)所对应函数值的符号。

（3）若\(r(h)=o(h^3)\)，判断\(r(h)/h^2\)、\(r(h)/h^3\)的极限能否确定。若仅知\(r(h)=o(h^2)\)，结论有何不同？
\end{exercise}

\subsection{Taylor公式的应用*}

有限展开的用途取决于余项的控制。求极限时，保留抵消后最先出现的有效项；估算时，使误差小于给定精度；判断函数值大小时，则要确定首个非零项的符号。

例如，
\begin{equation*}
 e^x=1+x+\frac{x^2}2+o(x^2),\qquad
 \cos x=1-\frac{x^2}2+o(x^2),
\end{equation*}
相减后可得
\begin{equation*}
 \lim_{x\to0}\frac{e^x-\cos x-x}{x^2}=1.
\end{equation*}
若只保留一次近似，两个余项都只是\(o(x)\)，便不足以确定除以\(x^2\)后的极限。所需阶数由抵消后的比较决定。

当\(f'(a)=0\)、\(f''(a)\ne0\)时，
\begin{equation*}
 f(a+h)-f(a)=\frac{f''(a)}2h^2+o(h^2).
\end{equation*}
充分小的非零\(h\)下，其符号与\(f''(a)\)相同。因此\(f''(a)>0\)给出严格极小值，\(f''(a)<0\)给出严格极大值。这是{\heiti{}二阶导数判别法}。

更一般地，若\(f'(a)=\cdots=f^{(m-1)}(a)=0\)，而\(f^{(m)}(a)\ne0\)，则首个非零项为\(f^{(m)}(a)h^m/m!\)。\(m\)为偶数时，它在两侧同号，决定极值；\(m\)为奇数时，两侧异号，因而不是极值点。这个判断把前面的\(x^2\)、\(x^3\)、\(x^4\)放在同一个局部模型中。

不等式也可以由余项符号得到。比如\(x>0\)时，对\(e^x\)取二阶展开，Lagrange余项为\(e^\xi x^3/6>0\)，于是
\begin{equation*}
 e^x>1+x+\frac{x^2}2.
\end{equation*}
这一结论比切线不等式多保留了一个正的二次贡献。

\begin{example}
求\(\lim_{x\to0}(\sin x-x\cos x)/x^3\)。再研究\(f(x)=\ln(1+x)-x+x^2/2\)在原点的极值情况，说明展开到哪一阶才足够判断。
\end{example}

\begin{example}[2023年新课标II卷第22题第（2）问]
设\(f(x)=\cos(ax)-\ln(1-x^2)\)。若\(x=0\)是\(f\)的极大值点，求\(a\)的取值范围。尝试使用有限阶Taylor展开，并处理二次项系数为零的情形。
\end{example}

\begin{exercise}
\mbox{}\par\nobreak

（1）求\(\lim_{x\to0}(\tan x-x)/x^3\)。可由正弦、余弦的有限展开得到，不必反复求\(\tan x\)的高阶导数。

（2）求\(\lim_{x\to0}\bigl((1+x)^\alpha-1-\alpha x\bigr)/x^2\)，其中\(\alpha\)为实常数。

（3）证明\(x>0\)时\(x-x^2/2<\ln(1+x)<x\)，分别比较中值定理、函数单调性与Taylor余项的作用。

（4）设\(f\)在\(a\)附近有直到四阶的连续导数，且\(f'(a)=f''(a)=f'''(a)=0\)、\(f^{(4)}(a)>0\)。证明\(a\)为严格极小值点。

（5）回看习题\ref{exs:calculus-parameter-gaokao}第（3）题的最后一问。先由连续性确定\(a\)，再令\(t=x-1\)。用\(\ln\frac{1+t}{1-t}\)在\(t=0\)处的三阶Taylor展开确定\(b\)的必要条件，再用余项符号证明其充分性，与原题的导数方法比较。
\end{exercise}

\subsection{迭代求根*}
\label{subsec:calculus-newton-iteration}

第一章用迭代数列近似求解方程，也讨论了平方根的迭代。现在可以利用切线，从一般函数的局部一次模型产生新的取值。

设要求\(f(x)=0\)的根，在近似值\(x_n\)处用切线代替曲线。令切线的纵坐标为零，得到
\begin{equation}\label{eq:calculus-newton-iteration}
 x_{n+1}=x_n-\frac{f(x_n)}{f'(x_n)}.
\end{equation}
这称为{\heiti{}Newton迭代}。例如，\(f(x)=x^2-A\)、\(A>0\)时，公式化为\(x_{n+1}=(x_n+A/x_n)/2\)，就是第一章的平方根迭代。切线模型说明了这个递推式的来源，但迭代能否一直进行、是否收敛，还需要证明。

\begin{proposition}\label{prop:calculus-newton-error}
设\(f\)在根\(\alpha\)附近二阶连续可导，且\(f'(\alpha)\ne0\)。则在足够小的邻域内，从任意初值出发的Newton迭代都趋于\(\alpha\)，并有
\begin{equation*}
 |x_{n+1}-\alpha|\le K|x_n-\alpha|^2,
\end{equation*}
其中\(K\)为与\(n\)无关的常数。
\end{proposition}

\begin{proof}
由连续性，在\(\alpha\)的一个闭邻域内可以取\(m>0\)、\(M>0\)，使\(|f'|\ge m\)、\(|f''|\le M\)。在\(x\)处展开并取终点为\(\alpha\)，Taylor公式给出
\begin{equation*}
 0=f(\alpha)=f(x)+f'(x)(\alpha-x)+\frac12f''(\xi)(\alpha-x)^2.
\end{equation*}
所以，对\(N(x)=x-f(x)/f'(x)\)，有
\begin{equation*}
 |N(x)-\alpha|\le\frac M{2m}|x-\alpha|^2.
\end{equation*}
取\(K=M/(2m)\)，再缩小半径\(\delta\)，使\(K\delta\le1/2\)。当\(|x-\alpha|\le\delta\)时，下一个点仍在该邻域内，且误差至多为原来的一半。归纳可知迭代始终有定义并趋于\(\alpha\)，上式同时保留了更精细的平方误差界。
\end{proof}

这个证明先保证点留在可以控制的范围内，再判断误差怎样缩小。平方误差意味着足够接近根以后，误差越小，下一步相对于当前误差的比例也越小；它比固定比例的收缩更快。结论是局部的，初值远离所求根时仍应重新研究迭代过程。

在上述邻域内，若\(K\le1\)，且某一步误差不超过\(10^{-3}\)，则后两步误差分别不超过\(10^{-6}\)、\(10^{-12}\)。平方误差界由此说明了接近单根后近似精度的提高。

\newpage
\paragraph{有限层分式} \mbox{}

求\(\sqrt2\)还可以从恒等式\(y=1+1/(1+y)\)出发，取
\begin{equation*}
 y_1=1,\qquad y_{n+1}=1+\frac1{1+y_n}.
\end{equation*}
逐次代入后，每个\(y_n\)都表示为有限层分式。由\(\varphi(y)=1+1/(1+y)\)将\([1,3/2]\)映入自身，且其上\(|\varphi'|\le1/4\)，中值定理给出
\begin{equation*}
 |y_{n+1}-\sqrt2|\le\frac14|y_n-\sqrt2|.
\end{equation*}
因此它也收敛，但这里首先得到的是固定比例的误差界。对这个例子，研究连分式的有限层截断，就是研究相应的迭代数列；收敛性和精度仍由明确的估计保证。

\begin{example}
分别用Newton迭代和上述分式迭代近似\(\sqrt2\)。从\(1\)开始各算三步，比较误差，再利用第一章平方根迭代的精确误差式及本小节的比例界估计下一步误差。
\end{example}

\begin{exercise}
\mbox{}\par\nobreak

（1）为\(x^3=A\)、\(A>0\)写出Newton迭代。若初值\(x_1\ge A^{1/3}\)，证明数列不低于所求根且单调不增，因而收敛。再研究误差怎样变化。

（2）用\(f(x)=\ln x-1\)写出求\(e\)的Newton迭代，从\(x_1=2\)计算前三个近似值。解释一次模型使用的基点为什么每一步都要更新。

（3）对\(f(x)=x^2\)，在非零初值处写出Newton递推式。求其误差比例，比较它与命题\ref{prop:calculus-newton-error}中的条件。

（4）对正文的分式迭代，证明
\begin{equation*}
 y_{n+1}-\sqrt2=-\frac{y_n-\sqrt2}{(1+y_n)(1+\sqrt2)}.
\end{equation*}
由此判断相邻近似值分居根的两侧，并给出使误差小于\(10^{-4}\)的一个充分迭代次数。
\end{exercise}

\subsection{从有限展开到无穷展开*}

有限阶Taylor公式对每个固定的阶数给出一个带余项的等式。要使多项式的次数不断增大后恢复原函数，还需要证明余项随阶数趋零。这里改变的是阶数\(n\)，不再是固定阶数下的增量\(h\)。

例如，对固定实数\(x\)，\(e^x\)的余项满足
\begin{equation*}
 |R_n(x)|\le e^{|x|}\frac{|x|^{n+1}}{(n+1)!}\longrightarrow0.
\end{equation*}
最后的数列相邻项之比为\(|x|/(n+2)\to0\)，所以趋零。因此
\begin{equation}\label{eq:calculus-exponential-series}
 e^x=\sum_{k=0}^{\infty}\frac{x^k}{k!}.
\end{equation}
在\(x=1\)处，它与第一章的\(e\)的级数表达一致。

正弦、余弦的各阶导数绝对值不超过\(1\)，同样由阶乘控制余项，得到
\begin{equation*}
 \sin x=\sum_{k=0}^{\infty}\frac{(-1)^kx^{2k+1}}{(2k+1)!},\qquad
 \cos x=\sum_{k=0}^{\infty}\frac{(-1)^kx^{2k}}{(2k)!}.
\end{equation*}
这些等式在每个固定实数\(x\)处成立。每个等式右边首先是一项项数构成的级数，第一章的收敛性工具仍然适用。

\paragraph{导数值是否总能恢复函数} \mbox{}

仅有各阶导数还不够。考察
\begin{equation*}
 f(x)=\begin{cases}e^{-1/x^2},&x\ne0,\\0,&x=0.\end{cases}
\end{equation*}
这个函数在原点的各阶导数均为零，故所有Taylor多项式都是零，但任意非零位置的函数值都为正。

为验证上述导数结论，先注意\(t^p e^{-t}\to0\)（\(t\to+\infty\)），其中\(p>0\)。由\(e^t\)的有限Taylor公式，对任意整数\(m>p\)有\(e^t\ge t^m/m!\)，故\(0\le t^pe^{-t}\le m!t^{p-m}\to0\)。所以\(e^{-1/x^2}\)乘以任意固定次的\(1/x\)，仍在\(x\to0\)时趋零。

在\(x\ne0\)处，每次求导都得到\(P(1/x)e^{-1/x^2}\)的形式，其中\(P\)为多项式。若前一阶导数在原点为零，其差商仍是这种形式乘以\(1/x\)，极限为零。归纳即得各阶导数存在且均为零。

这个例子说明，匹配基点处任意多阶导数，与在基点之外恢复函数，是不同的要求。前者由导数值确定，后者还取决于随阶数变化的余项。

\begin{exercise}
\mbox{}\par\nobreak

（1）分别写出\(e+e^{-1}\)与\(e-e^{-1}\)的级数表达，说明奇偶次项的作用。

（2）利用余项界，选取足够的Taylor多项式，使\(\sin(1/2)\)的近似误差小于\(10^{-6}\)。

（3）对上述\(f(x)=e^{-1/x^2}\)在原点补零的函数，证明对每个正整数\(n\)，均有\(f(x)=o(x^n)\)。解释这与\(f\)在原点附近不恒为零为何相容。
\end{exercise}
\endgroup

\FloatBarrier
\newpage
\section{函数研究的方法}
\label{sec:calculus-function-comparison}

\begingroup
\setlength{\parskip}{3.5pt}

前面已经由导数建立单调性、极值和凸性的判据。实际问题还需要确定研究哪个函数：题设可能给出两个数之间的关系，也可能只说明某个点满足一个方程；所求的量未必能直接算出。本节从这些关系出发，选择研究对象，利用已有性质完成比较。

\subsection{选择研究函数}

若\(H\)在区间上严格递增，则\(H(u)<H(v)\)与\(u<v\)等价，\(H(u)=H(v)\)则给出\(u=v\)。因此，有些等式或不等式可以先整理成同一个函数在两处的取值，再研究这个函数的单调性。

例如，\(u,v>0\)时，关系\(u+\ln u=v+\ln v\)可以写成\(H(u)=H(v)\)，其中\(H(x)=x+\ln x\)。由\(H'(x)=1+1/x>0\)，立即得到\(u=v\)。若各项分别比较，反而容易丢失整个表达式的关系。

比较幂也可以转为比较函数值。对\(u,v>0\)，取对数后，\(u^v<v^u\)等价于\(\ln u/u<\ln v/v\)。因此应研究\(H(x)=\ln x/x\)，其导数为\((1-\ln x)/x^2\)。这个函数在\((0,e)\)上递增、在\((e,+\infty)\)上递减；比较两个幂，便转为确定\(u,v\)所在的区间及相应函数值。

有时题设直接给出的是导数的组合。反过来观察乘积法则、商法则与链式法则，有
\begin{equation*}
 \begin{aligned}
 (e^{-ax}f(x))'&=e^{-ax}(f'(x)-af(x)),\\
 \left(\frac{f(x)}x\right)'&=\frac{xf'(x)-f(x)}{x^2},\\
 (\ln|f(x)|)'&=\frac{f'(x)}{f(x)}.
 \end{aligned}
\end{equation*}
这些函数可以沿求导公式反向寻找。若取\(H=wf\)，则\(H'=wf'+w'f\)；要使它包含\(f'-af\)，便希望\(w'=-aw\)，而\(w=e^{-ax}\)恰好满足。类似地，商的导数自然产生\(xf'-f\)，所以研究\(f/x\)。各式均在表达式有意义的区间上使用。

\begin{example}
设\(f(0)=1\)，且\(x\ge0\)时\(f'(x)\ge f(x)\)。证明\(f(x)\ge e^x\)。再设\(x>0\)时\(xf'(x)-f(x)>0\)，比较\(f(2x)\)与\(2f(x)\)。分别说明为什么选择乘积或商。
\end{example}

若新函数仍不容易判号，可以继续求导，再由已知点逐级回推。例如，若\(H(a)=H'(a)=0\)，且\(H''\ge0\)于\([a,b]\)，则先有\(H'\ge0\)，继而有\(H\ge0\)。初值确定了每一步比较的起点。

\begin{example}
比较\(e^\pi\)与\(\pi^e\)的大小。再设\(e<a<b\)，比较\(a^b\)与\(b^a\)。
\end{example}

\newpage
\begin{exercise}
\mbox{}\par\nobreak

\begin{samepage}
（1）设\(u,v\in\mathbb R\)，且\(u^3+u=-(v^3+v)\)，证明\(u+v=0\)。
\par
\end{samepage}

\begin{samepage}
（2）设\(f\)在\([0,1]\)上连续可导，\(f(0)=2\)，且\(f'+2f\le0\)。估计\(f(x)\)，并说明比较方向怎样由所选函数的导数确定。
\par
\end{samepage}

（3）设\(f>0\)，且\(f'/f\le1\)于\([a,b]\)。证明\(f(b)/f(a)\le e^{b-a}\)，比较研究\(\ln f\)与\(e^{-x}f\)两条路线。

（4）设\(H(0)=H'(0)=H''(0)=0\)，且\(H'''\ge0\)于\([0,1]\)。证明\(H\ge0\)。再将这一推理用于\(e^x-1-x-x^2/2\)，指出每次回推所需的零点信息。
\end{exercise}

\subsection{利用方程中的关系}

方程的根即使难以显式求出，所满足的等式仍能参与运算。把这个等式代入待研究的表达式，往往可以减少未知量，求出目标值或确定它的范围。

考察\(f_a(x)=e^x+x^2-ax\)。其导函数\(e^x+2x-a\)严格递增，两端分别趋于负无穷和正无穷，所以有唯一零点\(c\)，且\(c\)为原函数的最小值点。虽然不便显式写出\(c\)，仍有
\begin{equation*}
 a=e^c+2c,\qquad
 f_a(c)=e^c+c^2-c(e^c+2c)=(1-c)e^c-c^2.
\end{equation*}
原来的函数值含有\(a,c\)，利用驻点方程后只剩\(c\)。若还需限定\(a\)的范围，则应由\(a=e^c+2c\)及其单调性，把\(c\)的范围对应回去。这种设出而不显式求出的零点，常称为隐零点。

对于两个未知点，方程还可以同时约束它们的和与积。若\(c_1,c_2\)是二次方程的根，研究对称式时可以先使用韦达关系；研究单个函数值时，也可以用方程替换其中的二次项。第四章在曲线联立中保留根与系数关系的做法，在这里同样适用。

\begin{example}
设\(f_a(x)=e^x+x^2-ax\)，最小值点为\(c\)。确定\(c>0\)、\(c=0\)、\(c<0\)分别对应的\(a\)的范围，再证明最小值不超过\(1\)，并确定等号条件。比较直接使用\(f_a(0)\)与研究化简后的\(f_a(c)\)两种方法。
\end{example}

保留方程以后，还要注意所设点的范围。导函数的一个零点是否为极值点，取决于两侧符号；所得参数表达式是否可以反过来确定这个点，则取决于相应函数的单调性。方程、位置与函数变化共同组成了推理的依据。

\(G_a(x)=\ln x-ax^2+x\)、\(a>0\)的唯一驻点\(c\)满足\(2ac^2=c+1\)，最大值为\(G_a(c)=\ln c+(c-1)/2\)。后者随\(c\)递增，而\(a=(c+1)/(2c^2)\)随\(c\)递减，故最大值随\(a\)递减。

若要求\(G_a(x)\le M\)对所有\(x>0\)成立，就等价于\(\ln c+(c-1)/2\le M\)。先确定允许的\(c\)，再由驻点关系还原参数范围，恒成立问题便落到一个最值点上。

\begin{exercise}
\mbox{}\par\nobreak

（1）设\(c\)为\(e^x+x-2=0\)的根，证明\(0<c<1\)。利用\(e^c=2-c\)，将\(e^{2c}+2ce^c+c^2\)化为确定的数。

（2）设\(f(x)=x(\ln x-ax)\)有两个极值点\(c_1<c_2\)。由驻点方程化简\(f(c_i)\)，并研究如何利用\(c_i\)的位置估计这两个极值。

（3）设\(a>0\)，函数\(f(x)=-\ln x-ax^2+x\)有两个极值点\(c_1,c_2\)。先确定\(a\)的范围，再用韦达关系将\(f(c_1)+f(c_2)\)表示为\(a\)的函数。

（4）对正文中的\(f_a\)，若\(a\)只在\([1,e+2]\)内变化，确定最小值点的范围。说明先求\(c\)的范围，再研究\((1-c)e^c-c^2\)，怎样避免同时保留两个独立变量。
\end{exercise}

\subsection{减少变量}

表达式中出现两个字母，未必就有两个独立变化的量。若它们满足一个方程，可以先利用约束；若没有这样的约束，则可以固定其中一个，将另一个作为自变量研究。变量的选择需要服从题设中的关系。

例如，设\(0<u<v\)，且\(u-\ln u=v-\ln v\)，则
\begin{equation}\label{eq:calculus-equal-level-log}
 v-u=\ln\frac vu.
\end{equation}
若选择比值\(t=v/u>1\)，由\(u(t-1)=\ln t\)，得到
\begin{equation*}
 u=\frac{\ln t}{t-1},\qquad v=\frac{t\ln t}{t-1}.
\end{equation*}
若选择差\(d=v-u>0\)，则\(v=e^du\)，从而
\begin{equation*}
 u=\frac d{e^d-1},\qquad v=\frac{de^d}{e^d-1}.
\end{equation*}
两种选择都利用了\(v-u=\ln(v/u)\)：比值确定了差，差也确定了比值。两种表示满足\(t=e^d\)；将所求的和、积等表达式代入后，都可以转为单变量问题。

参数有时也可以通过换元从等值关系中消去。设\(k>0\)、\(F_k(x)=x-k\ln x\)，令\(x=kt\)，则
\begin{equation*}
 F_k(kt)=k\bigl(t-\ln t-\ln k\bigr).
\end{equation*}
若\(F_k(u)=F_k(v)\)，令\(u=kr\)、\(v=ks\)，即得\(r-\ln r=s-\ln s\)。参数只引起平移、伸缩；再由\(u+v=k(r+s)\)、\(uv=k^2rs\)还原两点的和、积。

若只给出\(0<u<v\)，则不能擅自加入式\eqref{eq:calculus-equal-level-log}这样的等式。此时可以固定\(u\)，研究\(v\)怎样变化。例如，令\(H(v)=\ln(v/u)-(v-u)/u\)，有\(H(u)=0\)、\(H'(v)=1/v-1/u<0\)，故\(\ln(v/u)<(v-u)/u\)。这里仍需保证所得结论对每个允许的\(u\)成立。

\begin{example}
设\(u<v\)且\(e^u-u=e^v-v\)。令\(d=v-u\)，把\(e^u,e^v\)表示为\(d\)的函数，再研究\(u+v\)。比较差值与比值在本题中的便利程度。
\end{example}

\begin{exercise}
\mbox{}\par\nobreak

（1）对满足式\eqref{eq:calculus-equal-level-log}的\(u<v\)，分别用\(t\)、\(d\)表示\(u+v\)与\(uv\)，核对其对应关系。

（2）对\(0<u<v\)，分别固定\(u\)或令\(t=v/u\)，证明\((v-u)/v<\ln(v/u)<(v-u)/u\)。

（3）设\(0<u<v\)、\(ue^{-u}=ve^{-v}\)。取对数后说明它与式\eqref{eq:calculus-equal-level-log}的关系。

（4）设\(u,v>0\)、\(u+v=s\)。用消元、对称性或对数凹性求\(\ln u+\ln v\)的最大值。
\end{exercise}

\subsection{比较零点的位置}

设函数先严格递减、后严格递增，在\(c\)处取得最小值，而\(u<c<v\)满足\(f(u)=f(v)\)。要比较\(u+v\)与\(2c\)，比较\(u\)与反射点\(2c-v\)。若反射点位于左侧单调区间，原来的数值比较就转成\(f(u)\)与\(f(2c-v)\)的比较；再用\(f(u)=f(v)\)，两者可统一为一个变量。

以\(f(x)=x-\ln x\)为例，最小值点为\(1\)。对\(0<t<1\)，考察
\begin{equation*}
 R(t)=f(1+t)-f(1-t)
 =2t-\ln\frac{1+t}{1-t}.
\end{equation*}
有\(R(0)=0\)、\(R'(t)=-2t^2/(1-t^2)<0\)，故右侧等距离位置的函数值更小。若\(1<v<2\)，则\(f(u)=f(v)<f(2-v)\)，由左侧严格递减，得到\(u>2-v\)；若\(v\ge2\)，则\(u+v>2\)直接成立。因此两根的中点在\(1\)的右侧。

另一条路线是直接利用两根满足的方程。为此先建立一个关于正数的平均值关系。

\begingroup
\setlength{\parskip}{3pt}
\setlength{\abovedisplayskip}{8pt}
\setlength{\belowdisplayskip}{8pt}
\begin{proposition}\label{prop:calculus-logarithmic-mean}
若\(0<a<b\)，其对数平均值满足
\begin{equation}\label{eq:calculus-logarithmic-mean}
 \sqrt{ab}<\frac{b-a}{\ln b-\ln a}<\frac{a+b}{2}.
\end{equation}
\end{proposition}

\begin{samepage}
\begin{proof}
令\(q=b/a>1\)。对\(t>1\)，函数\(t-t^{-1}-2\ln t\)在\(1\)处为零，导数为\((t-1)^2/t^2>0\)。取\(t=\sqrt q\)，得到\(\ln q<(q-1)/\sqrt q\)，即左侧不等式。

函数\(\ln q-2(q-1)/(q+1)\)在\(1\)处为零，导数为\((q-1)^2/[q(q+1)^2]>0\)，得到右侧不等式。
\end{proof}
\end{samepage}
\endgroup
\newpage
若\(u-\ln u=v-\ln v\)，式\eqref{eq:calculus-equal-level-log}使对数平均值恰为\(1\)。于是同时得到\(uv<1\)、\(u+v>2\)。取\(a=e^\alpha,b=e^\beta\)，还可将式\eqref{eq:calculus-logarithmic-mean}写成
\begin{equation*}
 e^{(\alpha+\beta)/2}
 <\frac{e^\beta-e^\alpha}{\beta-\alpha}
 <\frac{e^\alpha+e^\beta}{2}\qquad(\alpha<\beta).
\end{equation*}
这些关系把两点函数值之差、两点间的平均变化与几种平均值联系起来。后面建立积分以后，还可以从凸函数的积分估计重新理解它们。

\begin{example}
设\(f(x)=xe^{-x}\)，\(u<v\)且\(f(u)=f(v)>0\)。先用反射比较证明\(u+v>2\)，再分别用比值换元与式\eqref{eq:calculus-logarithmic-mean}证明\(u+v>2\)、\(uv<1\)。比较三条路线所用的数量关系。
\end{example}

非标准的根关系也可以先改变研究对象。例如，要比较\(u+v\)与一个依赖\(u\)的表达式，先把结论整理为\(v>\varphi(u)\)，再研究\(f(\varphi(u))-f(u)\)。构造后必须确认\(v,\varphi(u)\)位于可比较的同一单调区间。目标式决定比较点，原方程则负责把另一个根的函数值消去。

根和与根积对应的比较点也不同。要证明\(u+v>2c\)，可比较\(v\)与\(2c-u\)；要证明\(uv<K\)且\(u>0\)，则可比较\(v\)与\(K/u\)。把目标改写为两点的位置关系，便为辅助函数的选择提供了依据。

\begin{exercise}
\mbox{}\par\nobreak

（1）若\(u<v\)是\(x-\ln x=a\)的两个根，证明\(a>1\)、\(0<u<1<v\)，并由\(u+v=2a+\ln(uv)\)证明\(u+v<2a\)。

（2）在上一题条件下，证明\(u+v>a+1\)。可将目标改写为\(v>e^{1-u}\)，比较\(f(e^{1-u})\)与\(f(u)\)，其中\(f(x)=x-\ln x\)。

（3）对\(0<a<b\)，用对数平均不等式估计\(\ln b-\ln a\)。比较所得两侧界与中值定理给出的\((b-a)/b\)、\((b-a)/a\)。

（4）若\(f\)在\(c\)两侧分别严格递增、严格递减，且\(f(c+t)>f(c-t)\)，研究同一高度的两点\(u<c<v\)的中点位于\(c\)的哪一侧。写出所用比较点处于定义域的条件。

（5）设\(k>0\)，\(u<v\)满足\(u-k\ln u=v-k\ln v\)。证明\(u+v>2k\)、\(uv<k^2\)，并说明这两个界为什么随\(k\)分别作一次、二次伸缩。

（6）对\(0<t<1\)，证明\(2t<\ln\frac{1+t}{1-t}<\frac{2t}{\sqrt{1-t^2}}\)，并说明它与反射差\(R(t)\)的关系。

（7）（2021年新高考I卷第22题第（2）问）设\(a,b\)是两个不相等的正数，且\(b\ln a-a\ln b=a-b\)。证明\(2<1/a+1/b<e\)。
\end{exercise}

\endgroup
\FloatBarrier
\newpage
\subsection{零点随参数的变化*}

固定参数时，零点是一个数；让参数变化，同一单调支上的零点也可以形成一个函数。先确定这个函数存在，再研究它的变化，能够把一组方程的根联系起来。

若方程可以写成\(a=\varphi(x)\)，且\(\varphi\)在所选区间上连续严格单调，就可以用反函数\(x=x(a)\)描述这一支的根。在\(\varphi'(x)\ne0\)处，由反函数求导公式，
\begin{equation*}
 x'(a)=\frac1{\varphi'(x(a))}.
\end{equation*}
这既说明根怎样随参数移动，也给出移动速度。若研究两个根的和、积，分别确定两支后，再对组合求导。

以\(x-ae^x=0\)为例，写成\(a=xe^{-x}\)。在\(0<a<1/e\)时，两支分别给出\(0<u(a)<1<v(a)\)。对\(x-ae^x=0\)沿每一支求导，利用\(ae^x=x\)，得到
\begin{equation*}
 x'(a)=\frac{e^x}{1-ae^x}=\frac{x}{a(1-x)}.
\end{equation*}
因此左根随\(a\)增大而增大，右根则减小。对于它们的和，有
\begin{equation*}
 (u+v)'=\frac{u+v-2uv}{a(1-u)(1-v)}.
\end{equation*}
分母为负。两根还满足\(u-\ln u=v-\ln v\)，故由前面的平均值关系，\(u+v>2\)、\(uv<1\)，分子为正。于是两根之和随\(a\)增大而减小。

单独观察两根一增一减，并不能决定它们的和怎样变化。零点方程、分支导数和根之间的大小关系在这里分别提供了不同的信息。

\begin{example}
对上述两支根\(u(a),v(a)\)，研究当\(a\to0^+\)及\(a\to(1/e)^-\)时它们各自的极限。再研究\(v(a)/u(a)\)的单调性，比较直接求商的导数与先求对数导数的方法。
\end{example}

\begin{exercise}
\mbox{}\par\nobreak

（1）设\(u(a)<v(a)\)为\(x-\ln x=a\)的两个根。求两支的导数，判断根的间距\(v-u\)怎样随\(a\)变化。

（2）证明上题中的\(u+v\)随\(a>1\)增大而增大，并研究\(a\to1^+\)时的极限。与正文中的参数\(a\)比较：两个问题中的参数是否表示同一个量？

（3）对\(a=e^c+2c\)确定的\(c=c(a)\)，证明\(c'(a)=1/(e^c+2)\)。再对最小值\(m(a)=e^{c(a)}+c(a)^2-ac(a)\)求导，观察驻点方程怎样使含\(c'\)的项消去。
\end{exercise}

\subsection{比较函数的构造*}

前面的切线和Taylor多项式提供了简单的局部模型。若希望比较式在较大范围内仍容易处理，也可以选择多项式之比。选择时先匹配局部变化，再单独证明所需区间内的误差符号。

以\(\ln(1+h)\)在零点附近的二阶展开为例，尝试
\begin{equation*}
 Q(h)=\frac{Ah}{1+Bh}=Ah-ABh^2+o(h^2).
\end{equation*}
匹配\(h-h^2/2\)，得到\(A=1\)、\(B=1/2\)，即\(Q(h)=2h/(2+h)\)。这种通过匹配Taylor系数构造的有理逼近，称为Padé逼近。这里仅研究这个低阶例子。

局部匹配并不自行保证全区间的大小关系。令\(E(h)=\ln(1+h)-2h/(2+h)\)，则
\begin{equation*}
 E(0)=0,\qquad E'(h)=\frac{h^2}{(1+h)(2+h)^2}.
\end{equation*}
由此知道，在\(h>0\)时\(E(h)>0\)，在\(-1<h<0\)时\(E(h)<0\)。换成\(x=1+h\)，得到
\begin{equation*}
 \ln x\begin{cases}
 >2(x-1)/(x+1),&x>1,\\
 <2(x-1)/(x+1),&0<x<1.
 \end{cases}
\end{equation*}
现在这个比较已由区间上的导数符号证明，可以用于该区间中的零点定位。

\paragraph{由比较曲线约束零点} \mbox{}

设\(u<v\)是\(f(x)=x-\ln x=a\)、\(a>1\)的两个根。取比较函数
\begin{equation*}
 g(x)=x-\frac{2(x-1)}{x+1}
     =\frac{x^2-x+2}{x+1}.
\end{equation*}
在\((0,1)\)上\(g<f\)，在\((1,+\infty)\)上\(g>f\)，且\(g\)在这两段分别严格递减、严格递增。方程\(g(x)=a\)转为
\begin{equation*}
 x^2-(a+1)x+2-a=0.
\end{equation*}
它有两个根\(r_-<1<r_+\)。右支比较给出\(v>r_+\)；若\(r_->0\)，左支比较给出\(u>r_-\)，若\(r_-\le0\)，这个不等式也直接成立。因此
\begin{equation*}
 u+v>r_-+r_+=a+1.
\end{equation*}
超越方程不易显式求根，而比较方程是二次方程，韦达关系恰好给出所需的根和。这里先证明曲线的上下关系，再在各单调支上严格比较根的位置。

同一个结论也可以按目标直接构造\(f(\varphi(u))-f(u)\)。两条路线分别利用了目标表达式与简单比较曲线；哪一条更方便，要看能否找到易于判号的差，以及比较方程是否保留所需的和、积信息。

\begin{example}
设\(u<v\)为\(x-\ln x=a\)的两个根。取\(q(x)=(x+x^{-1})/2\)，证明它在两根所在的两段上分别位于\(f\)的上、下方，再由\(q(x)=a\)的根证明\(u+v<2a\)。与利用\(uv<1\)的证明比较。
\end{example}

\begin{exercise}
\mbox{}\par\nobreak

（1）设\(Q(x)=(1+Ax)/(1+Bx)\)。使它在零点匹配\(e^x\)直到二阶变化，求\(A,B\)。再在\((-2,2)\)内研究\(e^x\)与\(Q(x)\)的大小关系，可以先比较\(x\)与\(\ln Q(x)\)。

（2）在\(x=1\)附近比较\(\ln x\)与\(2(x-1)/(x+1)\)的Taylor多项式。两者从哪一阶开始不同？这能否直接代替正文中全区间的误差符号证明？

（3）设\(f,g\)在\(x>c\)上均严格递增，且\(g(x)>f(x)\)。若\(r,s>c\)且\(f(r)=g(s)\)，判断\(r,s\)的大小。再写出严格递减时的结论，用以核对正文两支比较的方向。

（4）在\(x=0\)处，用\((1+Ax)/(1+Bx)\)匹配\(\sqrt{1+x}\)直到二阶变化。研究所得有理函数与\(\sqrt{1+x}\)在\(x\ge0\)时的大小关系，并估计\(\sqrt{1.1}\)。

（5）（2018年全国I卷理科第21题第（2）问）设\(f(x)=1/x-x+a\ln x\)有两个极值点\(x_1,x_2\)。证明\(\bigl(f(x_1)-f(x_2)\bigr)/(x_1-x_2)<a-2\)。可先利用驻点方程，再选择对数的有理界。

（6）（2021年新高考I卷第22题再研究）回看习题5.28第（7）题。令\(u=1/a,v=1/b\)，不妨设\(u<v\)，再令\(t=v/u\)。将\(u+v\)表示为\(t\)的函数，用本小节的\(\ln t>2(t-1)/(t+1)\)研究其单调性与端点极限，重新证明原结论。
\end{exercise}

\FloatBarrier
\newpage
\section{定积分}
\label{sec:calculus-integrals}

导数研究一点附近的变化率。现在固定一段区间，研究各小段的贡献怎样合成一个总量。先把区间分成有限段，再使每段足够短，是建立这种累积运算的出发点。

\subsection{从有限求和到累积}

若质点在三个相邻时间段内的速度分别为\(v_1,v_2,v_3\)，持续时间分别为\(\Delta t_1,\Delta t_2,\Delta t_3\)，则总位移为
\begin{equation*}
 v_1\Delta t_1+v_2\Delta t_2+v_3\Delta t_3.
\end{equation*}
这里各项可以为负，正负表示位置变化的方向。若求路程，则应把各段速度换成速率\(|v_i|\)。

速度连续变化时，可以先在每个短时间段内选取一个时刻，用该时刻的速度近似整段速度，仍得到同样形式的和式。分段越短，局部近似越准确；还需要判断不同的分段与取样，是否会给出同一个极限。

曲边图形的面积也提出同一问题。以\(y=x\)在\([0,1]\)上与横轴围成的区域为例，将区间\(n\)等分，分别用每小段的左、右端点高度作矩形，得到
\begin{equation*}
 L_n=\frac1n\sum_{k=0}^{n-1}\frac kn=\frac{n-1}{2n},\qquad
 U_n=\frac1n\sum_{k=1}^{n}\frac kn=\frac{n+1}{2n}.
\end{equation*}
两种近似之差为\(1/n\)，都趋于\(1/2\)。区域的面积被夹在两者之间，因而也被确定。一般累积量未必已有几何公式，但有限和式及其误差仍然可以研究。

\begin{figure}[!ht]
\centering
\includegraphics[width=.83\linewidth]{calculus-riemann-sums}
\caption{在各小段内取样形成矩形和；分割加细后，矩形更贴近曲线。一般定义还要求结果不依赖分割与取样。}
\label{fig:calculus-riemann-sums}
\end{figure}

\begin{example}
某质点在\(0\le t<1\)时速度为\(2\)，在\(1\le t\le3\)时速度为\(-1\)。求总位移与总路程。再对速度\(v(t)=t\)、\(0\le t\le2\)，分别用四个等长时间段的左、右端点速度估计位移。
\end{example}

\subsection{定积分的定义}

\begin{definition}
设\(a<b\)。一组分点
\begin{equation*}
 P:\quad a=x_0<x_1<\cdots<x_n=b
\end{equation*}
称为\([a,b]\)的一个{\heiti{}分割}。记\(\Delta x_i=x_i-x_{i-1}\)，称\(\|P\|=\max_i\Delta x_i\)为分割的{\heiti{}最大子区间长度}。在每段中任选\(\xi_i\in[x_{i-1},x_i]\)，得到{\heiti{}Riemann和}
\begin{equation*}
 S(f;P,\xi)=\sum_{i=1}^n f(\xi_i)\Delta x_i.
\end{equation*}
\end{definition}

\begin{definition}
设\(f\)在\([a,b]\)上有界。若存在实数\(J\)，使对任意\(\varepsilon>0\)，均有\(\delta>0\)，只要\(\|P\|<\delta\)，无论怎样选取分割及取样点，都有
\begin{equation*}
 |S(f;P,\xi)-J|<\varepsilon,
\end{equation*}
就称\(f\)在\([a,b]\)上{\heiti{}Riemann可积}，称\(J\)为\(f\)在该区间上的{\heiti{}定积分}，记为
\begin{equation}\label{eq:calculus-riemann-definition}
 \int_a^b f(x)\,dx=\lim_{\|P\|\to0}\sum_{i=1}^n f(\xi_i)\Delta x_i.
\end{equation}
\end{definition}

最大子区间长度趋零，要求所有小段都变短，仅使分段数增加还不够。定义也要求全部足够细的取样和接近同一个数，仅算出一列等分和的极限，并未证明可积性。

式\eqref{eq:calculus-riemann-definition}中的\(x\)是积分变量，可以换成\(t\)或其他记号；固定函数和上下限后，定积分是一个数。

常值函数\(f(x)=c\)的任意取样和都是\(c(b-a)\)，所以积分等于这个数。对\(f(x)=x\)，注意到
\begin{equation*}
 \frac{b^2-a^2}{2}=\sum_{i=1}^n\frac{x_{i-1}+x_i}{2}\Delta x_i.
\end{equation*}
每个取样点与相应中点的距离不超过\(\Delta x_i/2\)，故
\begin{equation*}
 \left|S(f;P,\xi)-\frac{b^2-a^2}{2}\right|
 \le\frac12\sum_i(\Delta x_i)^2
 \le\frac{b-a}{2}\|P\|\longrightarrow0.
\end{equation*}
这证明了\(\int_a^b x\,dx=(b^2-a^2)/2\)，且结果不依赖分割和取样。

\begin{exercise}
\mbox{}\par\nobreak

（1）用定义求\(\int_a^b(kx+c)\,dx\)，并给出取样和误差关于\(\|P\|\)的界。

（2）设分割\(P_n\)始终保留\([0,1/2]\)这一整段，只将\([1/2,1]\)不断细分。分段数趋于无穷时，\(\|P_n\|\)是否趋零？对\(f(x)=x\)，改变第一段的取样能否产生不消失的差异？

（3）对\(f(x)=x^2\)，计算\([0,1]\)上等分分割的左右端点和及其极限。区分这个计算与证明一般定义所需完成的任务。
\end{exercise}

\subsection{可积函数类}

定义给出了可积性的统一要求，还需要认识哪些函数满足它。以下记\(\mathcal R[a,b]\)为\([a,b]\)上全体Riemann可积函数组成的集合；连续函数是其中首先研究的一类。

\begin{theorem}\label{thm:calculus-continuous-integrable}
闭区间上的连续函数Riemann可积。
\end{theorem}

连续函数在每一点附近变化很小。要控制整段的取样和，还需要让所有小段的波动同时变小：若每段内最大值与最小值之差都不超过\(\eta\)，则同一分割的任意两个取样和相差不超过
\begin{equation*}
 \sum_{i=1}^n\eta\Delta x_i=\eta(b-a).
\end{equation*}
闭区间连续性保证，当各段都足够短时，可以得到这样的统一控制。再比较不同分割，就能建立同一个和式极限。完整证明在下一节用第二章的一致连续性完成；以下主讲内容直接使用本定理。

定理保证了\(x^2\)、\(\sin x\)、\(e^{-x^2}\)等连续函数在任意有限闭区间上的积分存在，即使暂时没有算出其值，也可以研究其大小和性质。

\begin{theorem}\label{thm:calculus-other-integrable-classes}
在闭区间上，有界单调函数以及只有有限个间断点的有界函数都可积。
\end{theorem}

单调函数可以发生跳跃，但各段的上、下波动不能任意累积，因为它们受整个区间内总增量的控制。有限个间断点附近的函数值则未必接近，但可以将这些位置放进总长度很小的区间中，用有界性控制其贡献。这两种控制与连续函数的小段波动控制不同，证明统一放在下一节。

有限分段连续且有界的函数属于后一类。定理的适用范围还更广：只要在有限个分点之外连续并保持有界，分点处甚至可以没有单侧极限。例如
\begin{equation*}
 f(x)=\begin{cases}\sin(1/x),&0<x\le1,\\0,&x=0\end{cases}
\end{equation*}
在零点附近不断振荡，却只有这一个间断点，因而在\([0,1]\)上可积。可积性要求累积误差能够控制，并不要求每一点附近都没有振荡。

反过来，有界还不足以保证可积。考察\([0,1]\)上的函数
\begin{equation*}
 D(x)=\begin{cases}1,&x\in\mathbb Q,\\0,&x\notin\mathbb Q.\end{cases}
\end{equation*}
每个非零长度区间内都有有理数和无理数。因此同一分割上，全部选有理取样点的和为\(1\)，全部选无理取样点的和为\(0\)，无论分割多细都不能趋于同一个数。这个函数不Riemann可积。

\begin{example}
判断下列函数在\([0,1]\)上是否可积，并说明所用的函数类：\(|x-1/2|\)；在\(x<1/2\)时为\(0\)、其余位置为\(1\)的阶跃函数；在正半轴为\(\sin(1/x)\)、零点补零的函数。比较它们的连续性与可积性。
\end{example}

有限个间断点是一个充分条件，不是全部可积函数的共同限制。例如，令\(g(0)=0\)，并在\(1/(n+1)<x\le1/n\)时取\(g(x)=1/n\)。它在\([0,1]\)上有界且单调不减，在无穷多个分点处发生跳跃，仍由单调函数的可积性得到积分存在。

\begin{exercise}
\mbox{}\par\nobreak

（1）画出上述\(g\)在\([1/4,1]\)上的图像，证明它单调不减，并判断在零点是否连续。

（2）在\([0,1]\)上，令\(h(1/2)=1\)，其余位置\(h(x)=0\)。用取样和定义证明\(h\)可积且积分为零。

（3）分别说明“连续”“只有有限个间断点”“单调”“有界”与Riemann可积之间已经建立的蕴含关系，指出哪些条件不能反过来推得。
\end{exercise}

\subsection{积分的基本性质}

以下各式均在相应积分存在的范围内使用。由有限和式的加法、数乘及大小关系，取极限便得到积分的共同性质。

\begin{theorem}\label{thm:calculus-integral-properties}
若\(f,g\)在\([a,b]\)上可积，则对常数\(\alpha,\beta\)，有
\begin{equation*}
 \int_a^b(\alpha f+\beta g)=\alpha\int_a^b f+\beta\int_a^b g.
\end{equation*}
若\(f(x)\le g(x)\)，则\(\int_a^b f\le\int_a^b g\)。对\(a<c<b\)，积分还满足区间可加性：
\begin{equation}\label{eq:calculus-integral-additivity}
 \int_a^b f(x)\,dx=\int_a^c f(x)\,dx+\int_c^b f(x)\,dx.
\end{equation}
函数\(|f|\)也可积，且
\begin{equation}\label{eq:calculus-integral-absolute-bound}
 \left|\int_a^b f(x)\,dx\right|\le\int_a^b|f(x)|\,dx.
\end{equation}
\end{theorem}

\begin{proof}
线性性与保序性分别来自同一分割、同一取样下的和式运算。对于区间可加性，先将分割在\(c\)处分开，两部分的和分别趋于对应积分；若原分割没有\(c\)，插入它只影响包含\(c\)的一段，其误差受常数乘\(\|P\|\)控制，趋于零。此论证也说明，两侧分别可积的有界函数可以拼接为可积函数。\(|f|\)的可积性将在下一节由小段波动证明；再由\(-|f|\le f\le|f|\)得到绝对值估计。
\end{proof}

规定
\begin{equation*}
 \int_a^a f=0,\qquad \int_b^a f=-\int_a^b f.
\end{equation*}
这样式\eqref{eq:calculus-integral-additivity}对各端点的不同次序仍然成立。积分上下限的方向，与被积函数的正负共同决定累积的符号。

改变有限个点的函数值不改变积分：这些点只影响有限个小区间，所对应的和式差随最大子区间长度趋零。因此，有限分段连续的有界函数也可逐段积分，分点处的取值不影响结果。

若\(f\ge0\)，定积分给出图像与横轴之间的面积；若\(f\)有正有负，则积分是上方区域面积减去下方区域面积，几何总面积由\(\int_a^b|f|\)给出。

\begin{example}
不使用原函数，利用几何及区间可加性计算\(\int_{-1}^1x\,dx\)、\(\int_{-1}^1|x|\,dx\)和\(\int_0^2|x-1|\,dx\)。解释三个积分的符号与面积意义。
\end{example}

\subsection{积分估计}

若\(m\le f(x)\le M\)，保序性立即给出
\begin{equation}\label{eq:calculus-integral-bounds}
 m(b-a)\le\int_a^b f(x)\,dx\le M(b-a).
\end{equation}
特别地，若\(|f|\le K\)，则\(|\int_a^b f|\le K(b-a)\)。

函数近似的误差也可以这样累积。若\(|f(x)-g(x)|\le\eta\)，则
\begin{equation}\label{eq:calculus-integral-approximation-bound}
 \left|\int_a^b f-\int_a^b g\right|\le\eta(b-a).
\end{equation}
一个对全区间有效的函数值误差界，乘以区间长度，就得到总累积量的误差界。

若一个全局界过于粗糙，可以先拆分区间。设第\(i\)段上\(m_i\le f\le M_i\)，则
\begin{equation*}
 \sum_i m_i\Delta x_i\le\int_a^b f(x)\,dx\le\sum_i M_i\Delta x_i.
\end{equation*}
不同位置使用不同的界，可以保留函数在各段上的实际大小。积分的计算与估计都从同一个累积过程产生，估计并不以已经求得原函数为前提。

若连续函数\(f\ge0\)，且在某点取正值，则在一小段上仍有正的下界，故\(\int_a^b f>0\)。因此，连续函数\(f\le g\)只要不恒等，就有严格的积分大小关系。

\begin{example}
不求精确值，利用\(0\le x\le1\)时\(1-x^2\le e^{-x^2}\le1\)，估计\(\int_0^1 e^{-x^2}\,dx\)。再将区间分成两半，分别用每段的函数值上下界估计，比较所得范围。
\end{example}

\subsection{积分中值定理}

\begin{definition}
称\(\displaystyle\frac1{b-a}\int_a^b f(x)\,dx\)为\(f\)在\([a,b]\)上的{\heiti{}平均值}。
\end{definition}

\begin{theorem}[积分中值定理]\label{thm:calculus-integral-mean-value}
若\(f\)在\([a,b]\)上连续，则存在\(\xi\in[a,b]\)，使
\begin{equation*}
 \int_a^b f(x)\,dx=f(\xi)(b-a).
\end{equation*}
\end{theorem}

\begin{proof}
记\(f\)的最小值、最大值为\(m,M\)。由式\eqref{eq:calculus-integral-bounds}，平均值介于\(m,M\)之间，再由介值定理，它等于某点的函数值。
\end{proof}

平均值使一个变化的函数与一个常值函数具有相同总累积。它由整个区间决定，一般不等于端点值的平均。例如，\(x^2\)在\([0,1]\)上的平均值为\(1/3\)，两个端点值的平均却为\(1/2\)。这里\(\int_0^1x^2\,dx=1/3\)可由前面的等分和与可积性定理确定。

\begin{exercise}
\mbox{}\par\nobreak

（1）已知\(\int_0^1f=2\)、\(\int_1^3f=-1\)，求\(\int_0^3f\)、\(\int_3^0f\)及\(\int_0^1(3f-2)\)。

（2）设\(f\)连续且\(1\le f(x)\le3\)，\(x\in[0,2]\)。估计其积分和平均值。若积分等于\(2\)，能否确定\(f\)？

（3）设\(f\)在\([a,b]\)上连续。证明\(\int_a^b f^2=0\)当且仅当\(f\equiv0\)。

（4）若\(|f(x)-g(x)|\le10^{-3}\)，\(x\in[-2,3]\)，估计两积分之差。若只知道某一点的误差很小，能否得到同样结论？

（5）构造一个有界的分段常值函数，其平均值不等于任何一点的函数值，说明积分中值定理中连续性的作用。
\end{exercise}

\FloatBarrier
\newpage
\section{函数的可积性*}
\label{sec:calculus-integrability}

上一节列出了连续函数、有界单调函数和只有有限个间断点的有界函数。本节先建立共同的可积性判据，再分别研究这几类函数怎样控制取样误差。这样既能完成存在性证明，也能说明各类条件在证明中承担什么作用。

\subsection{上下和}

连续函数在每个闭子区间上取得最值，一般有界函数却未必如此。为比较全部取样值，可以使用最紧的上下界。

\begin{definition}
非空实数集合\(E\)的最小上界称为它的{\heiti{}上确界}，记为\(\sup E\)；最大下界称为它的{\heiti{}下确界}，记为\(\inf E\)。若最大值存在，它就等于上确界；上确界本身不必属于\(E\)。
\end{definition}

例如，\((0,1)\)没有最大值和最小值，但上确界为\(1\)、下确界为\(0\)。非空有界实数集合总有上下确界，这可以由第一章的单调有界原理得到：从一个非上界\(l_0\)和一个上界\(u_0\)出发，每次二分，若中点是上界就用它代替\(u_n\)，否则代替\(l_n\)。于是\(l_n\)递增、\(u_n\)递减，间距趋零，共同极限\(M\)仍是上界，而任何小于\(M\)的数最终小于某个非上界\(l_n\)，故不可能是上界。下确界同理。

设\(f\)在\([a,b]\)上有界。对分割\(P\)，定义
\begin{equation*}
 m_i=\inf_{[x_{i-1},x_i]}f,\qquad M_i=\sup_{[x_{i-1},x_i]}f,
 \qquad \omega_i=M_i-m_i.
\end{equation*}
称\(\omega_i\)为该段上的{\heiti{}振幅}。它描述全部函数值的波动范围，不要求上下界都能取到。定义{\heiti{}下和}与{\heiti{}上和}
\begin{equation*}
 L(f;P)=\sum_i m_i\Delta x_i,\qquad U(f;P)=\sum_i M_i\Delta x_i.
\end{equation*}
任意取样和满足\(L(f;P)\le S(f;P,\xi)\le U(f;P)\)，而上下和之差就是\(\sum_i\omega_i\Delta x_i\)。

\begin{figure}[!ht]
\centering
\includegraphics[width=.49\linewidth]{calculus-darboux}
\caption{蓝灰色矩形给出下和，叠加橙色部分得到上和；两者之差累积了各小段的振幅乘段长。}
\label{fig:calculus-darboux}
\end{figure}

若加点得到分割\(R\)，子段下界不低于原段下界，上界不高于原段上界，所以
\begin{equation*}
 L(f;P)\le L(f;R)\le U(f;R)\le U(f;P).
\end{equation*}
任意两个分割都可以把分点合在一起，得到共同加细分割。因此一个分割的下和不会超过另一个分割的上和。

\begin{example}
对\(f(x)=x^2\)，求\([0,1]\)上等分分割的上下和之差。再考察只在\(x=1/2\)处取\(1\)、其余位置为零的函数：哪些小段有非零振幅？说明这些小段的总贡献为何趋零。
\end{example}

\begingroup
\setlength{\parskip}{4pt}
\setlength{\abovedisplayskip}{10pt}
\setlength{\belowdisplayskip}{10pt}
\subsection{可积性判据}

\begin{theorem}[Darboux判据]\label{thm:calculus-darboux}
有界函数\(f\)在\([a,b]\)上Riemann可积，当且仅当对任意\(\varepsilon>0\)，都存在一个分割\(P\)，使
\begin{equation}\label{eq:calculus-darboux-criterion}
 U(f;P)-L(f;P)<\varepsilon.
\end{equation}
\end{theorem}

这个判据只要求找到上下和接近的分割，而积分定义要求任意足够细的分割及取样都接近同一个数。两者之间还需要一道比较。

\begin{proof}
先证必要性。若积分为\(J\)，足够细的分割上，所有取样和都与\(J\)相差小于\(\varepsilon/4\)。在每段选取接近上确界或下确界的函数值，可以让取样和分别任意接近上和或下和。因此\(U-L\le\varepsilon/2<\varepsilon\)。

再证充分性。设\(|f|\le B\)，先选定分割\(P_0\)，使上下和之差很小，记其内部分点个数为\(m\)。对另一个分割\(Q\)，加入\(P_0\)的分点得到\(R\)。除跨过这些固定分点的小段外，\(Q\)与\(R\)没有区别；受影响小段的总长度不超过\(m\|Q\|\)，每段振幅不超过\(2B\)，故
\begin{equation*}
 \begin{aligned}
 U(f;Q)-L(f;Q)&\le U(f;R)-L(f;R)+2Bm\|Q\|\\
 &\le U(f;P_0)-L(f;P_0)+2Bm\|Q\|.
 \end{aligned}
\end{equation*}
先使第一项小，再使\(\|Q\|\)足够小，就说明任意足够细的分割都有任意小的上下和之差。

对两个这样的细分割\(P,Q\)，在共同加细分割上任选取样和\(T\)。它与\(S(f;P,\xi)\)都夹在\(P\)的上下和之间，对\(Q\)也一样，所以
\begin{equation*}
 |S(f;P,\xi)-S(f;Q,\zeta)|
 \le U(f;P)-L(f;P)+U(f;Q)-L(f;Q).
\end{equation*}
右端可以任意小，且不依赖取样。

取一列最大子区间长度趋零的等分分割，在每个分割上任选和式\(S_n\)。上式表明\(\{S_n\}\)满足Cauchy准则，因而收敛于某个\(J\)。给定\(\varepsilon>0\)，先把细分程度选到任意两个和式相差小于\(\varepsilon/2\)，再取其中一个足够靠后的\(S_n\)，使\(|S_n-J|<\varepsilon/2\)。任意满足该细分要求的和式\(S\)均有
\begin{equation*}
 |S-J|\le|S-S_n|+|S_n-J|<\varepsilon.
\end{equation*}
这正是Riemann积分的定义。
\end{proof}

判据把一个极限存在性问题转为有限和的误差问题。证明中固定分点附近的误差控制，也说明为什么不能只研究一列特别选定的分割。
\endgroup

\subsection{连续函数的可积性}

现在证明定理\ref{thm:calculus-continuous-integrable}。闭区间上的连续函数一致连续。给定\(\eta>0\)，存在同一个\(\delta>0\)，使\(|u-v|<\delta\)时总有\(|f(u)-f(v)|<\eta\)。

选择\(\|P\|<\delta\)，每小段上振幅都小于\(\eta\)，于是
\begin{equation}\label{eq:calculus-upper-lower-gap}
 U(f;P)-L(f;P)=\sum_i\omega_i\Delta x_i<\eta(b-a).
\end{equation}
令\(\eta\)足够小，便满足Darboux判据，可积性得证。这里用到的不只是每个点附近函数值接近，而是所有小段同时有同一个波动上界。

若还知道\(|f'|\le K\)，则中值定理给出\(\omega_i\le K\Delta x_i\)，从而
\begin{equation*}
 U(f;P)-L(f;P)\le K\sum_i(\Delta x_i)^2\le K(b-a)\|P\|.
\end{equation*}
这提供了可计算的细分要求，但连续函数未必可导，一般证明仍依赖一致连续性。

上述估计还控制同一分割上的取样差异。任意两组取样和都在上下和之间，故其差的绝对值不超过\(U-L\)。当\(|f'|\le K\)时，对\([a,b]\)作\(n\)等分，只须\(n>K(b-a)^2/\varepsilon\)，就能使全部取样和之间的差小于\(\varepsilon\)。

\begin{exercise}
\mbox{}\par\nobreak

（1）设\(|f(u)-f(v)|\le C\sqrt{|u-v|}\)。给出上下和之差关于\(\|P\|\)的界，并证明可积性。

（2）用上题的方法处理\(f(x)=\sqrt{x}\)、\(0\le x\le1\)。为什么不能直接在整个区间上使用一个有限的导数上界？

（3）对\(f(x)=|x-1/2|\)、\(0\le x\le1\)，估计\(n\)等分时的上下和之差，并求使任意两组取样和相差小于\(10^{-3}\)的一个充分条件。
\end{exercise}

\subsection{单调函数的可积性}

设\(f\)在\([a,b]\)上单调不减。各段的下、上确界分别为\(f(x_{i-1})\)、\(f(x_i)\)，所以
\begin{equation*}
 \begin{aligned}
 U(f;P)-L(f;P)
 &=\sum_i(f(x_i)-f(x_{i-1}))\Delta x_i\\
 &\le\|P\|\sum_i(f(x_i)-f(x_{i-1}))
 =\|P\|(f(b)-f(a)).
 \end{aligned}
\end{equation*}
分割足够细时，该差趋零，故\(f\)可积。单调不增的情形应用于\(-f\)即可。这证明了定理\ref{thm:calculus-other-integrable-classes}中的第一项。

连续函数的证明使每段振幅都小；单调函数的证明则利用各段增量相加后恰为总增量。某些包含跳点的小段振幅并不随段长趋零，但它们对和式的影响仍可控制。

\begin{example}
对上一节的阶跃函数和无穷阶梯函数\(g\)，分别解释上述估计如何证明可积性。是否需要先数清全部间断点？
\end{example}

\subsection{有限个间断点的情形}

设\(f\)有界，\(|f|\le B\)，只有有限个间断点。用有限个小区间把这些点包住，使这些区间的总长度小于\(\ell\)。在剩下的有限个闭区间上，\(f\)连续，可以各自取足够细的分割，使振幅小于\(\eta\)。把所有分点合并得到\(P\)。

对包含间断点的小区间，只使用振幅不超过\(2B\)；对其余小段，使用连续性得到的振幅界。于是
\begin{equation*}
 U(f;P)-L(f;P)\le2B\ell+\eta(b-a).
\end{equation*}
先使\(\ell\)足够小，再使\(\eta\)足够小，就满足可积性判据。这证明了定理\ref{thm:calculus-other-integrable-classes}的第二项，也包括有限分段连续的有界函数。

这一证明允许间断点附近的振幅始终很大，只要求它们占据的区间总长度很小。例如，\(\sin(1/x)\)在零点附近无法用一个常值精确近似，但\([0,\ell]\)上振幅不超过\(2\)，其上下和之差的贡献不超过\(2\ell\)。

有限点的函数值也不会改变积分。若\(q\)只在\(N\)个点可能非零，且\(|q|\le B\)，任意取样和的非零贡献只能来自接触这些点的小段，其总长度不超过\(2N\|P\|\)，故
\begin{equation*}
 |S(q;P,\xi)|\le2NB\|P\|\longrightarrow0.
\end{equation*}
所以\(q\)可积且积分为零。由线性性，任意可积函数加上这样的\(q\)，积分保持不变。

\begin{exercise}
\mbox{}\par\nobreak

（1）对\(f(x)=\sin(1/x)\)、\(x>0\)，零点补零，具体说明怎样先选\([0,\ell]\)，再在\([\ell,1]\)上细分，使上下和之差小于给定误差。

（2）某函数除有限个点外等于\(x^2\)。证明它在\([0,1]\)上可积并求积分。若改为所有有理点处任意改变取值，原结论是否还必然成立？
\end{exercise}

\subsection{可积函数的运算}

可积函数的线性组合已由和式证明可积。Darboux判据还可以处理乘积与绝对值，从而补足主线使用的运算性质。

先补充区间上的限制。若\(f\)在\([a,b]\)上可积，向一个上下和接近的分割中加入\(c\)，再截取\([a,c]\)、\([c,b]\)两部分，每部分的上下和之差都不超过整个分割的差。因此\(f\)在这两个子区间上也可积，这给出了区间可加性中各项存在的依据。

设\(|f|\le A\)、\(|g|\le B\)，在任一小段内，对两点\(u,v\)，有
\begin{equation*}
 |f(u)g(u)-f(v)g(v)|
 \le A|g(u)-g(v)|+B|f(u)-f(v)|.
\end{equation*}
故乘积的振幅满足\(\omega_{fg}\le A\omega_g+B\omega_f\)。对\(f,g\)分别选择上下和接近的分割，再取共同加细分割，便得到
\begin{equation*}
 U(fg;P)-L(fg;P)
 \le A\bigl(U(g;P)-L(g;P)\bigr)+B\bigl(U(f;P)-L(f;P)\bigr).
\end{equation*}
右端可以任意小，所以\(fg\)可积。

又由\(\bigl||f(u)|-|f(v)|\bigr|\le|f(u)-f(v)|\)，知\(\omega_{|f|}\le\omega_f\)，故\(|f|\)可积。于是\(f^2\)、\(\max\{f,g\}\)、\(\min\{f,g\}\)也可积，后两者可写成\((f+g\pm|f-g|)/2\)。

这些结论说明\(\mathcal R[a,b]\)在常用运算下具有稳定的结构。连续函数还有额外性质：对连续\(f\)，\(\int f^2=0\)能推出\(f\equiv0\)；若只要求可积，单点取值的改变就使这一推断失效。

\begin{exercise}
\mbox{}\par\nobreak

（1）若\(f,g\)可积，证明\(|f-g|\)、\((f+g)^2\)可积。说明积分\(\int fg\)与\((\int f)(\int g)\)一般没有相等关系。

（2）构造不可积函数\(h\)，使\(|h|\)和\(h^2\)都可积。可考虑在有理点、无理点取不同符号。

（3）对上一节的\(D\)，直接求任意分割的上和与下和，用判据再次说明它不可积。
\end{exercise}

\FloatBarrier
\newpage
\section{微积分基本定理}
\label{sec:calculus-fundamental-theorem}

导数由差商极限定义，定积分由分割求和定义。现在让累积的终点变化，研究所得函数的导数，从而建立这两种运算之间的联系。

\subsection{累积函数}

设\(f\)在\([a,b]\)上Riemann可积。固定起点\(a\)，定义
\begin{equation}\label{eq:calculus-accumulation}
 A(x)=\int_a^x f(t)\,dt,\qquad x\in[a,b].
\end{equation}
这里\(x\)决定累积到哪里，\(t\)是在积分内部使用的变量。改变\(x\)会改变区间，因而定积分的值随之形成一个新的函数。

显然\(A(a)=0\)。由区间可加性，
\begin{equation}\label{eq:calculus-accumulation-increment}
 A(x+h)-A(x)=\int_x^{x+h}f(t)\,dt.
\end{equation}
即累积函数的增量，恰好等于新增加的有向小段上的累积。若\(f\ge0\)，增大终点会使累积函数单调不减；若\(f\)变号，则新增部分也可能使已有总量减小。

\begin{figure}[!ht]
\centering
\includegraphics[width=.83\linewidth]{calculus-accumulation}
\caption{左图为\(f(t)=t-1\)，右图为\(A(x)=\int_0^x f(t)\,dt\)。橙色窄带给出\(A(x+h)-A(x)\)，其除以\(h\)后的极限为\(f(x)\)。}
\label{fig:calculus-accumulation}
\end{figure}

又设\(|f|\le M\)，则\(|A(y)-A(x)|\le M|y-x|\)，所以\(A\)连续。这一结论只用了积分估计，还没有使用导数。

即使\(f\)在少数位置跳跃，累积函数也不会因此发生跳跃。单个点没有宽度，真正改变累积量的是一段区间上的贡献。若两个可积函数只在有限点取值不同，它们从同一起点得到的累积函数完全相同。

\begin{example}
设\(f(t)=t-1\)，\(0\le t\le2\)。先不求\(A(x)=\int_0^x f(t)\,dt\)的表达式，根据面积与符号判断它的增减、零点及最小值位置，再用几何面积求出\(A(x)\)。
\end{example}

\subsection{累积函数的导数}

当\(h\)很小时，\([x,x+h]\)内的函数值接近\(f(x)\)，所以这段累积应接近\(f(x)h\)。下面把这个近似写成误差估计。

\begin{theorem}[微积分基本定理]\label{thm:calculus-ftc-derivative}
若\(f\)在\([a,b]\)上连续，则式\eqref{eq:calculus-accumulation}所定义的函数在内部可导，且
\begin{equation*}
 A'(x)=f(x).
\end{equation*}
端点处相应的单侧导数也满足此式。
\end{theorem}

\begin{proof}
由式\eqref{eq:calculus-accumulation-increment}，对\(h\ne0\)，有
\begin{equation*}
 \frac{A(x+h)-A(x)}h-f(x)
 =\frac1h\int_x^{x+h}(f(t)-f(x))\,dt.
\end{equation*}
给定\(\varepsilon>0\)，由\(f\)在\(x\)处连续，当\(|h|\)足够小时，介于\(x\)与\(x+h\)之间的所有\(t\)都满足\(|f(t)-f(x)|<\varepsilon\)。故上式绝对值不超过\(\varepsilon|h|/|h|=\varepsilon\)，无论\(h\)取正或负都成立。因此差商趋于\(f(x)\)。
\end{proof}

\begin{definition}
若区间上\(F'(x)=f(x)\)，则称\(F\)为\(f\)在该区间上的一个{\heiti{}原函数}。
\end{definition}

上述差商证明实际上只用了\(f\)在所考察点\(x\)处连续，其余位置只须保证积分存在。因此，对一般可积函数，累积函数在被积函数的每个连续点处仍满足\(A'(x)=f(x)\)。但若\(f\)在一点跳跃，结论未必延伸到该点。

例如，在\([-1,1]\)上令\(f(t)=0\)（\(t<0\)）、\(f(t)=1\)（\(t\ge0\)），则\(\int_0^xf(t)\,dt\)在\(x\le0\)时为零，在\(x>0\)时为\(x\)。这个累积函数始终连续，却在零点出现折角。可积性保证累积量存在，连续性进一步使局部平均贡献趋于该点的函数值。

基本定理说明，每个连续函数都有原函数，其中一个就是固定起点的累积函数。原函数的存在由此得到保证，并不要求它能够用现有的初等函数写出。

例如，\(\int_0^x e^{-t^2}\,dt\)定义了一个可导函数，其导数为\(e^{-x^2}\)。这个积分表达式本身就可以作为研究函数的依据。

\subsection{Newton--Leibniz公式}

\begin{theorem}\label{thm:calculus-newton-leibniz}
设\(f\)在\([a,b]\)上Riemann可积，\(F\)在该区间连续且在内部满足\(F'=f\)。则
\begin{equation}\label{eq:calculus-newton-leibniz}
 \int_a^b f(x)\,dx=F(b)-F(a).
\end{equation}
\end{theorem}

\begin{proof}
对任意分割\(P\)，在各小段上对\(F\)使用中值定理，得到\(\xi_i\in(x_{i-1},x_i)\)，使
\begin{equation*}
 F(x_i)-F(x_{i-1})=f(\xi_i)\Delta x_i.
\end{equation*}
相加后左端逐项相消，故\(F(b)-F(a)=\sum_i f(\xi_i)\Delta x_i\)。由\(f\)可积，当\(\|P\|\to0\)时，右端趋于\(\int_a^bf\)，便得公式。
\end{proof}

当\(f\)连续时，上一小节还保证原函数存在。此时取累积函数\(A\)，由\((F-A)'=0\)知\(F-A\)为常数，再用\(A(a)=0\)得到\(F(x)-A(x)=F(a)\)。一个证明逐段累积总变化，另一个比较两个具有相同导数的函数，二者解释同一联系。

对于一般可积函数，Newton--Leibniz公式要求另有原函数\(F\)满足条件，不能仅凭可积就假定存在这样的\(F\)。同样，若只知道\(F\)可导，也没有自动得到\(F'\)的Riemann可积性。本章始终在已明确的函数类中使用公式。

记\(\left.F(x)\right|_a^b=F(b)-F(a)\)。例如，由\((x^3/3)'=x^2\)，有
\begin{equation*}
 \int_0^1x^2\,dx=\left.\frac{x^3}{3}\right|_0^1=\frac13.
\end{equation*}
这与先前由和式得到的结果相同。现在无须对每个积分重新计算分割和式，寻找一个原函数便能完成计算。

定积分的定义给出累积量是什么，可积性保证它存在，式\eqref{eq:calculus-newton-leibniz}则给出一种计算方式。前面的证明使这三个问题联系起来。

\begin{example}
计算\(\int_1^e dx/x\)、\(\int_0^\pi\sin x\,dx\)及\(\int_{-1}^1(3x^2-1)\,dx\)。再用几何意义核对最后一个积分的值是否意味着图像与横轴之间没有面积。
\end{example}

\subsection{变限积分}

设\(f\)连续，\(u,v\)可导。将变限积分写为两个累积函数之差，链式法则给出
\begin{equation}\label{eq:calculus-variable-limits}
 \frac d{dx}\int_{u(x)}^{v(x)}f(t)\,dt
 =f(v(x))v'(x)-f(u(x))u'(x).
\end{equation}
右端分别记录上下限移动的贡献；交换上下限时，公式仍按有向积分的约定成立。

例如，\(H(x)=\int_x^{x^2}e^{-t^2}\,dt\)的导数为\(2xe^{-x^4}-e^{-x^2}\)。这里的\(x\)通过上下限影响积分，\(t\)始终是积分变量；被积表达式内部另含自由参数时，不能直接套用式\eqref{eq:calculus-variable-limits}。

\begin{example}
设\(G(x)=\int_0^x e^{-t^2}\,dt\)，证明其为奇函数且严格递增，并求\(H(x)=G(\sin x)\)的导数和单调区间。
\end{example}

\begin{exercise}
\mbox{}\par\nobreak

（1）求\(\int_0^{x^2}\cos t\,dt\)、\(\int_x^{2x}\ln(1+t^2)\,dt\)的导数，并说明每一步的取值位置。

（2）设\(A(x)=\int_0^x(t^2-1)\,dt\)。先由被积函数判断\(A\)的增减与极值，再求其表达式。

（3）证明\(\displaystyle\lim_{h\to0}\frac1h\int_a^{a+h}f(t)\,dt=f(a)\)，其中\(f\)在\(a\)附近连续。比较它与积分中值定理的关系。

\begin{samepage}
（4）设连续函数\(f\)满足\(\int_0^x f(t)\,dt=x^2+\sin x\)，求\(f\)。若只知道\(\int_0^1 f=1\)，能否作同样判断？
\par
\end{samepage}

（5）设\(f\)在\([0,1]\)上连续，且对所有\(x\in[0,1]\)均有\(\int_0^x f(t)\,dt=0\)。证明\(f\equiv0\)。

（6）设\(f\)在含\(0\)的区间上连续、严格递增，\(A(x)=\int_0^x f(t)\,dt\)。证明\(A\)为凸函数，并解释图像的弯曲如何由局部累积的变化决定。

（7）设\(f\)为上文的阶跃函数。证明它在包含\(0\)的开区间上没有原函数。可以先在零点两侧使用“导数相同的函数相差常数”，再考察拼接处的可导性。

（8）设\(f\)在\(\mathbb R\)上连续，\(B(x)=\int_x^{x+1}f(t)\,dt\)。求\(B'\)，并证明\(B\)为常值函数当且仅当\(f(x+1)=f(x)\)对所有\(x\)成立。

（9）不求原函数，利用\(\sin u/u\to1\)与积分估计，求\(\displaystyle\lim_{x\to0}\frac1{x^3}\int_0^x\sin(t^2)\,dt\)。为什么仅把被积函数在零点的值代入，还不足以确定这个极限？

（10）设连续函数\(f\)满足\(f(x)=1+\int_0^xf(t)\,dt\)。先由等式说明\(f\)可导，再确定\(f\)，并代回原等式验证。这个关系是否唯一确定了\(f\)？

（11）设\(f\)在\([0,1]\)上连续且严格为正。证明存在唯一的\(c\in(0,1)\)，使\(\int_0^cf=\tfrac12\int_0^1f\)。若\(f\)还严格递增，判断\(c\)与\(1/2\)的大小关系，并用累积函数的变化解释。

（12）设\(H(x)=\int_{\cos x}^{\sin x}dt/(1+t^2)\)，\(0\le x\le\pi/2\)。不求\(H\)的显式表达式，判断它的单调性与零点，再比较\(H(\pi/6)\)与\(H(\pi/3)\)的符号。
\end{exercise}

\FloatBarrier
\newpage
\section{积分的计算}
\label{sec:calculus-integration-methods}

\begingroup
\setlength{\parskip}{3pt}
\setlength{\abovedisplayskip}{10pt}
\setlength{\belowdisplayskip}{10pt}

基本定理把定积分的计算与原函数联系起来。本节由求导结构组织寻找原函数的方法，再回到分割求和，研究没有现成原函数表达式时的数值近似。

\subsection{不定积分}

\begin{definition}
函数\(f\)在一个区间上的全部原函数，称为它在该区间上的{\heiti{}不定积分}，记为\(\int f(x)\,dx\)。
\end{definition}

若\(F'=f\)，则\(F+C\)也是原函数。反过来，任意两个原函数之差的导数为零，由中值定理，它们只相差常数。因此
\begin{equation*}
 \int f(x)\,dx=F(x)+C,
\end{equation*}
其中\(C\)为任意常数。定积分是确定的数，不定积分表示一个函数族；写出其中一个原函数，就能通过加常数得到全部原函数。

将基本求导公式反过来使用，得到
\begin{equation*}
 \begin{aligned}
 \int x^\alpha\,dx&=\frac{x^{\alpha+1}}{\alpha+1}+C\quad(\alpha\ne-1),&
 \int\frac{dx}{x}&=\ln|x|+C,\\
 \int e^x\,dx&=e^x+C,&\int a^x\,dx&=\frac{a^x}{\ln a}+C,\\
 \int\sin x\,dx&=-\cos x+C,&\int\cos x\,dx&=\sin x+C,\\
 \int\frac{dx}{1+x^2}&=\arctan x+C,&
 \int\frac{dx}{\sqrt{1-x^2}}&=\arcsin x+C,\\
 \int\frac{dx}{\cos^2x}&=\tan x+C,&
 \int\frac{dx}{\sin^2x}&=-\cot x+C,\\
 \int\sinh x\,dx&=\cosh x+C,&\int\cosh x\,dx&=\sinh x+C.
 \end{aligned}
\end{equation*}
各式沿用相应函数的定义范围，\(a>0\)、\(a\ne1\)。涉及\(1/x\)时，在不跨过零点的区间上使用；两个不同区间上的积分常数可以独立选择。

由求导的线性性质，不定积分也可以按线性组合寻找原函数。核对结果最直接的方法，是对所得表达式求导。

\begin{example}
求\(\int(3x^2-2\cos x+1/x)\,dx\)。再求满足\(F'(x)=2x+e^x\)、\(F(0)=3\)的函数\(F\)，说明初值怎样确定常数。
若另要求\(F(1)=5\)，这些条件能否同时满足？
再求\(\int(x+1)^2\,dx\)与\(\int(1+x^2)^{-1}\,dx\)，比较先作代数展开与直接识别基本公式的选择。
\end{example}

\subsection{换元积分}

若\(F'=f\)，链式法则给出
\begin{equation*}
 \frac d{dx}F(u(x))=f(u(x))u'(x).
\end{equation*}
因此，将内层函数的导数与外层表达式一同识别，就得到换元积分：
\begin{equation*}
 \int f(u(x))u'(x)\,dx=F(u(x))+C.
\end{equation*}
微分记号\(du=u'(x)\,dx\)将它写为\(\int f(u)\,du\)，依据是链式法则及微分形式不变性。

例如，在\(\int2x\cos(x^2)\,dx\)中令\(u=x^2\)，则\(du=2x\,dx\)，得到\(\sin(x^2)+C\)。若缺少这个\(2x\)，就不能把同一替换直接当作完整的换元。

\begin{theorem}[定积分换元公式]\label{thm:calculus-substitution}
设\(f\)在包含\(\varphi([\alpha,\beta])\)的区间上连续，\(\varphi\)在\([\alpha,\beta]\)上连续可导，则
\begin{equation}\label{eq:calculus-substitution}
 \int_\alpha^\beta f(\varphi(t))\varphi'(t)\,dt
 =\int_{\varphi(\alpha)}^{\varphi(\beta)}f(x)\,dx.
\end{equation}
\end{theorem}

\begin{proof}
取\(f\)的原函数\(F\)，左端由链式法则及Newton--Leibniz公式等于\(F(\varphi(\beta))-F(\varphi(\alpha))\)，也就是右端。
\end{proof}

这个公式不要求\(\varphi\)单调。例如\(\varphi(t)=\sin t\)、\(0\le t\le\pi\)时，代换后的起点和终点相同，有向贡献在往返中相消。若把一个给定积分改用新参数计算，则应选好对应的参数范围，并同时变换微分和上下限。

\begin{example}
计算\(\int_0^1 2x e^{x^2}\,dx\)、\(\int_0^1 dx/(1+x)^2\)及\(\int_{-\pi/2}^{\pi/2}\cos^2t\,dt\)。再用\(x=R\sin t\)计算\(\int_{-R}^R\sqrt{R^2-x^2}\,dx\)，与半圆面积比较。
\end{example}

\begin{exercise}
\mbox{}\par\nobreak

（1）求\(\int x/(1+x^2)\,dx\)、\(\int e^{2x+1}\,dx\)和\(\int\cos x/(1+\sin^2x)\,dx\)。

（2）计算\(\int_0^\pi f(\sin t)\cos t\,dt\)，其中\(f\)在\([0,1]\)上连续。解释结果与参数往返的关系。

（3）在\(\int_1^e(\ln x)^2/x\,dx\)中令\(u=\ln x\)，写出新积分的上下限并计算。

（4）求\(\int(\sqrt{x}+1/\sqrt{x})\,dx\)，其中\(x>0\)。再求满足\(F'(x)=3x^2-2x+1\)、\(F(1)=0\)的原函数，说明两问中积分常数的作用。

（5）计算\(\int x\sqrt{1+x^2}\,dx\)与\(\int_0^1x/(1+x^2)^2\,dx\)，指出所选中间变量及对应微分。

（6）计算\(\int_1^e dx/[x(1+\ln^2x)]\)与\(\int_0^{\pi/4}\tan x\,dx\)，说明所选的中间变量。
\end{exercise}

\subsection{分部积分}

乘积法则\((uv)'=u'v+uv'\)给出
\begin{equation}\label{eq:calculus-integration-by-parts}
 \int_a^b u(x)v'(x)\,dx
 =\left.u(x)v(x)\right|_a^b-\int_a^b u'(x)v(x)\,dx,
\end{equation}
其中\(u,v\)连续可导。不定积分的相应形式为\(\int u\,dv=uv-\int v\,du\)。

选择分部的目的，是让新的积分更容易处理。例如，对\(\int x e^x\,dx\)，将\(x\)求导、将\(e^x\)积分，得到
\begin{equation*}
 \int x e^x\,dx=xe^x-\int e^x\,dx=(x-1)e^x+C.
\end{equation*}
多项式因子求导后次数降低，因而这一选择可以继续使用。对\(\ln x\)，则可以把它看成\(\ln x\cdot1\)，通过积分\(1\)引入\(x\)。

对\(m\)次多项式\(P\)，连续分部直到多项式因子成为常数，得到
\begin{equation*}
 \int P(x)e^x\,dx
 =e^x\bigl(P(x)-P'(x)+\cdots+(-1)^mP^{(m)}(x)\bigr)+C.
\end{equation*}
对右端求导时，括号内的每一项都分别产生自身与下一阶导数，相邻项恰好抵消，只留下\(P(x)e^x\)。多项式求导最终为零，使分部过程有限终止；这种次数的降低，比只看表达式含有哪些因子更能说明选择的依据。

分部后也可能重新出现原积分。这时可以先计算一个定积分并记为\(J\)，将重复出现的\(J\)移到同一边，通过一个关于\(J\)的方程求值。对不定积分也可作相应整理，最后用求导核对，并保留任意常数。

\begin{example}
求\(\int\ln x\,dx\)、\(\int x\sin x\,dx\)及\(\int_0^1 x^2e^x\,dx\)。再计算\(\int e^x\cos x\,dx\)，观察两次分部后原积分如何重新出现。
\end{example}

\begin{exercise}
\mbox{}\par\nobreak

（1）计算\(\int_0^1x\arctan x\,dx\)及\(\int_1^e x\ln x\,dx\)。

（2）记\(J_n=\int_0^1x^ne^x\,dx\)。建立\(J_n\)与\(J_{n-1}\)的关系，并求\(J_3\)。

（3）设\(F\)在\([0,1]\)上连续可导，记\(f=F'\)。证明\(\int_0^1F(x)\,dx=F(1)-\int_0^1xf(x)\,dx\)。

（4）计算\(\int_0^1x\ln(1+x)\,dx\)。分部后得到的有理式可以先做怎样的代数整理？

（5）记\(I_n=\int_0^{\pi/2}\sin^nx\,dx\)。对\(n\ge2\)，用分部积分建立\(I_n\)与\(I_{n-2}\)的关系，并求\(I_4\)。

（6）分别计算\(\int xe^{x^2}\,dx\)与\(\int xe^x\,dx\)，比较换元与分部的选择。两个表达式都含有因子\(x\)，为什么不宜据此选择同一种方法？
\end{exercise}

\newpage
\subsection{定积分的计算}

定积分有时可以在寻找原函数之前，先由区间与函数的关系简化。作\(x=-t\)的换元，得到
\begin{equation*}
 \int_{-a}^a f(x)\,dx=\int_0^a(f(x)+f(-x))\,dx.
\end{equation*}
所以奇函数在对称区间上的积分为零，偶函数的积分为正半段积分的两倍。

同样，在\([a,b]\)上作\(x=a+b-t\)，有
\begin{equation}\label{eq:calculus-integral-reflection}
 \int_a^b f(x)\,dx=\int_a^b f(a+b-x)\,dx.
\end{equation}
若\(f(x)+f(a+b-x)=c\)，将两式相加就得到\(\int_a^b f=c(b-a)/2\)，无需求出原函数。

例如，\(f(x)=1/(1+e^x)\)满足\(f(x)+f(-x)=1\)，故\(\int_{-a}^a dx/(1+e^x)=a\)。这个计算使用了函数在整个区间上的配对关系。

若\(f(a+b-x)=f(x)\)，记\(J=\int_a^bxf(x)\,dx\)，反射换元给出\(J=\int_a^b(a+b-x)f(x)\,dx\)。相加即得
\begin{equation*}
 \int_a^bxf(x)\,dx=\frac{a+b}{2}\int_a^bf(x)\,dx.
\end{equation*}
对称关系使一次因子\(x\)可以由区间中点代替。这里同样没有要求先求出\(xf(x)\)的原函数。

绝对值、分段表达式和几何面积则通常需要先拆分区间，确定每段上的表达式或符号，再使用基本定理。

\begin{example}
计算\(\int_{-2}^2(x^3+\cos x)\,dx\)、\(\int_0^2|x^2-1|\,dx\)及\(\int_0^{\pi/2}\sin x/(\sin x+\cos x)\,dx\)。分别指出计算中起主要作用的关系。
\end{example}

\begin{exercise}
\mbox{}\par\nobreak

（1）设\(f\)连续且周期为\(T>0\)。证明\(\int_c^{c+T}f\)与\(c\)无关。可以拆分区间，也可以把它看成变限积分求导。

（2）设\(f\)在\([0,1]\)上可积，且\(f(x)+f(1-x)=x^2+(1-x)^2\)，求\(\int_0^1f(x)\,dx\)。

（3）计算\(\int_{-1}^2|x(x-1)|\,dx\)，并区分结果与\(\int_{-1}^2x(x-1)\,dx\)。

（4）用\(x=\pi-t\)的换元计算\(\int_0^\pi x\sin x\,dx\)，再用分部积分核对。比较两种方法分别利用了什么结构。

（5）不先求原函数，计算\(\int_0^1\dfrac{\ln(1+x)}{\ln(1+x)+\ln(2-x)}\,dx\)。说明\(x\)与\(1-x\)处的函数值怎样配对。
\end{exercise}

\subsection{数值积分}

即使不能从基本公式中直接找到原函数，也可以回到Riemann和近似积分。对\([a,b]\)作\(n\)等分，令\(h=(b-a)/n\)、\(x_i=a+ih\)，常用的近似有
\begin{equation*}
 L_n=h\sum_{i=0}^{n-1}f(x_i),\qquad
 R_n=h\sum_{i=1}^nf(x_i),\qquad
 M_n=h\sum_{i=1}^nf\!\left(\frac{x_{i-1}+x_i}{2}\right).
\end{equation*}
它们依次为左、右端点与中点矩形公式。各段改用弦近似，得到{\heiti{}梯形公式}
\begin{equation*}
 T_n=h\left(\frac{f(a)+f(b)}2+\sum_{i=1}^{n-1}f(x_i)\right)
 =\frac{L_n+R_n}{2}.
\end{equation*}
连续性保证这些近似都趋于积分，但固定\(n\)时的误差还取决于函数的变化。

若\(f\)单调递增，则\(L_n\le\int_a^b f\le R_n\)，且
\begin{equation*}
 R_n-L_n=h(f(b)-f(a)).
\end{equation*}
单调递减时上下界交换。若\(f\)为凸函数，图像不高于弦，故积分不超过梯形和；中点切线不高于图像，且其斜率项在对称小段上积分为零，故
\begin{equation*}
 M_n\le\int_a^b f(x)\,dx\le T_n.
\end{equation*}
凹函数的方向相反。

\begin{figure}[!ht]
\centering
\includegraphics[width=.78\linewidth]{calculus-quadrature}
\caption{凸函数的中点与梯形估计。左图各段中点切线的积分等于对应矩形面积，右图用弦给出上界。}
\label{fig:calculus-quadrature}
\end{figure}

若函数的单调性或凹凸性在区间内改变，可以先在变化点拆分，再对各段分别建立上、下界，最后相加。所需信息由当前的估计方法决定：单调性提供矩形界，凹凸性提供中点与梯形界；仅有有限个取样值时，不能自动获得这些界的方向。

若\(A\le I\le B\)，以\((A+B)/2\)近似\(I\)，误差至多为\((B-A)/2\)。故两界之差不超过\(2\varepsilon\)时，该近似的误差不超过\(\varepsilon\)。

\FloatBarrier
\newpage

将单段的中点与梯形估计用于\(e^t\)在\([\ln a,\ln b]\)上的积分，就得到
\begin{equation*}
 \sqrt{ab}\le\frac1{\ln b-\ln a}\int_{\ln a}^{\ln b}e^t\,dt
 =\frac{b-a}{\ln b-\ln a}\le\frac{a+b}{2}.
\end{equation*}
因指数函数严格凸，两侧均取严格不等号。这重新解释了命题\ref{prop:calculus-logarithmic-mean}中的平均值关系：对数平均也是指数函数在相应区间上的积分平均值。前面用导数建立的比较，在累积量中获得了同一个来源。

\begin{example}
对\(J=\int_0^1e^{-x^2}\,dx\)，写出四等分的左右端点估计，给出一个包含\(J\)的区间。再对\(\int_1^2dx/x\)写出两等分的中点与梯形估计，利用凸性判断它们的方向。
\end{example}

\begin{exercise}
\mbox{}\par\nobreak

（1）对\(\int_0^1x^2\,dx\)，求\(M_n,T_n\)及其精确误差，比较误差随\(n\)的变化。

（2）只使用\(e^{-x^2}\)的单调性，要使左右矩形和之差小于\(10^{-2}\)，可以取怎样的分段数？

（3）一个连续函数的端点值相同，是否足以断定梯形公式精确？举例说明。

（4）对\(\int_0^1\ln(1+x)\,dx\)，写出两等分的中点与梯形近似，利用凹性判断上下界，再计算积分值核对。

（5）已知\(f\)在\([0,1]\)上连续、单调递减，且\(f(0)=2\)、\(f(1/2)=1\)、\(f(1)=1/2\)。仅由这些数据给出积分的上下界，再写出梯形近似。这个近似是否必为精确值？

（6）设\(f(x)=x^3\)，\(-1\le x\le2\)。在零点拆分区间，各用一个中点矩形和一个梯形，分别求整个积分的上下界。若直接在\([-1,2]\)上只用一个梯形，它与真实积分的大小能否仅由整段凸性判断？

（7）用两等分的中点与梯形公式估计\(\int_0^1\sqrt{1+x}\,dx\)。取两界平均值作近似，给出无需计算精确积分的误差界。

（8）用梯形公式近似\(\ln2=\int_0^1dx/(1+x)\)。只利用单调性给出误差界，并确定使误差不超过\(10^{-3}\)的一个充分分段数。若只取左端点和，所需的分段数有何变化？

（9）设\(f\)在\([0,2]\)上连续且为凸函数，已知\(f(0)=1\)、\(f(1)=2\)、\(f(2)=5\)。用一个中点矩形和一个梯形给出积分的上下界，再改用两个梯形求上界。若还知\(f(1/2)=5/4\)、\(f(3/2)=13/4\)，写出两个中点矩形给出的下界，并估计两界平均值的误差。

\end{exercise}

\endgroup
\FloatBarrier
\newpage
\begingroup
\setlength{\parskip}{4pt}
\setlength{\abovedisplayskip}{10pt}
\setlength{\belowdisplayskip}{10pt}
\subsection{Taylor公式的积分余项*}

分部积分还可以把局部展开的余项写成累积量。先看一次近似：由基本定理，\(f(x)-f(a)=\int_a^xf'(t)\,dt\)。若进一步把\(f'\)的变化纳入积分，就能把线性项分离出来；重复这一过程，便得到任意有限阶的积分余项。

\begin{theorem}\label{thm:calculus-taylor-integral}
设\(f\)在包含\(a,x\)的区间上具有直到\(n+1\)阶的连续导数，则
\begin{equation}\label{eq:calculus-taylor-integral-remainder}
 f(x)=\sum_{k=0}^n\frac{f^{(k)}(a)}{k!}(x-a)^k
 +\frac1{n!}\int_a^x(x-t)^nf^{(n+1)}(t)\,dt.
\end{equation}
\end{theorem}

\begin{proof}
记右端积分项为\(Q_n\)。对\(n\ge1\)，分部积分得到
\begin{equation*}
 Q_n=-\frac{f^{(n)}(a)}{n!}(x-a)^n
 +\frac1{(n-1)!}\int_a^x(x-t)^{n-1}f^{(n)}(t)\,dt
 =Q_{n-1}-\frac{f^{(n)}(a)}{n!}(x-a)^n.
\end{equation*}
而\(Q_0=\int_a^xf'(t)\,dt=f(x)-f(a)\)。逐次代回，即得公式。
\end{proof}

记\(h=x-a\)，在余项中作\(t=a+sh\)的换元，可将不同长度的积分区间统一到\([0,1]\)：
\begin{equation}\label{eq:calculus-taylor-scaled-remainder}
 Q_n=\frac{h^{n+1}}{n!}\int_0^1(1-s)^n f^{(n+1)}(a+sh)\,ds.
\end{equation}
这个表示对正、负\(h\)都成立。因子\(h^{n+1}\)记录误差的阶数，其余部分是高阶导数沿基点与终点之间的加权累积。

若\(|f^{(n+1)}|\le M\)，由积分绝对值估计得到
\begin{equation*}
 |Q_n|\le\frac{M}{(n+1)!}|x-a|^{n+1}.
\end{equation*}
因为\(\int_0^1(1-s)^n\,ds=1/(n+1)\)，上界中的阶乘正是对这个权函数积分后得到的。若高阶导数在不同位置大小差别明显，保留积分表达式还可以得到比统一上界更细的估计。

\paragraph{与Lagrange余项的联系} \mbox{}

权函数\((1-s)^n\)非负。设\(f^{(n+1)}\)在所考察线段上的最小值、最大值为\(m,M\)，则
\begin{equation*}
 \frac m{n+1}\le\int_0^1(1-s)^nf^{(n+1)}(a+sh)\,ds\le\frac M{n+1}.
\end{equation*}
由连续性与介值定理，可在\(a,x\)之间取\(\xi\)，使中间积分等于\(f^{(n+1)}(\xi)/(n+1)\)。代回式\eqref{eq:calculus-taylor-scaled-remainder}，便得到
\begin{equation*}
 Q_n=\frac{f^{(n+1)}(\xi)}{(n+1)!}(x-a)^{n+1}.
\end{equation*}
Lagrange余项将这一段高阶变化的加权平均，由某一点的值实现；积分余项则保留各位置对误差的贡献。

\newpage
\paragraph{局部阶数与余项符号} \mbox{}

当\(n\ge1\)时，若\(f\)具有连续的\(n\)阶导数，对\(n-1\)阶积分展开再减去第\(n\)项，可得
\begin{equation*}
 f(a+h)-T_n(a+h)=\frac{h^n}{(n-1)!}
 \int_0^1(1-s)^{n-1}\bigl(f^{(n)}(a+sh)-f^{(n)}(a)\bigr)\,ds.
\end{equation*}
由\(f^{(n)}\)在\(a\)处连续，括号内的差在整段\(0\le s\le1\)上一致趋零，故右端为\(o(h^n)\)。这在当前连续高阶导数条件下再次得到Peano形式；5.5中的Peano定理还允许更弱的条件。

对一次近似，式\eqref{eq:calculus-taylor-scaled-remainder}给出特别直观的关系：
\begin{equation*}
 f(a+h)-f(a)-f'(a)h=h^2\int_0^1(1-s)f''(a+sh)\,ds.
\end{equation*}
若\(m\le f''\le M\)，则无论\(h\)的正负，都有
\begin{equation}\label{eq:calculus-quadratic-tangent-bounds}
 \frac m2h^2\le f(a+h)-f(a)-f'(a)h\le\frac M2h^2.
\end{equation}
二阶导数非负时，图像位于切线上方；二阶导数还有正下界时，图像与切线的距离至少含一个确定的二次贡献。这也接回了5.5的凸性与习题中的定量估计。

\begin{example}
设\(f(x)=\ln(1+x)\)，\(0\le x\le1\)。写出一次近似的积分余项，并由\(f''\)的上下界估计\(f(x)-x\)。再对\(e^x\)使用同一关系，比较这两个函数与零点切线的位置。
\end{example}

\begin{exercise}
\mbox{}\par\nobreak

（1）用积分余项证明\(x>0\)时\(e^x>1+x+x^2/2\)，与Lagrange余项的证明比较。

（2）对\(f(x)=\ln(1+x)\)，写出在零点的二阶积分余项，证明\(x\ge0\)时\(x-x^2/2\le\ln(1+x)\le x-x^2/2+x^3/3\)。

（3）设\(f''\ge m>0\)。对区间中点的左右增量分别使用式\eqref{eq:calculus-quadratic-tangent-bounds}，重新证明习题中\(\bigl(f(a)+f(b)\bigr)/2-f((a+b)/2)\)的下界。

（4）设\(f\)在\([0,1]\)上三阶连续可导。用积分余项比较\(f(1)\)与二阶Taylor多项式在\(1\)处的值。当\(f'''\)在前半段与后半段大小不同时，权函数如何影响两部分的贡献？

（5）设\(f'''\ge0\)且连续。用积分余项判断\(f(x)\)与\(T_2(x)\)在基点\(a\)左右两侧的大小关系，解释为什么两侧方向不同。

\begin{samepage}
（6）设\(f(0)=1\)、\(f'(0)=0\)，且在\([0,1]\)上\(f''\)连续并满足\(2\le f''\le4\)。先估计\(f(x)\)，再给出\(\int_0^1f(x)\,dx\)的上下界。
\par
\end{samepage}
\end{exercise}

\subsection{数值积分的误差估计*}

中点公式的线性误差可以在每个小段内相消，因此比单纯估计函数波动更精确。下面只研究这一种公式。

\begin{proposition}\label{prop:calculus-midpoint-error}
若\(f\)在\([a,b]\)上二阶连续可导，且\(|f''|\le M\)，则\(n\)等分的中点公式满足
\begin{equation*}
 \left|\int_a^b f(x)\,dx-M_n\right|\le\frac{M(b-a)^3}{24n^2}.
\end{equation*}
\end{proposition}

\begin{proof}
对长度为\(h\)、中点为\(m\)的小段，Taylor公式给出
\begin{equation*}
 |f(x)-f(m)-f'(m)(x-m)|\le\frac M2(x-m)^2.
\end{equation*}
线性项在\([m-h/2,m+h/2]\)上的积分为零。因此该小段的误差不超过
\begin{equation*}
 \frac M2\int_{m-h/2}^{m+h/2}(x-m)^2\,dx=\frac{Mh^3}{24}.
\end{equation*}
把\(n\)段误差相加，再代入\(h=(b-a)/n\)，便得结论。
\end{proof}

证明先用局部多项式近似，再把各段误差累积。真正进入最终上界的是整个区间上统一有效的\(M\)，而每一段线性项的相消解释了\(n^{-2}\)的误差阶。

若分割不等长，同一证明给出\(\frac M{24}\sum_i(\Delta x_i)^3\)的误差界；若各段使用各自的二阶导数上界\(M_i\)，还可以写为\(\frac1{24}\sum_iM_i(\Delta x_i)^3\)。因此函数弯曲较明显的位置可以分得更细，平缓部分则不必采用同样的段长。

又因为每段中点切线的线性项积分为零，当\(f''\ge0\)时，中点近似不超过真实积分。这将数值近似的方向、误差大小与前面凸性的几何判断联系起来。

\begin{exercise}
\mbox{}\par\nobreak

（1）对\(f(x)=e^{-x^2}\)、\(0\le x\le1\)，验证\(|f''(x)|\le2\)。要使中点公式误差小于\(10^{-4}\)，可以选取怎样的\(n\)？

（2）将上述证明用于\(f(x)=x^2\)，比较误差界与精确误差。为什么在这个例子中等号能够成立？

\begin{samepage}
（3）对\(\int_0^1x^2\,dx\)，分别使用分割\(0,1/2,3/4,1\)与三等分分割的中点近似，求精确误差，并与\(\frac M{24}\sum_i(\Delta x_i)^3\)核对。相同分段数下，哪一个分割更准确？
\par
\end{samepage}

（4）若\(f''\)连续且\(0<m\le f''\le M\)，给出任意分割的中点近似误差的正下界与上界。说明二阶导数的符号和大小分别提供什么信息。
\end{exercise}

\endgroup
\newpage
\subsection{特殊换元*}

二次表达式可以通过配方、因式分解或三角参数转为已知形式。方法的选择取决于当前表达式具有的结构。

对\(a>0\)，在\(\sqrt{a^2-x^2}\)中令\(x=a\sin t\)、\(-\pi/2\le t\le\pi/2\)，根式化为\(a\cos t\)。相应的\(dx=a\cos t\,dt\)仍须同时代入。类似地，\(x=a\tan t\)可以处理\(a^2+x^2\)，也可由双曲函数恒等式选择\(x=a\sinh t\)。

有理式则可以先分解。例如
\begin{equation*}
 \frac1{(x+1)(x+2)}=\frac1{x+1}-\frac1{x+2},
\end{equation*}
从而原积分化为两个对数。不可约二次因子先配方，则可以与反正切公式联系。这里的代数整理和几何参数选择，分别延续了第四章对二次表达式的研究。

\begin{example}
求\(\int dx/(x^2+2x+5)\)、\(\int dx/(x^2+3x+2)\)及\(\int\sqrt{a^2-x^2}\,dx\)。对最后一式，比较三角代换与圆的几何关系。
\end{example}

\begin{exercise}
\mbox{}\par\nobreak

（1）计算\(\int_0^{a/2} x^2/\sqrt{a^2-x^2}\,dx\)，其中\(a>0\)，并解释所选参数区间。

（2）令\(x=\sinh t\)，求\(\int\sqrt{1+x^2}\,dx\)，并对结果求导核对。

（3）计算\(\int_0^1dx/(x^2+4x+5)\)及\(\int_0^1dx/[(x+1)(x+3)]\)，比较配方与部分分式所利用的代数结构。

（4）对\(a>0\)，用\(x=a\sinh t\)求\(\int dx/\sqrt{x^2+a^2}\)，再将结果写成对数形式，并求导验证。

（5）计算\(\int_0^a(a^2-x^2)^{3/2}\,dx\)，其中\(a>0\)。利用三角代换化为三角函数幂的积分，并与本节分部积分中的递推关系联系起来。
\end{exercise}

\FloatBarrier
\newpage
\section{函数空间*}
\label{sec:calculus-function-spaces}

前面已经用函数近似另一个函数，也用误差界控制相应的积分。本节把整个函数作为研究对象：先确定函数之间的线性运算，再比较整体误差，最后研究以函数为输入的映射。以下固定\(a<b\)。

\subsection{连续函数空间}

记\(C[a,b]\)为\([a,b]\)上全体实值连续函数的集合。对\(f,g\in C[a,b]\)及实数\(\lambda\)，规定
\begin{equation*}
 (f+g)(x)=f(x)+g(x),\qquad(\lambda f)(x)=\lambda f(x).
\end{equation*}
和与数乘仍为连续函数，其余线性空间公理由实数运算逐点得到。因此\(C[a,b]\)是一个实线性空间。它的零向量是恒等于零的函数，向量相等则意味着在每一点取值都相同。

第三章的坐标向量有有限个分量。这里一个函数在整个区间上的取值共同确定一个元素，但加法、数乘和线性组合仍具有相同的代数意义。例如，次数不超过\(n\)的多项式构成一个子空间，由\(1,x,\ldots,x^n\)张成。

在\(f(x)+g(x)\)中，先选定一个位置再相加；在\(f+g\)中，则把所有位置的相加结果组成一个新函数。前者是数的运算，后者是空间中元素的运算。点值可以帮助计算，但一个点的取值不能确定整个元素。

前面引入的\(\mathcal R[a,b]\)也是实线性空间，且\(C[a,b]\subset\mathcal R[a,b]\)。阶跃函数说明这个包含关系是严格的。连续性并非线性运算成立的必要条件，但它会影响后面哪些误差量确实成为范数，以及积分误差量能否辨认不同函数。

\begin{example}
判断\(\{f\in C[a,b]:f(a)=0\}\)与\(\{f\in C[a,b]:\int_a^b f=1\}\)是否为线性子空间。若将后一条件的\(1\)改为\(0\)，结论如何变化？
\end{example}

\subsection{范数}

为了比较整个函数，需要把误差在全区间上的大小概括为一个数。对连续函数，最值定理使最大误差确实能够取到。

\begin{definition}
实线性空间\(V\)上的函数\(\|\cdot\|:V\to\mathbb R\)，若满足
\begin{equation*}
 \|v\|\ge0,\quad\|v\|=0\Longleftrightarrow v=0;\qquad
 \|\lambda v\|=|\lambda|\|v\|;\qquad
 \|u+v\|\le\|u\|+\|v\|,
\end{equation*}
就称为\(V\)上的一个{\heiti{}范数}。
\end{definition}

在\(C[a,b]\)上定义{\heiti{}最大值范数}
\begin{equation*}
 \|f\|_\infty=\max_{x\in[a,b]}|f(x)|.
\end{equation*}
非负性、零元素条件与数乘性质直接成立。又因为每一点都有\(|f(x)+g(x)|\le\|f\|_\infty+\|g\|_\infty\)，取最大值得到三角不等式。因此这确实是一个范数。

\(\|f-g\|_\infty\)衡量两函数在整个区间上的最大误差。例如，\(\|f-g\|_\infty\le\eta\)意味着图像在每个横坐标处的高度差都不超过\(\eta\)。这与只在某一点接近，或只在基点附近比较余项，是不同的要求。

例如，在\([0,1]\)上用\(p(x)=1+x\)近似\(e^x\)。两函数之差在零点为零，其导数为\(e^x-1\ge0\)，故
\begin{equation*}
 \|e^x-(1+x)\|_\infty=e-2.
\end{equation*}
若只在\([0,h]\)、\(0<h\le1\)上比较，同样得到最大误差\(e^h-1-h\)，并由Taylor公式知它不超过\(eh^2/2\)。区间缩小时的局部近似，与固定区间上的整体误差，由此可以放在同一例子中辨认。

不同范数也会记录不同信息。\(\int_a^b|f|\)在连续函数空间上是范数，因为非零的连续函数必在一小段上保持非零，绝对值积分为正。但若直接把同一表达式放在\(\mathcal R[a,b]\)上，单点取值为\(1\)、其余为零的函数就破坏了“范数为零只对应零函数”的要求。研究对象的选择参与决定哪些表达式能够充当距离的依据。

\begin{exercise}
\mbox{}\par\nobreak

（1）在\([0,1]\)上求\(\|x\|_\infty\)、\(\|x-x^2\|_\infty\)及\(\|x^2-1\|_\infty\)。

（2）在\(C[0,1]\)上，\(|f(0)|\)能否作为范数？证明\(\int_0^1|f(x)|\,dx\)是另一个范数，并说明它衡量误差的方式与最大值范数有何不同。
\end{exercise}

\subsection{逐点收敛}

\begin{definition}
设\(f_n\)是在同一区间上定义的一列函数。若对每个固定的\(x\)，数列\(\{f_n(x)\}\)都收敛于\(f(x)\)，就称\(f_n\)在该区间上{\heiti{}逐点收敛}于\(f\)。
\end{definition}

此时先固定位置，再让函数的指标\(n\)变化。对给定误差，不同位置可能需要不同的\(N\)。写出量词，就是
\begin{equation*}
 \forall x, \forall \varepsilon>0, \exists N=N(x,\varepsilon),
 \quad n>N\Longrightarrow |f_n(x)-f(x)|<\varepsilon.
\end{equation*}
函数指标是数列的指标，\(x\)则标记正在观察哪一列数。不同位置处的这些数列，一起构成函数列的逐点研究。

例如，\(f_n(x)=x^n\)在\([0,1]\)上逐点收敛于
\begin{equation*}
 f(x)=\begin{cases}0,&0\le x<1,\\1,&x=1.\end{cases}
\end{equation*}
每个\(f_n\)连续，极限函数却在\(1\)处不连续。固定一个\(x<1\)后，\(x^n\)确实趋零；但越靠近\(1\)，达到同一误差要求所需的指标可能越大。

还要注意，选择随\(n\)变化的位置，已经超出了固定点极限所直接保证的范围。比如\(x_n=2^{-1/n}\)虽然始终小于\(1\)，却满足\(x_n^n=1/2\)。这不会反驳每个固定\(x<1\)处的极限，却暴露了不同位置不能同时受控的问题。

\subsection{一致收敛}

\begin{definition}
若对任意\(\varepsilon>0\)，存在\(N\)，使\(n>N\)时，对区间内所有\(x\)同时有
\begin{equation*}
 |f_n(x)-f(x)|<\varepsilon,
\end{equation*}
就称\(f_n\){\heiti{}一致收敛}于\(f\)。
\end{definition}

一致收敛要求同一个\(N\)适用于整个区间，所以必然推出逐点收敛。若\(f_n,f\)都属于\(C[a,b]\)，这一要求等价于\(\|f_n-f\|_\infty\to0\)。

例如，\(f_n(x)=x/(n+1)\)在\([0,1]\)上与零函数的最大误差为\(1/(n+1)\)，故一致收敛。前面的\(x^n\)则不一致收敛于其逐点极限：取\(x_n=2^{-1/n}<1\)，始终有\(|f_n(x_n)-f(x_n)|=1/2\)。

第二章的一致连续，比较同一个函数在两个接近位置的取值；这里的一致收敛，比较一列函数与极限函数在相同位置的取值。两者都要求统一控制，但所变化的对象不同。

判断一致收敛时，可以先寻找不依赖\(x\)的上界\(a_n\)，使\(|f_n(x)-f(x)|\le a_n\to0\)。若要否定一致收敛，则可以寻找一列位置\(x_n\)，让误差始终不小于某个固定正数。两种方法分别对应定义中的统一上界与反例。

\begin{figure}[!ht]
\centering
\includegraphics[width=.87\linewidth]{calculus-convergence}
\caption{左图\(x^n\)逐点收敛，靠近\(1\)的误差却不能统一变小；右图\(x/(n+1)\)的整条图像最终进入任意给定的零函数误差带。}
\label{fig:calculus-convergence}
\end{figure}

例如，\(f_n(x)=xe^{-nx}\)在\([0,1]\)上逐点趋零。对\(x\)求导可知最大值在\(x=1/n\)取得，故\(\|f_n\|_\infty=1/(en)\to0\)。最大误差的位置也在移动，但误差本身趋零，所以这次确实是一致收敛。

\begin{exercise}
\mbox{}\par\nobreak

（1）判断\(x^n\)在\([0,r]\)、\(0<r<1\)上是否一致收敛，并给出误差界。

（2）求\(f_n(x)=x/(1+nx)\)在\([0,1]\)上的逐点极限，判断是否一致收敛。

（3）对\(f_n(x)=x^n\)、\(x\in[0,1]\)，说明为什么逐点极限\(f\)不能直接当作\(C[0,1]\)中的元素。使用一致收敛的定义判断时，是否需要先假定\(f\)连续？
\end{exercise}

\subsection{一致极限的连续性}

\begin{theorem}\label{thm:calculus-uniform-limit-continuity}
若连续函数列\(f_n\)在\([a,b]\)上一致收敛于\(f\)，则\(f\)连续。
\end{theorem}

\begin{proof}
固定\(x_0\)及\(\varepsilon>0\)。先选定一个足够靠后的\(f_N\)，使处处有\(|f_N-f|<\varepsilon/3\)。再由\(f_N\)的连续性取\(\delta>0\)，使\(|x-x_0|<\delta\)时有\(|f_N(x)-f_N(x_0)|<\varepsilon/3\)。于是
\begin{equation*}
 \begin{aligned}
 |f(x)-f(x_0)|&\le|f(x)-f_N(x)|+|f_N(x)-f_N(x_0)|+|f_N(x_0)-f(x_0)|<\varepsilon.
 \end{aligned}
\end{equation*}
故\(f\)在\(x_0\)处连续。
\end{proof}

这里先固定同一个近似函数，再移动\(x\)。一致收敛使两端的误差都受同一个指标控制；若只有逐点收敛，指标可能随\(x\)变化，就无法直接使用这个固定函数的连续性。

证明并未要求所有\(f_n\)在\(x_0\)附近使用同一个连续性范围，只须其中一个已经选定的函数连续。两侧的近似误差与中间的连续性误差，分别由先选指标、后选位置范围控制。这种先后次序使函数列的逼近与一个函数的局部变化分开处理。

\begin{example}
在\([-1,1]\)上，令\(f_n(x)=\sqrt{x^2+1/n}\)。证明它一致收敛于\(|x|\)，并说明极限连续而在零点不可导，与本小节定理是否矛盾。可先估计\(0\le\sqrt{x^2+1/n}-|x|\le1/\sqrt n\)。
\end{example}

\subsection{一致逼近的积分误差}

对\(f,g\in C[a,b]\)，主线的积分估计给出
\begin{equation}\label{eq:calculus-function-integral-bound}
 \left|\int_a^b f(x)\,dx-\int_a^b g(x)\,dx\right|
 \le(b-a)\|f-g\|_\infty.
\end{equation}
因此，连续函数列一致收敛于\(f\)时，先由上一小节知\(f\)连续，再由此估计得到
\begin{equation*}
 \lim_{n\to\infty}\int_a^b f_n(x)\,dx=\int_a^b f(x)\,dx.
\end{equation*}
整个函数的误差趋零，使对应累积量的误差也趋零。

例如，\(e^x\)在\([0,1]\)上的Taylor多项式\(T_n\)满足
\begin{equation*}
 \|e^x-T_n(x)\|_\infty\le\frac e{(n+1)!}.
\end{equation*}
同一个余项界在整个区间上成立，所以既控制函数近似，也控制积分近似。这里的统一估计，比只知道每个固定点的余项趋零提供了更多信息。

积分也可以保留变化的终点。对\(A_n(x)=\int_a^xf_n(t)\,dt\)、\(A(x)=\int_a^xf(t)\,dt\)，有
\begin{equation*}
 |A_n(x)-A(x)|\le|x-a|\|f_n-f\|_\infty
 \le(b-a)\|f_n-f\|_\infty.
\end{equation*}
同一个右端同时控制全部终点，所以得到的是累积函数的一致逼近，而不只是一个固定积分值的收敛。

若累积中带有连续权函数\(w\)，还可用\(\int_a^b|w|\)代替区间长度：
\begin{equation*}
 \left|\int_a^bw(x)(f_n(x)-f(x))\,dx\right|
 \le\|f_n-f\|_\infty\int_a^b|w(x)|\,dx.
\end{equation*}
权的总绝对贡献决定了函数误差对结果的放大程度。

\begin{exercise}
\mbox{}\par\nobreak

（1）用\(e^x\)的二阶Taylor多项式近似\(\int_0^1e^x\,dx\)，利用最大误差给出积分误差界，再与精确结果比较。

（2）设\(f_n\)在\([a,b]\)上一致收敛于\(f\)，各函数连续。固定\(c\in[a,b]\)，证明\(\int_c^xf_n(t)\,dt\)一致收敛于\(\int_c^xf(t)\,dt\)，并给出同时适用于所有起点\(c\)和终点\(x\)的误差界。

（3）设\(f_n(x)=nx(1-x)^n\)、\(0\le x\le1\)。证明\(f_n\)逐点趋于零，且\(\int_0^1f_n(x)\,dx\to0\)，但\(f_n\)不一致收敛于零。说明一致收敛在积分极限结论中是充分条件，却不是必要条件。
\end{exercise}

\subsection{度量空间}

实数差的绝对值、向量之差的长度、函数之间的最大误差，都在量化两个对象的接近程度。现在从这些已有对象中提取共同条件。

\begin{definition}
集合\(X\)上的函数\(d:X\times X\to\mathbb R\)，若满足
\begin{equation*}
 d(x,y)\ge0,\quad d(x,y)=0\Longleftrightarrow x=y;\qquad
 d(x,y)=d(y,x);\qquad
 d(x,z)\le d(x,y)+d(y,z),
\end{equation*}
就称为\(X\)上的一个{\heiti{}度量}，称\((X,d)\)为{\heiti{}度量空间}。
\end{definition}

范数产生度量\(d(u,v)=\|u-v\|\)，三项条件由范数性质直接得到。因此实数空间、欧氏向量空间以及带最大值范数的\(C[a,b]\)，都可以用距离描述接近。

以三角不等式为例，\(u-w=(u-v)+(v-w)\)，所以\(d(u,w)\le d(u,v)+d(v,w)\)。这一形式与第二章证明连续性时拆分误差的步骤相同，只是现在的点也可以是函数。

在度量空间中，\(x_n\to x\)表示\(d(x_n,x)\to0\)。因而同一个定义既包含实数数列，也包含向量列与函数列。不过距离必须先确定：在最大值距离下，函数列收敛要求全区间同时接近；逐点收敛只逐个比较位置，不能直接替代这个距离条件。

例如，函数列在最大值距离下趋于\(f\)，正是\(d(f_n,f)=\|f_n-f\|_\infty\to0\)，也就是前面研究的一致收敛。距离语言把已经完成的误差比较组织为同一种表述。

\subsection{连续映射}

\begin{definition}
设\((X,d_X)\)、\((Y,d_Y)\)为度量空间。若对任意\(\varepsilon>0\)，存在\(\delta>0\)，使
\begin{equation*}
 d_X(x,x_0)<\delta\quad\Longrightarrow\quad
 d_Y(T(x),T(x_0))<\varepsilon,
\end{equation*}
就称映射\(T:X\to Y\)在\(x_0\)处{\heiti{}连续}。
\end{definition}

当输入输出都是实数，距离取绝对值时，这就是第二章的连续性定义。现在也可以让输入是整个函数。

固定\(c\in[a,b]\)，点值映射\(E_c:C[a,b]\to\mathbb R\)、\(E_c(f)=f(c)\)满足
\begin{equation*}
 |E_c(f)-E_c(g)|\le\|f-g\|_\infty,
\end{equation*}
所以连续。积分映射\(J(f)=\int_a^bf\)由式\eqref{eq:calculus-function-integral-bound}也连续。两者还都是线性映射，但连续性的证明来自距离估计，线性性本身不替代这一估计。

连续映射也可以是非线性的。例如，\(M(f)=\max_{[a,b]}f\)把函数送到它的最大值。由\(f(x)\le g(x)+\|f-g\|_\infty\)，先对\(x\)取最大值，再交换\(f,g\)，得到
\begin{equation*}
 |M(f)-M(g)|\le\|f-g\|_\infty.
\end{equation*}
所以整个函数小幅变化时，最值也只会小幅变化。这并不意味着最值的位置同样稳定，最值和取得最值的位置是不同的对象。

\begin{exercise}
\mbox{}\par\nobreak

（1）在\(C[0,1]\)上取\(f(x)=x\)、\(g(x)=1-x\)，说明最大值映射不满足可加性。

（2）设\(f_n(x)=x/n\)，\(g_n(x)=-x/n\)，\(0\le x\le1\)。两列函数之间的最大误差趋零，它们各自唯一最大值点之间的距离是否趋零？比较这与最大值映射的连续性。
\end{exercise}

\begin{example}
对固定连续函数\(w\)，证明\(J_w(f)=\int_a^bw(x)f(x)\,dx\)是\(C[a,b]\)到\(\mathbb R\)的连续线性映射，并给出误差界。
\end{example}

\subsection{微积分的线性算子}

记\(C^1[a,b]\)为一阶导数连续的函数组成的线性空间，端点导数按单侧理解。定义
\begin{equation*}
 D:C^1[a,b]\to C[a,b],\quad DF=F',\qquad
 I:C[a,b]\to C^1[a,b],\quad(If)(x)=\int_a^xf(t)\,dt.
\end{equation*}
这两个作用在函数空间上的线性映射，称为{\heiti{}线性算子}。积分算子的值确实属于\(C^1[a,b]\)，其依据是微积分基本定理。

在明确输入输出之后，基本定理及Newton--Leibniz公式分别给出
\begin{equation*}
 DIf=f\quad(f\in C[a,b]),\qquad
 IDF=F-F(a)\quad(F\in C^1[a,b]).
\end{equation*}
这里\(F(a)\)作为一个常值函数参与相减。求导的核是所有常值函数，因此从导数恢复函数时，必须另外确定一个常数。

若把求导的定义域限制为
\begin{equation*}
 C^1_a=\{F\in C^1[a,b]:F(a)=0\},
\end{equation*}
则\(D:C^1_a\to C[a,b]\)与\(I:C[a,b]\to C^1_a\)互为逆映射。微积分的互逆关系在这个明确的范围内成立，其依据是前面分别定义并证明的两条复合恒等式。

例如，在\([0,1]\)上取\(F(x)=x^2+3\)，则\(DF=2x\)，从零点积分回来得到\(IDF=x^2\)。所缺的常数正是\(F(0)=3\)。若改从\(a\)累积，缺失的量就相应变为\(F(a)\)，所以起点不仅是记号选择，也决定了积分恢复哪个代表函数。

从线性空间看，任意\(F\in C^1[a,b]\)都有唯一分解
\begin{equation*}
 F=\bigl(F-F(a)\bigr)+F(a),
\end{equation*}
前一部分属于\(C^1_a\)，后一部分属于常值函数空间。求导保留前一部分的信息，而把后一部分送到零函数。

\begingroup
\setlength{\parskip}{4pt}
\setlength{\abovedisplayskip}{10pt}
\setlength{\belowdisplayskip}{10pt}
\subsection{求导对近似的影响}

在\([-1,1]\)上考察\(f_n(x)=\sin(nx)/n\)。由于\(\|f_n\|_\infty\le1/n\)，它一致收敛于零函数，但
\begin{equation*}
 f_n'(x)=\cos(nx),\qquad f_n'(0)=1.
\end{equation*}
导函数甚至在零点也不趋于零函数的导数。函数值的振幅虽不断缩小，变化却越来越快，因此小的函数值误差不自动给出小的导数误差。

\begin{figure}[!ht]
\centering
\includegraphics[width=.87\linewidth]{calculus-derivative-error}
\caption{左图\(\sin(nx)/n\)的振幅趋零，右图导函数\(\cos(nx)\)的振幅仍为\(1\)。图中取\(n=5,12\)。}
\label{fig:calculus-derivative-error}
\end{figure}

从差商也能看到误差为什么会放大。若\(\|f-g\|_\infty\le\eta\)，则对\(h\ne0\)，
\begin{equation*}
 \left|\frac{f(x+h)-f(x)}h-\frac{g(x+h)-g(x)}h\right|
 \le\frac{2\eta}{|h|}.
\end{equation*}
当计算局部变化率所用的\(h\)缩小时，这个由函数值误差得到的界反而增大。因此仅让函数图像靠近，不足以保证斜率靠近。

反过来，若同时掌握导函数误差与一个初值误差，则基本定理给出
\begin{equation*}
 \|F-G\|_\infty\le|F(a)-G(a)|+(b-a)\|F'-G'\|_\infty.
\end{equation*}
导数控制各段变化，初值确定整体高度。这与5.3中“导数相同仍可能相差常数”的认识在近似问题中形成对应。

相比之下，积分误差受区间长度乘最大误差控制。求导和积分都是线性运算，但它们对同一种函数近似的反应不同；要在求导后继续控制近似，必须掌握比函数值误差更多的信息。

\begin{exercise}
\mbox{}\par\nobreak

（1）求算子\(D\)的核与像，说明固定起点的积分算子为什么是单射。

（2）令\(QF=F-F(a)\)。证明\(Q^2=Q\)，求\(Q\)的核与像，并把它与\(ID\)联系起来。

（3）在\(C^1[-1,1]\)上，以函数最大值距离比较输入，以\(C[-1,1]\)的同类距离比较输出。用\(f_n(x)=\sin(nx)/n\)证明求导映射在零函数处不连续。

（4）设\(F_n\in C^1[a,b]\)，\(F_n(a)\to c\)，且\(F_n'\)一致收敛于连续函数\(g\)。利用\(F_n(x)=F_n(a)+\int_a^xF_n'(t)\,dt\)，证明\(F_n\)一致收敛，并求极限函数的导数。

（5）在\([0,1]\)上令\(f_n(x)=\cos(nx)\)。证明\(f_n\)不一致趋于零，但\(If_n\)一致趋于零。由此说明积分结果的接近，能否反过来保证被积函数接近？

（6）设\(F,G\in C^1[a,b]\)，且\(F(c)=G(c)\)，其中\(c\in[a,b]\)。证明\(\|F-G\|_\infty\le\max\{c-a,b-c\}\|F'-G'\|_\infty\)。怎样选择\(c\)，可以使右侧的常数最小？

（7）给定\(g\in C[a,b]\)，证明条件\(F'=g\)、\(\int_a^bF(x)\,dx=0\)唯一确定一个\(F\in C^1[a,b]\)，并用\(Ig\)表示\(F\)。比较这一条件与\(F(a)=0\)确定积分常数的方式。
\end{exercise}
\endgroup

\FloatBarrier
\newpage
\section{积分的应用}
\label{sec:calculus-integral-applications}

\begingroup
\setlength{\parskip}{3pt}
\setlength{\abovedisplayskip}{10pt}
\setlength{\belowdisplayskip}{10pt}

积分应用的共同任务，是确定每小段贡献什么、沿哪个变量累积，再把局部贡献相加。不同问题的表达式来自数量关系，累积过程则由同一种积分运算完成。

\subsection{平面图形的面积}

设连续曲线\(y=f(x)\)、\(y=g(x)\)在\([a,b]\)上满足\(f\ge g\)。第\(i\)条窄带的宽度为\(\Delta x_i\)，高度可以近似为\(f(\xi_i)-g(\xi_i)\)，故区域面积由矩形和的极限给出：
\begin{equation}\label{eq:calculus-area-between-graphs}
 S=\int_a^b(f(x)-g(x))\,dx.
\end{equation}
若上下关系改变，则先求交点，再分段；也可以统一写为\(\int_a^b|f-g|\)。若区域更适合按水平方向切片，就改用\(y\)为积分变量，以右边界横坐标减去左边界横坐标作为窄带长度。

例如，\(y=x\)与\(y=x^2\)相交于\(x=0,1\)，在中间直线位于抛物线上方，所以
\begin{equation*}
 S=\int_0^1(x-x^2)\,dx=\frac16.
\end{equation*}

\begin{figure}[!ht]
\centering
\includegraphics[width=.44\linewidth]{calculus-area-slices}
\caption{两条曲线之间的竖直窄带：局部面积由高度差乘横向宽度近似。}
\label{fig:calculus-area-slices}
\end{figure}

切片方向影响计算的简便程度。若竖直切片截出几段，改用水平切片后可能只剩一段。因此应根据区域的形状选择积分变量。

\begin{example}
求\(y=x^2\)与\(y=2-x^2\)围成的面积。再求椭圆\(x^2/a^2+y^2/b^2=1\)内部的面积，利用\(x=au\)的换元与第四章伸缩变换的面积倍率相互核对。
另求\(y=|x|\)与\(y=x^2\)围成的总面积，说明积分区间怎样由交点确定。
\end{example}

\begin{exercise}
\mbox{}\par\nobreak

（1）求\(y=x^3\)与\(y=x\)围成的全部有限区域的总面积，说明为什么不能只计算\(\int_{-1}^1(x-x^3)\,dx\)。

（2）求\(x=y^2\)与\(x=2-y\)围成的面积，比较水平与竖直切片的计算。

（3）对\(0<a<1\)，求\(y=x^2\)、横轴及直线\(x=a\)、\(x=1\)围成区域的面积，并求\(a\)，使它恰为\([0,1]\)上对应总面积的一半。

（4）求\(y=\ln x\)、横轴及直线\(x=e\)围成的面积。再求\(y=e^x\)、纵轴及直线\(y=e\)围成的面积，比较两区域关于\(y=x\)的对称关系。
\end{exercise}

\subsection{总变化量}

若位置函数\(s\)连续可导，速度为\(v=s'\)，则基本定理给出
\begin{equation*}
 s(b)-s(a)=\int_a^b v(t)\,dt.
\end{equation*}
这是位移。在不改变运动方向的时间段，路程等于位移的绝对值；各段相加，得到
\begin{equation*}
 \text{总路程}=\int_a^b|v(t)|\,dt.
\end{equation*}
速度变号的位置决定怎样拆分，而积分的上下限由所研究的时间范围决定。

位移累积各段有正负的变化，路程累积各段经过的长度。

同样，若\(q(t)\)表示每单位时间流入的量，则\(q(t)\Delta t\)近似一小段时间的流入量，\(\int_a^bq(t)\,dt\)给出总流入量。若同时有流出，则净变化应使用流入率减去流出率。

对细杆，设\(\rho(x)\)表示每单位长度的质量，称为线密度。短段质量近似为\(\rho(x)\Delta x\)，故整根杆的质量为\(\int_a^b\rho(x)\,dx\)。被积量与微分长度的单位相乘，应与所求总量一致。

若物体沿直线运动，沿运动坐标方向的力为\(F(x)\)，恒力在位移\(\Delta x\)上做功\(F\Delta x\)。变力可以在短段内用一个取样值近似，所以从\(a\)运动到\(b\)的功为
\begin{equation*}
 W=\int_a^bF(x)\,dx.
\end{equation*}
例如，弹簧的拉力为\(F(x)=kx\)，其中\(x\)是伸长量，功由各段的力乘位移累积。

\begin{example}
质点速度为\(v(t)=t^2-2t\)，\(0\le t\le3\)，初始位置为\(s(0)=1\)。求位置函数、总位移和总路程。再设一个容器的流入率为\(2+t\)、流出率为\(t^2\)，求\([0,2]\)内储量的净变化，并判断储量何时最大。
\end{example}

若位置为\(x=x(t)\)，由换元公式，功也可写成\(\int F(x(t))x'(t)\,dt\)。因子\(x'(t)\)将时间增量转为有方向的位移增量。

\begin{example}
一根细杆位于\([0,L]\)，从左端到位置\(x\)的累计质量为\(M(x)=kx^2\)，其中\(k>0\)。求线密度及区间\([L/2,L]\)内的质量，用密度积分核对。累计质量与线密度记录了什么？
\end{example}

\subsection{加权平均}

将不同部分赋予不同的非负权重，可以研究整体的平均位置或平均水平。设\(w\)连续、非负，且\(W=\int_a^bw(x)\,dx>0\)，定义\(f\)的{\heiti{}加权平均值}为
\begin{equation}\label{eq:calculus-weighted-average}
 \overline f_w=\frac{\int_a^bf(x)w(x)\,dx}{\int_a^bw(x)\,dx}.
\end{equation}
分子累积各段的加权贡献，分母累积总权重。权为正常数时，就得到普通积分平均值。

从有限分割看，若第\(i\)小段的权为\(w(\xi_i)\Delta x_i\)，离散加权平均就是
\begin{equation*}
 \frac{\sum_i f(\xi_i)w(\xi_i)\Delta x_i}{\sum_iw(\xi_i)\Delta x_i}.
\end{equation*}
分割加细后，分子与分母分别趋于对应积分，便得到式\eqref{eq:calculus-weighted-average}。

若\(m\le f\le M\)，将不等式乘以\(w\ge0\)并积分，得到\(m\le\overline f_w\le M\)。若\(f\)连续，再由介值定理，存在\(\xi\in[a,b]\)，使\(\overline f_w=f(\xi)\)。

例如，线密度为\(\rho\)的细杆可看成许多小质量的组合。其平均位置，即质心坐标，为
\begin{equation*}
 \overline x=\frac{\int_a^bx\rho(x)\,dx}{\int_a^b\rho(x)\,dx}.
\end{equation*}
它由位置乘对应质量的总和除以总质量得到，与有限个质点的质量加权平均一致。

合并两个部分时也应按各自总权重加权。若两段总权重为\(W_1,W_2\)，平均值为\(m_1,m_2\)，则由积分的区间可加性，整体平均值为\((W_1m_1+W_2m_2)/(W_1+W_2)\)。只有两部分总权重相等时，才能对两个平均值直接取算术平均。

\begin{example}
一根细杆占据\([0,L]\)，线密度为\(\rho(x)=k(1+x/L)\)，其中\(k,L>0\)。求总质量与质心位置。再说明同时将所有位置的密度乘以同一个正常数时，质心为什么不变。
\end{example}

\begin{exercise}
\mbox{}\par\nobreak

（1）设\(v(t)=\cos t\)，\(0\le t\le2\pi\)。求位移、路程和平均速度，并与平均速率比较。

（2）杆的线密度满足\(\rho(a+b-x)=\rho(x)\)。用积分换元证明质心位于中点。

（3）设\(w\ge0\)、\(\int_a^bw>0\)。证明\(\int_a^b(f-c)^2w\)在\(c=\overline f_w\)时最小，并解释加权平均与平方偏差之间的关系。

（4）某弹簧的拉力为\(F(x)=kx\)、\(k>0\)。求将其伸长量从\(a\)增加到\(b\)所需的功，其中\(0\le a<b\)。比较从\(0\)到\(a\)与从\(a\)到\(2a\)的功，解释为何相等位移未必对应相等的功。

\end{exercise}

\endgroup
\FloatBarrier
\newpage
\subsection{参数曲线围成的面积*}

直角坐标下的面积可以经换元写成参数形式。若一段曲线表示为\(x=\varphi(t)\)、\(y=\psi(t)\)，且\(x\)沿该段递增，则曲线与横轴之间的有向面积为
\begin{equation*}
 \int y\,dx=\int\psi(t)\varphi'(t)\,dt.
\end{equation*}
若方向反过来，积分相应变号。

以下设闭曲线除首尾相接外不自交，其坐标在有限分段上连续可导，并能拆成横坐标单调的弧段。沿逆时针方向绕区域一周，在各条带上配对上、下边界，累加切片面积公式，得到
\begin{equation*}
 S=-\oint y\,dx.
\end{equation*}
符号\(\oint\)表示沿闭曲线一周，把各参数段的积分相加。逆时针时下边界从左向右、上边界从右向左，因而负号恰好使上、下高度之差为正。

一条带的下、上边界若为\(g,f\)，横坐标范围为\([u,v]\)，沿边界的贡献为
\begin{equation*}
 \int_u^vg(x)\,dx+\int_v^uf(x)\,dx
 =-\int_u^v(f(x)-g(x))\,dx.
\end{equation*}
把各条带相加，得到整个面积；人为切开的边界两侧方向相反，其贡献相消，竖直连接段的\(dx\)为零。这说明闭曲线公式来自切片面积和边界方向的配合。

又因为一周后\(xy\)回到原值，\(\oint d(xy)=\oint(x\,dy+y\,dx)=0\)，所以也有
\begin{equation}\label{eq:calculus-parametric-area}
 S=\oint x\,dy=\frac12\oint(x\,dy-y\,dx)
 =\frac12\int_\alpha^\beta(\varphi\psi'-\psi\varphi')\,dt.
\end{equation}
本式要求沿边界逆时针走一周。反向走得到负的有向面积，多走几周则重复累积。

对椭圆\(x=a\cos t\)、\(y=b\sin t\)，\(0\le t\le2\pi\)，有\(xy'-yx'=ab\)，故面积为\(\pi ab\)。参数表示把四段边界统一在同一个积分中。

\begin{exercise}
\mbox{}\par\nobreak

（1）对上述椭圆分别取\(t\in[0,4\pi]\)及反向参数，判断式\eqref{eq:calculus-parametric-area}的结果，解释覆盖次数与方向的作用。

（2）设\(r(t)\)表示一条正向简单闭曲线。证明平移后面积不变；再用\(x=au+cv\)、\(y=bu+dv\)代入面积公式，核对有向面积乘以\(ad-bc\)的规律。
\end{exercise}

\subsection{极坐标下的面积*}

设\(r=r(\theta)\ge0\)连续，\(\alpha\le\theta\le\beta\)，其中\(\beta-\alpha\le2\pi\)。研究每条射线上从原点到\(r(\theta)\)之间的区域。把角度分成许多小段，在每段用最小、最大半径作内外扇形。

弧度制下，半径为\(r\)、圆心角为\(\Delta\theta\)的扇形面积是\(r^2\Delta\theta/2\)。因此内外扇形和分别是\(r^2/2\)的下和与上和。连续性使两者趋于同一积分，得到
\begin{equation}\label{eq:calculus-polar-area}
 S=\frac12\int_\alpha^\beta r(\theta)^2\,d\theta.
\end{equation}
角度区间决定覆盖的射线，半径决定每条射线上的长度。若参数重复描绘同一区域，应先选出恰好覆盖一次的范围。

\begin{figure}[!ht]
\centering
\includegraphics[width=.47\linewidth]{calculus-polar-area}
\caption{用取样半径作小扇形，局部面积近似为\(r^2\Delta\theta/2\)，再沿角度累积。}
\label{fig:calculus-polar-area}
\end{figure}

当\(r\)连续可导时，可用参数曲线公式核对。令\(x=r\cos\theta\)、\(y=r\sin\theta\)，求导得
\begin{equation*}
 xy'-yx'=r^2(\cos^2\theta+\sin^2\theta)=r^2.
\end{equation*}
两条连接原点的径向边界不贡献\(x\,dy-y\,dx\)，所以参数面积公式也给出式\eqref{eq:calculus-polar-area}。扇形和仅要求半径函数连续，参数求导则联系了更光滑情形下的两种面积表示。

\begin{example}
求心形线\(r=a(1+\cos\theta)\)、\(0\le\theta\le2\pi\)围成的面积。再用\(r=2a\cos\theta\)、\(-\pi/2\le\theta\le\pi/2\)计算圆的面积，并与直角坐标下的圆方程核对。
\end{example}

\begin{exercise}
\mbox{}\par\nobreak

（1）求曲线\(r=a\cos2\theta\)在\(-\pi/4\le\theta\le\pi/4\)上所围成一瓣的面积，再由对称性求全部四瓣的总面积。

（2）若区域在每条射线上介于\(r_1(\theta)\)与\(r_2(\theta)\)之间，且\(0\le r_1\le r_2\)，推导其面积公式。用同心圆环核对。
\end{exercise}

\subsection{体积*}

对本节研究的几何体，可以把沿某个方向的薄层近似为柱体。设在位置\(x\)处垂直于横轴的截面积为连续函数\(A(x)\)，薄层厚度为\(\Delta x\)，则局部体积近似为\(A(x)\Delta x\)。由薄层和的极限，得到
\begin{equation}\label{eq:calculus-volume}
 V=\int_a^b A(x)\,dx.
\end{equation}
需要先确定截面，再计算它的面积；截面形状不同，对应的\(A(x)\)也不同。

\paragraph{旋转体的截面积} \mbox{}

若将\(0\le y\le f(x)\)、\(a\le x\le b\)的区域绕横轴旋转，垂直于旋转轴的截面为半径\(f(x)\)的圆盘，因此
\begin{equation*}
 V=\pi\int_a^bf(x)^2\,dx.
\end{equation*}
若\(0\le g(x)\le f(x)\)，旋转的是两曲线之间的区域，则截面为圆环，外半径、内半径分别为\(f(x),g(x)\)，故
\begin{equation}\label{eq:calculus-washer-volume}
 V=\pi\int_a^b\bigl(f(x)^2-g(x)^2\bigr)\,dx.
\end{equation}
圆环面积是两个圆面积之差，不能把\(f-g\)当作半径后再平方。若区域跨过旋转轴，则应从旋转后的实际截面重新判断半径。

绕纵轴旋转时，改用水平截面更直接：若在\(c\le y\le d\)上，区域介于\(0\le u(y)\le x\le v(y)\)之间，则
\begin{equation*}
 V=\pi\int_c^d\bigl(v(y)^2-u(y)^2\bigr)\,dy.
\end{equation*}
半径始终是到旋转轴的距离。绕\(y=k\)或\(x=k\)旋转，只须根据位置关系重新写出内、外半径；累积变量仍沿轴的方向选取。

例如，半径为\(R\)的球在横坐标\(x\)处截面积为\(\pi(R^2-x^2)\)，所以
\begin{equation*}
 V=\pi\int_{-R}^R(R^2-x^2)\,dx=\frac43\pi R^3.
\end{equation*}
熟悉的体积公式由截面的累积得到。

截面不必是圆。例如，椭球\(x^2/a^2+y^2/b^2+z^2/c^2\le1\)在横坐标\(x\)处的截面为椭圆，两条半轴分别为\(b\sqrt{1-x^2/a^2}\)、\(c\sqrt{1-x^2/a^2}\)。它的截面积是\(\pi bc(1-x^2/a^2)\)，再沿\([-a,a]\)累积，就将椭球体积化成了一个多项式积分。圆与椭圆的截面虽不同，确定截面、求面积、沿厚度累积的次序相同。

\begin{example}
求\(y=x\)与\(y=x^2\)在\([0,1]\)上围成区域绕横轴旋转所得几何体的体积。再用截面积法求底面半径为\(R\)、高为\(h\)的正圆锥体积。
\end{example}

\paragraph{圆柱壳法} \mbox{}

同一个旋转体也可以按到轴的距离分层。先设\(0\le a<b\)、\(f\ge0\)，将\(a\le x\le b\)、\(0\le y\le f(x)\)的区域绕纵轴旋转。竖直窄条旋转后形成薄圆柱壳，其半径接近\(x\)，高度接近\(f(x)\)，厚度为\(\Delta x\)。

对固定高度\(H\)，内、外半径为\(x_{i-1},x_i\)的圆柱壳体积恰为
\begin{equation*}
 \pi(x_i^2-x_{i-1}^2)H
 =2\pi\frac{x_i+x_{i-1}}2H\Delta x_i.
\end{equation*}
在每段分别以\(f\)的最小值、最大值取代\(H\)，便得到内外近似体积。连续性使高度的波动统一趋零，而圆环底面积之和为固定的\(\pi(b^2-a^2)\)，所以两种近似之差趋零，得到
\begin{equation}\label{eq:calculus-shell-volume}
 V=2\pi\int_a^bxf(x)\,dx.
\end{equation}
若旋转的是\(g(x)\le y\le f(x)\)之间的区域，壳高改为\(f-g\)，于是\(V=2\pi\int_a^bx(f-g)\,dx\)。这里区域位于轴的一侧，各半径层只累积一次。

截面积法沿轴累积圆盘或圆环，圆柱壳法垂直于轴累积半径层。两者可以处理同一个几何体，选择哪一种，取决于边界表达和需要拆分的次数。

\begin{figure}[!ht]
\centering
\includegraphics[width=.83\linewidth]{calculus-rotation}
\caption{左图用圆盘截面沿轴累积，右图的竖直窄条绕纵轴生成圆柱壳；半径、截面积与壳高分别由几何关系确定。}
\label{fig:calculus-rotation}
\end{figure}

\begin{example}
将\(y=x\)与\(y=x^2\)围成的区域绕纵轴旋转，分别用圆环截面与圆柱壳计算体积。再将它绕\(x=-1\)旋转，说明圆柱壳的半径与高度分别怎样变化。
\end{example}

\begin{exercise}
\mbox{}\par\nobreak

（1）将\(y=\sqrt{x}\)、横轴及直线\(x=1\)、\(x=4\)围成的区域绕横轴旋转，求体积。再改绕\(y=-1\)旋转，写出圆环的内、外半径。

（2）将\(y=x^2\)、\(y=1\)围成的区域绕纵轴旋转，分别用水平截面和圆柱壳计算。为什么使用圆柱壳时，不能把左右两半区域各自旋转所得体积直接相加？

（3）补完上文椭球体积的计算，核对结果在\(a=b=c\)时是否回到球体积。
\end{exercise}

\subsection{弧长*}

对平面参数曲线\(r(t)=(x(t),y(t))\)，在参数区间\([\alpha,\beta]\)上作分割\(P\)，将曲线上对应的点依次连接，得到折线长度
\begin{equation*}
 L(P)=\sum_i|r(t_i)-r(t_{i-1})|.
\end{equation*}
若当\(\|P\|\to0\)时，这些长度趋于同一个有限极限，就称该极限为这段曲线的{\heiti{}弧长}。

现在设两个坐标函数连续可导。由局部一次近似，短弦长度应接近
\begin{equation*}
 |r'(t)|\Delta t=\sqrt{x'(t)^2+y'(t)^2}\,\Delta t.
\end{equation*}
把它累积，就得到弧长公式
\begin{equation}\label{eq:calculus-arc-length}
 L=\int_\alpha^\beta\sqrt{x'(t)^2+y'(t)^2}\,dt.
\end{equation}

为说明误差如何控制，对每段分别使用坐标函数的中值定理，短弦增量可写为\(\Delta t_i(x'(\xi_i),y'(\eta_i))\)。两个中间点未必相同，但导函数的一致连续性保证，当段长足够小时，这个向量与\(\Delta t_i r'(\tau_i)\)之差的长度不超过\(\varepsilon\Delta t_i\)。由向量长度的三角不等式，弦长与\(|r'(\tau_i)|\Delta t_i\)的差也受此界控制。累积后总误差不超过\(\varepsilon(\beta-\alpha)\)，故折线长度与速度模的Riemann和趋于同一极限，证明了公式。

这里沿用第二章的一致连续性，是将逐点的线性近似用于整个曲线长度的依据。对图像\(y=f(x)\)，取参数为\(x\)，便得到
\begin{equation*}
 L=\int_a^b\sqrt{1+f'(x)^2}\,dx.
\end{equation*}
几何长度累积的是速度的模，而不是某一个坐标变化率。

\begin{figure}[!ht]
\centering
\includegraphics[width=.83\linewidth]{calculus-arc-length}
\caption{左图用折线近似弧长；右图将短弦的两个坐标增量合成为长度，再过渡到弧长积分。}
\label{fig:calculus-arc-length}
\end{figure}

若曲线更适合写成\(x=g(y)\)，则取\(y\)为参数，得到\(L=\int_c^d\sqrt{1+g'(y)^2}\,dy\)。同一段曲线可以有不同的坐标表示，弧长公式始终累积两个坐标变化合成的长度。

\paragraph{弧长函数} \mbox{}

固定参数起点，令
\begin{equation*}
 s(t)=\int_\alpha^t\sqrt{x'(u)^2+y'(u)^2}\,du.
\end{equation*}
它表示沿参数递增方向，从起点到当前位置已经走过的长度。由基本定理，
\begin{equation*}
 s'(t)=\sqrt{x'(t)^2+y'(t)^2},\qquad
 ds=\sqrt{x'(t)^2+y'(t)^2}\,dt.
\end{equation*}
若\(t\)是时间，右端的导数就是速率。这个长度元也使“沿曲线按长度累积”获得具体表达，例如非均匀细线的质量可以写成\(\int\rho\,ds\)，其中\(\rho\)为单位弧长的质量，计算时再代入参数形式。

参数改变时，走过同一位置所需的参数增量会变化，速度模也相应变化。若\(t=\varphi(s)\)、\(\varphi'>0\)，则新的速度模为\(|r'(\varphi(s))|\varphi'(s)\)。换元积分使它沿新参数的累积仍等于原来的弧长。面积公式保留方向，弧长公式使用长度，这解释了反向走同一段曲线时面积积分变号、长度却不变。

\begin{example}
用参数公式求半径为\(R\)的圆周长。再求曲线\(x=t^2\)、\(y=t^3\)、\(0\le t\le1\)的弧长，比较直接使用参数与先消去参数的计算。
\end{example}

\paragraph{极坐标曲线的弧长} \mbox{}

若\(r=r(\theta)\)连续可导，代入\(x=r\cos\theta\)、\(y=r\sin\theta\)，有
\begin{equation*}
 \begin{aligned}
 x'&=r'\cos\theta-r\sin\theta,\\
 y'&=r'\sin\theta+r\cos\theta,
 \end{aligned}
 \qquad x'^2+y'^2=r'^2+r^2.
\end{equation*}
交叉项相消，因而
\begin{equation}\label{eq:calculus-polar-arc-length}
 L=\int_\alpha^\beta\sqrt{r(\theta)^2+r'(\theta)^2}\,d\theta.
\end{equation}
其中\(r'\)记录半径的变化，\(r\)对应角度变化造成的横向位移。即使\(r\)为常数，角度变化仍产生长度，因此弧长不能只累积\(|r'|\)。

取\(r\equiv R\)，公式给出\(L=R(\beta-\alpha)\)，与圆弧的弧度定义一致；对一般曲线，还需要同时计入半径的变化。

\begin{example}
求心形线\(r=a(1+\cos\theta)\)、\(0\le\theta\le2\pi\)的周长，并与前面的面积计算比较。化简根式时注意\(\sqrt{\cos^2(\theta/2)}=|\cos(\theta/2)|\)。
\end{example}

\begin{exercise}
\mbox{}\par\nobreak

（1）求\(x=y^2\)、\(1\le y\le2\)的弧长，比较以\(x\)为变量和以\(y\)为变量的计算。

（2）摆线的一拱由\(x=a(t-\sin t)\)、\(y=a(1-\cos t)\)、\(0\le t\le2\pi\)表示。求其弧长。

（3）对同一段光滑曲线作保持方向的连续可导参数替换，且替换导数为正。用换元公式证明弧长不变，并与参数曲线切线方向的研究联系起来。

（4）一根细线沿四分之一圆周\(x=R\cos t\)、\(y=R\sin t\)、\(0\le t\le\pi/2\)放置，单位弧长的质量为\(\rho(t)=k(1+\sin t)\)。求总质量，并说明为什么不能直接计算\(\int_0^{\pi/2}\rho(t)\,dt\)。

（5）求极坐标曲线\(r=a\theta\)、\(0\le\theta\le1\)的弧长。把结果与径向距离的总增量比较。
\end{exercise}

\subsection{曲率*}
\label{subsec:calculus-curvature}

第四章用切线和法线研究二次曲线，一阶导数进一步给出了切线方向。不同曲线在一点可以有同一条切线，却以不同的程度离开它。若要比较这种弯曲，就需要研究切线方向沿曲线怎样改变。

同一段曲线可以快走，也可以慢走，参数变化的快慢并不改变曲线本身。因此先用已经建立的弧长衡量走过的距离，再比较单位弧长内切线方向的变化。

\begin{definition}
设平面曲线\(r(t)=(x(t),y(t))\)的坐标二阶连续可导，且\(r'(t)\ne0\)。记速率\(v(t)=|r'(t)|\)，弧长函数为\(s(t)\)，则\(ds/dt=v(t)>0\)。定义单位切向量
\begin{equation*}
 T(t)=\frac{r'(t)}{v(t)},
\end{equation*}
称\(\kappa=|dT/ds|\)为该点的{\heiti{}曲率}。若\(\kappa>0\)，称\(\rho=1/\kappa\)为{\heiti{}曲率半径}。
\end{definition}

局部选取切线方向角\(\theta\)，使\(T=(\cos\theta,\sin\theta)\)，则\(|dT/ds|=|d\theta/ds|\)。所以曲率也就是切线方向角对弧长的变化率的绝对值。直线的切向不变，曲率为零；半径\(R\)的圆走过弧长\(R\Delta\theta\)时，切线方向也转过\(\Delta\theta\)，故曲率恒为\(1/R\)。

\paragraph{参数公式} \mbox{}

由\(x'=v\cos\theta\)、\(y'=v\sin\theta\)，对参数求导并消去\(v'\)，得到
\begin{equation*}
 x'y''-y'x''=v^2\theta'.
\end{equation*}
又因为\(ds/dt=v\)，于是
\begin{equation}\label{eq:calculus-parametric-curvature}
 \kappa=\frac{|\theta'|}{v}
 =\frac{|x'y''-y'x''|}{(x'^2+y'^2)^{3/2}}.
\end{equation}
分子是第三章中两个平面向量\(r',r''\)的行列式的绝对值；它记录速度变化中使方向转动的部分。分母消除了参数速度的影响。反向走同一曲线时，方向角变化的符号可以改变，曲率的大小仍相同。

平移不改变坐标导数，旋转则保持向量长度和行列式，因而这两种变换都保持曲率。第四章中改变坐标原点或旋转坐标轴，不会改变所研究曲线的弯曲程度。

对图像\(y=f(x)\)，取\(x\)为参数，得到
\begin{equation}\label{eq:calculus-graph-curvature}
 \kappa=\frac{|f''(x)|}{(1+f'(x)^2)^{3/2}}.
\end{equation}
因此，\(f''\)的大小本身并不能直接充当弯曲程度；还要计入斜率对弧长的影响。在水平切线处，公式才简化为\(\kappa=|f''|\)。

\paragraph{密切圆} \mbox{}

在曲率非零的点\(P\)，令\(N=(dT/ds)/\kappa\)。由\(T\cdot T=1\)求导，知\(T\cdot N=0\)，且\(|N|=1\)。它指向切线方向正在转向的一侧。取
\begin{equation*}
 C=P+\rho N,
\end{equation*}
以\(C\)为圆心、\(\rho\)为半径作圆，称为曲线在\(P\)处的{\heiti{}密切圆}，也称{\heiti{}曲率圆}。

圆心位于法线上，半径也有确定的选择。以\(P\)处为弧长起点，曲线的二阶Taylor展开为
\begin{equation*}
 r(s_0+h)=P+Th+\frac12\kappa Nh^2+o(h^2).
\end{equation*}
上述圆可在\(P\)附近写为
\begin{equation*}
 c(h)=P+\rho\sin(h/\rho)T+\rho\bigl(1-\cos(h/\rho)\bigr)N.
\end{equation*}
展开后同样得到\(P+Th+\kappa Nh^2/2+o(h^2)\)。切线匹配了一阶方向，密切圆还匹配二阶弯曲；这使第四章的圆与法线，成为一般曲线局部近似的工具。

\paragraph{圆锥曲线中的曲率} \mbox{}

对椭圆\(x=a\cos t\)、\(y=b\sin t\)，其中\(a>b>0\)，有\(x'y''-y'x''=ab\)，代入式\eqref{eq:calculus-parametric-curvature}得
\begin{equation*}
 \kappa(t)=\frac{ab}{(a^2\sin^2t+b^2\cos^2t)^{3/2}}.
\end{equation*}
长轴端点的曲率半径为\(b^2/a\)，短轴端点则为\(a^2/b\)。前者弯曲较急，后者较平缓。把圆沿两个方向作不同倍率的伸缩，虽然仍保持切线和接触关系，却不保持圆形与曲率；曲率还依赖长度和角度的度量。

对抛物线\(y^2=2px\)、\(p>0\)，沿用第四章的参数\(x=pt^2/2\)、\(y=pt\)，得到
\begin{equation*}
 \kappa(t)=\frac1{p(1+t^2)^{3/2}}.
\end{equation*}
顶点处的曲率半径为\(p\)，密切圆圆心为\((p,0)\)。它与焦点\((p/2,0)\)位于同一条轴上，但由不同的几何关系确定：焦点参与整条曲线的定义，密切圆描述顶点处的局部弯曲。

\begin{figure}[!ht]
\centering
\includegraphics[width=.83\linewidth]{calculus-curvature}
\caption{椭圆长轴端点及抛物线顶点处的密切圆。圆心位于法线上，圆与曲线在接触点匹配切向和曲率。}
\label{fig:calculus-curvature}
\end{figure}

\begin{example}
对椭圆\(x^2/4+y^2=1\)，求长、短轴端点处的密切圆方程。再求抛物线\(y^2=4x\)顶点处的密切圆，比较圆心与焦点的位置。
\end{example}

\begin{exercise}
\mbox{}\par\nobreak

（1）将圆\(x=R\cos t\)、\(y=R\sin t\)代入参数曲率公式，核对\(\kappa=1/R\)。再求\(y=x^2\)在原点及\(x=1\)处的曲率，比较它与常值二阶导数\(2\)。

（2）用\(x=a\cosh t\)、\(y=b\sinh t\)求双曲线\(x^2/a^2-y^2/b^2=1\)的曲率。求右顶点处密切圆的圆心与半径，并与椭圆长轴端点比较。

（3）对保持方向的二阶连续可导参数替换\(t=\varphi(u)\)、\(\varphi'>0\)，由式\eqref{eq:calculus-parametric-curvature}证明曲率不变。二阶求导产生的哪一项在行列式中消失？

（4）将平面图形整体放大\(c>0\)倍，证明曲率变为原来的\(1/c\)，曲率半径变为原来的\(c\)倍。再用圆变椭圆的例子，说明第四章的一般仿射变换为什么没有同一个曲率倍率。

（5）椭圆\(x^2/a^2+y^2/b^2=1\)、\(a>b>0\)在哪些点弯曲最急，哪些点最平缓？求两处曲率的比值，并讨论当\(a/b\)增大时怎样变化。
\end{exercise}

\subsection{旋转曲面的面积*}

曲线绕直线旋转时，每个点扫过一个圆，整条母线生成一片曲面。旋转体体积累积的是截面积乘厚度，曲面面积则需要同时考虑旋转半径与母线的长度。以下选取位于旋转轴一侧、恰好生成曲面一次的母线。

\paragraph{从圆台侧面积到曲面面积} \mbox{}

先设\(f\ge0\)且连续可导，将\(y=f(x)\)、\(a\le x\le b\)绕横轴旋转。用母线的折线近似曲线，每条线段旋转后生成圆台侧面；若两端半径相同，就是圆柱侧面。

第\(i\)段母线长为
\begin{equation*}
 \ell_i=\sqrt{(\Delta x_i)^2+(f(x_i)-f(x_{i-1}))^2},
\end{equation*}
由圆台侧面积公式，该段旋转后的面积为\(\pi(f(x_{i-1})+f(x_i))\ell_i\)。对本节的光滑旋转曲面，由这些圆台侧面积和在分割加细时的极限确定面积。

\begin{figure}[!ht]
\centering
\includegraphics[width=.48\linewidth]{calculus-surface-frustum}
\caption{短弦绕轴生成圆台侧面，其面积为两端圆周长的平均值乘短弦长度；加细后得到\(2\pi r\,ds\)的累积。}
\label{fig:calculus-surface-frustum}
\end{figure}

对\(f\)使用中值定理，可写\(\ell_i=\sqrt{1+f'(\eta_i)^2}\Delta x_i\)。连续性使两端半径的平均与\(f(\eta_i)\)的差一致趋零。又因\(f'\)有界，设\(|f'|\le K\)，则\(\sum_i\ell_i\le\sqrt{1+K^2}(b-a)\)，所以替换半径平均造成的总误差也趋零。圆台和因此与下列Riemann和有相同极限，得到
\begin{equation}\label{eq:calculus-surface-x}
 S=2\pi\int_a^bf(x)\sqrt{1+f'(x)^2}\,dx.
\end{equation}
其中\(2\pi f(x)\)是旋转一周的周长，\(\sqrt{1+f'(x)^2}\,dx\)是母线的长度元。直接写成\(2\pi\int f(x)\,dx\)，只计入横向宽度，会漏掉母线倾斜造成的长度。

令\(f\equiv R\)，公式得到圆柱侧面积\(2\pi R(b-a)\)。若母线从半径\(r_1\)线性变化到\(r_2\)，轴向长度为\(h\)，则斜率恒定，公式给出
\begin{equation*}
 S=\pi(r_1+r_2)\sqrt{h^2+(r_2-r_1)^2},
\end{equation*}
与圆台的几何公式一致。

\paragraph{参数表示与旋转轴} \mbox{}

若母线用\(x=x(t)\)、\(y=y(t)\)表示，坐标连续可导，且\(y\ge0\)，绕横轴旋转时同样以短弦构成圆台。弧长证明中的一致误差控制，加上半径函数的连续性，得到
\begin{equation}\label{eq:calculus-parametric-surface}
 S_x=2\pi\int_\alpha^\beta y(t)\sqrt{x'(t)^2+y'(t)^2}\,dt.
\end{equation}
母线可以先分成有限段沿旋转轴坐标单调的弧段，再把各段面积相加，参数不必在整段上都能反解成一个横坐标函数。

绕纵轴旋转、取\(x\ge0\)的母线时，旋转半径改为\(x(t)\)，所以
\begin{equation*}
 S_y=2\pi\int_\alpha^\beta x(t)\sqrt{x'(t)^2+y'(t)^2}\,dt.
\end{equation*}
更一般地，若母线上各点到所选旋转轴的距离为\(\rho(t)\)，面积统一由\(2\pi\rho\,ds\)累积。改变旋转轴只改变半径，长度元仍由同一条母线决定。

这里计算母线扫出的曲面。若求封闭几何体的总表面积，还须加上未由母线生成的底面或端面；若重复扫过同一片曲面，应先选取只覆盖一次的范围。

\begin{example}
将上半圆\(x=R\cos t\)、\(y=R\sin t\)、\(0\le t\le\pi\)绕横轴旋转，求球面积。再求其中两平行平面\(x=c_1\)、\(x=c_2\)之间的球带面积，观察它与球带轴向高度的关系。
\end{example}

\begin{example}
将\(y=\sqrt{x}\)、\(1\le x\le4\)绕横轴旋转，求生成的曲面面积。再写出该曲线绕直线\(y=-1\)旋转时的面积积分，说明哪些因子改变、哪些因子保持不变。
\end{example}

\begin{exercise}
\mbox{}\par\nobreak

（1）用积分求底面半径为\(R\)、高为\(h\)的正圆锥侧面积，再求总表面积。比较体积与侧面积中局部贡献的表达式。

（2）将\(y=x\)、\(1\le x\le2\)分别绕横轴、纵轴旋转，求曲面面积，用圆台公式核对。

（3）设\(R>r>0\)。圆\(x=R+r\cos t\)、\(y=r\sin t\)、\(0\le t\le2\pi\)绕纵轴旋转生成圆环面，求其面积。解释为何此处应选整圆作为母线，而求球面时只选一个半圆。

（4）同一条母线改变参数后，利用换元公式证明旋转面积不变。若沿母线走两遍，积分结果如何变化？
\end{exercise}

\FloatBarrier
\newpage
\section{反常积分*}
\label{sec:calculus-improper-integrals}

此前积分的区间有限，被积函数也有界。若范围延伸到无穷远，或者函数在端部无界，就先在正常的有限范围内积分，再研究这些积分的极限。无穷累积是否有限，由这个新的极限过程决定。

\subsection{无穷区间上的积分}

\begin{definition}
设\(f\)在\([a,+\infty)\)的每个有限闭子区间上可积。若有限极限存在，定义
\begin{equation*}
 \int_a^{+\infty}f(x)\,dx=\lim_{R\to+\infty}\int_a^Rf(x)\,dx,
\end{equation*}
并称该{\heiti{}反常积分收敛}；否则称其发散。左侧无穷区间的定义类似。
\end{definition}

两端都无穷时，选取一个有限分点\(c\)，分别定义\(\int_{-\infty}^cf\)与\(\int_c^{+\infty}f\)；只有两者都收敛，才定义其和为\(\int_{-\infty}^{+\infty}f\)。同时对称地取两端极限，是另一种取法，不能替代两侧分别收敛。

例如，\(\int_0^{+\infty}e^{-x}\,dx=\lim_{R\to\infty}(1-e^{-R})=1\)。虽然区间无限长，总累积仍有限。对幂函数，直接计算有限上限积分得到
\begin{equation*}
 \int_1^{+\infty}\frac{dx}{x^p}
 \quad\text{在且仅在 }p>1\text{ 时收敛，此时值为 }\frac1{p-1}.
\end{equation*}
\(p=1\)时有限积分为\(\ln R\)，其余情形由幂函数的端部行为判断。

\begin{example}
讨论\(\int_0^{+\infty}xe^{-x}\,dx\)的敛散性并求值。再考察\(f(x)=x\)：对称截取区间的积分极限，能否说明其在全实轴上的反常积分收敛？
\end{example}

\subsection{无界函数的积分}

若\(f\)在\(a\)右侧无界，但在每个\([a+\varepsilon,b]\)上可积，则定义
\begin{equation*}
 \int_a^bf(x)\,dx=\lim_{\varepsilon\to0^+}\int_{a+\varepsilon}^bf(x)\,dx,
\end{equation*}
其中右端为有限极限时称为收敛。在右端无界的情形类似；区间内部有奇点时，先在奇点两侧拆开，分别要求收敛。

对\(x^{-p}\)，有
\begin{equation*}
 \int_0^1\frac{dx}{x^p}
 \quad\text{在且仅在 }p<1\text{ 时收敛，此时值为 }\frac1{1-p}.
\end{equation*}
这与无穷远的判别方向不同：靠近零点时，指数越大增长越快；趋向无穷远时，指数越大衰减越快。

在避开零点的有限积分中令\(x=1/t\)，再取极限，就能联系这两种情形：
\begin{equation*}
 \int_0^1x^{-p}\,dx=\int_1^{+\infty}t^{p-2}\,dt
\end{equation*}
（两边按相应反常积分理解）。因此零点附近的\(p<1\)，正好对应无穷远处幂指数\(2-p>1\)。端部位置改变时，微分\(dx=-dt/t^2\)也参与改变衰减阶数。

例如，\(1/\sqrt{x}\)虽在零点附近无界，积分仍为\(2\)。另一方面，\(\int_{-1}^1dx/x\)不收敛，即使对称避开零点的两个有限积分之和始终为零。

\begin{figure}[!ht]
\centering
\includegraphics[width=.83\linewidth]{calculus-improper}
\caption{左图先截断无穷区间，再估计橙色尾部；右图先避开零点附近的奇性，再令截断位置趋近零点。}
\label{fig:calculus-improper}
\end{figure}

\begin{exercise}
\mbox{}\par\nobreak

（1）判断\(\int_0^1\ln x\,dx\)、\(\int_1^{+\infty}dx/(1+x^2)\)的敛散性，收敛时求值。

（2）分别讨论\(\int_0^1dx/(x\sqrt{x})\)与\(\int_1^{+\infty}dx/(x\sqrt{x})\)。

（3）说明为什么不能直接把\(\ln|x|\)在\(-1,1\)两端的值相减，就断定\(\int_{-1}^1dx/x=0\)。
\end{exercise}

\subsection{收敛判别}

若\(f\ge0\)，有限上限积分随上限单调不减。因此，反常积分收敛，当且仅当这些有限积分有统一上界。这使第一章的单调有界思想再次成为存在性工具。

\begin{theorem}[比较判别]\label{thm:calculus-improper-comparison}
设\(0\le f\le g\)，各函数在有限闭子区间上可积。若\(\int_a^{+\infty}g\)收敛，则\(\int_a^{+\infty}f\)收敛；若后者发散，则前者也发散。在有限奇点附近，同样可以按大小比较。
\end{theorem}

\begin{proof}
由保序性，\(\int_a^Rf\le\int_a^Rg\)。右侧有统一上界时，左侧也有；第二个结论为第一个结论的逆否命题。奇点附近同理。
\end{proof}

比较只需在足够靠近所研究端部的范围内成立，剩下的正常有限积分不影响敛散性。若正函数满足\(f/g\to c\in(0,+\infty)\)，则最终有\((c/2)g\le f\le2cg\)，所以两积分同敛散。

例如，\(x\ge1\)时\(e^{-x^2}\le e^{-x}\)，故\(\int_1^{+\infty}e^{-x^2}\,dx\)收敛，无须先求积分值。

一个积分可能在多个位置需要判别。先用正常的有限分点隔开，逐段处理；一段收敛不能抵消另一段发散。比较函数应按当前端部的行为选择。

\begin{example}
讨论\(\int_0^{+\infty}dx/[x^\alpha(1+x)]\)的敛散性。分别在零点附近与无穷远处比较幂函数；当\(\alpha=1/2\)时，用\(x=t^2\)求积分值。再讨论\(\int_0^1dx/[x^\alpha(1-x)^\beta]\)的敛散条件。
\end{example}

对可能变号的函数，可以使用Cauchy准则：\(\int_a^{+\infty}f\)收敛，当且仅当对任意\(\varepsilon>0\)，存在\(R\)，使任意\(v>u>R\)都有
\begin{equation}\label{eq:calculus-improper-cauchy}
 \left|\int_u^vf(x)\,dx\right|<\varepsilon.
\end{equation}
必要性由收敛极限直接得到。反过来，取整数上限的积分，由数列Cauchy准则得到极限，再用式\eqref{eq:calculus-improper-cauchy}比较任意足够大的上限与整数上限，便得到相同极限。这与可积性证明中先取一列、再比较任意取样的组织相似。

\subsection{绝对收敛}

\begin{definition}
若\(\int_a^{+\infty}|f(x)|\,dx\)收敛，就称\(\int_a^{+\infty}f(x)\,dx\){\heiti{}绝对收敛}。
\end{definition}

\begin{theorem}
绝对收敛的反常积分收敛。
\end{theorem}

\begin{proof}
由有限区间上的绝对值估计，\(|\int_u^vf|\le\int_u^v|f|\)。右侧由收敛性满足Cauchy条件，因此原积分也满足。
\end{proof}

例如，\(|\sin x|/x^2\le1/x^2\)，故\(\int_1^{+\infty}\sin x/x^2\,dx\)绝对收敛。原积分收敛而绝对值积分发散时，称为{\heiti{}条件收敛}。

一个典型例子是\(\int_1^{+\infty}\sin x/x\,dx\)。对\(v>u\ge1\)，分部积分得到
\begin{equation*}
 \int_u^v\frac{\sin x}{x}\,dx
 =\left.-\frac{\cos x}{x}\right|_u^v-\int_u^v\frac{\cos x}{x^2}\,dx,
 \qquad
 \left|\int_u^v\frac{\sin x}{x}\,dx\right|\le\frac2u.
\end{equation*}
所以原积分收敛。但对每个正整数\(k\)，在\([k\pi+\pi/6,k\pi+5\pi/6]\)上\(|\sin x|\ge1/2\)，故这一段的绝对值积分至少为\(1/[3(k+1)]\)。将这些互不重叠的区间相加，与调和级数比较，便知绝对值积分发散。

\begin{exercise}
\mbox{}\par\nobreak

（1）判断\(\int_1^{+\infty}dx/(x^2+\sqrt{x})\)及\(\int_0^1dx/\sqrt{x+x^2}\)的敛散性，指出比较对象。

（2）证明\(\int_0^{+\infty}e^{-x}|\cos x|\,dx\)收敛，进而说明\(\int_0^{+\infty}e^{-x}\cos x\,dx\)收敛。

（3）将上述分部积分用于\(\int_1^{+\infty}\cos x/x\,dx\)，判断它是否收敛、是否绝对收敛。
\end{exercise}

\subsection{尾部估计}

积分收敛后，还可以研究只计算有限范围造成的误差。若\(|f(x)|\le g(x)\)且\(g\)非负可积，则
\begin{equation*}
 \left|\int_a^{+\infty}f-\int_a^Rf\right|
 \le\int_R^{+\infty}g(x)\,dx.
\end{equation*}
例如，\(p>1\)时\(\int_R^{+\infty}x^{-p}\,dx=R^{1-p}/(p-1)\)。对\(R\ge1\)，又有
\begin{equation*}
 \int_R^{+\infty}e^{-x^2}\,dx
 \le\frac1R\int_R^{+\infty}xe^{-x^2}\,dx
 =\frac{e^{-R^2}}{2R}.
\end{equation*}
这既证明尾部很小，也给出选择截断位置的依据。

无穷区间上不能直接把最大误差乘以总长度。比如
\begin{equation*}
 f_n(x)=\begin{cases}\dfrac1n\left(1-\dfrac xn\right),&0\le x\le n,\\0,&x>n\end{cases}
\end{equation*}
在整个半轴上与零函数的误差不超过\(1/n\)，但积分恒为\(1/2\)。误差虽小，发生误差的范围却在扩大。

因此，比较无穷累积时，可以先固定\(R\)，控制\([a,R]\)上的误差，再分别估计两函数在\([R,+\infty)\)上的尾部。有限范围与尾部承担不同的误差控制任务。

例如，要近似\(\int_0^{+\infty}e^{-x^2}\,dx\)，先选\(R\)，使\(e^{-R^2}/(2R)\)小于给定总误差的一半；再对固定的\([0,R]\)使用中点公式，使有限区间求积误差小于另一半。先截断范围、后选分段数，两种误差便可分别验证，而不要求预先知道整个反常积分的精确值。

\subsection{积分判别法}

设\(f\)在\([1,+\infty)\)上连续、非负且单调不增。每个单位区间上的矩形与曲边面积满足
\begin{equation*}
 f(n+1)\le\int_n^{n+1}f(x)\,dx\le f(n).
\end{equation*}
累加得到
\begin{equation*}
 \int_1^{N+1}f(x)\,dx\le\sum_{n=1}^Nf(n)
 \le f(1)+\int_1^Nf(x)\,dx.
\end{equation*}
由正项级数与非负反常积分的单调有界判据，\(\sum_{n=1}^{\infty}f(n)\)与\(\int_1^{+\infty}f(x)\,dx\)同敛散。这称为{\heiti{}积分判别法}。

收敛时，还得到尾和的两侧估计：
\begin{equation*}
 \int_{N+1}^{+\infty}f(x)\,dx
 \le\sum_{n=N+1}^{\infty}f(n)
 \le\int_N^{+\infty}f(x)\,dx.
\end{equation*}
例如，\(p>1\)时，\(\sum_{n>N}n^{-p}\)介于\((N+1)^{1-p}/(p-1)\)与\(N^{1-p}/(p-1)\)之间。这接回了第一章的\(p\)级数，并进一步给出截断误差。

令\(D_N=\sum_{n=1}^Nf(n)-\int_1^Nf(x)\,dx\)。前面的估计给出\(0\le D_N\le f(1)\)，且
\begin{equation*}
 D_{N+1}-D_N=f(N+1)-\int_N^{N+1}f(x)\,dx\le0,
\end{equation*}
故\(D_N\)单调有界，必收敛。取\(f(x)=1/x\)，即得\(H_N-\ln N\to\gamma\)；调和级数与积分都发散，差值仍有极限，误差见式\eqref{eq:calculus-euler-constant-error}。

\begin{exercise}
\mbox{}\par\nobreak

（1）估计\(\sum_{n=N+1}^{\infty}1/n^2\)。要使尾和小于\(10^{-3}\)，可以选取怎样的\(N\)？

（2）讨论\(\sum_{n=2}^{\infty}1/[n(\ln n)^p]\)的敛散性，并说明被积函数在哪一段起单调不增。

（3）给出\(\int_1^{+\infty}\sin x/x\,dx\)截断到\(R\)的误差界。与绝对收敛情形的尾部比较方法有何不同？

（4）若要用数值积分计算\(\int_0^{+\infty}e^{-x^2}\,dx\)，设计一种将总误差分成有限区间求积误差与尾部误差的方案。

（5）计算\(\int_e^{+\infty}dx/[x(\ln x)^2]\)，并给出级数\(\sum_{n=2}^{\infty}1/[n(\ln n)^2]\)截断到\(n=N\)后的尾和上下界。要保证尾和小于\(1/10\)，可以取怎样的\(N\)？比较它与\(\sum1/n^2\)的截断速度。

（6）讨论\(\int_0^1x^{-p}\ln x\,dx\)的敛散性，收敛时求其值。若从\(\varepsilon\)开始积分，求遗漏\((0,\varepsilon)\)所造成的误差，并说明它怎样随\(\varepsilon\to0^+\)趋零。

（7）设\(p>0\)。用分部积分证明\(\int_1^{+\infty}\sin x/x^p\,dx\)收敛，并给出截断到\(R\)的误差界。再按\(p\)的范围判断绝对收敛与条件收敛；讨论绝对值时，可以在每个周期中选取\(|\sin x|\ge1/2\)的子区间。

（8）设\(f_n(x)=n^2xe^{-nx}\)，\(x\ge0\)。求逐点极限及每个\(f_n\)在半轴上的积分。说明为何不能直接交换积分与极限，并比较它与正文中图像向远处铺开的例子：本题的面积集中到了哪里？
\end{exercise}

\FloatBarrier
\newpage
\section{微分方程}
\label{sec:calculus-differential-equations}

前面给定一个函数，研究它怎样变化；现在反过来，给定变化规律，寻找符合规律的函数。例如，已知一个种群的增长速度与当前数量成正比，能否确定数量随时间的变化？变化率与未知函数同时出现，就形成了新的方程。

\subsection{微分方程}

设\(y(t)\)表示时刻\(t\)的种群数量。若单位数量在单位时间内的增长率为常数\(k\)，则短时间内的增量近似为\(ky(t)\Delta t\)。将这一规律写成瞬时变化率，就是
\begin{equation*}
 y'(t)=ky(t).
\end{equation*}
方程中的未知对象是函数，等式同时约束它的取值与变化率。

\begin{definition}
含有未知函数及其导数的方程，称为{\heiti{}微分方程}。未知函数只有一个自变量时，称为{\heiti{}常微分方程}；方程中出现的最高阶导数的阶数，称为方程的{\heiti{}阶数}。
\end{definition}

先研究可以写成
\begin{equation}\label{eq:calculus-ode-normal-form}
 y'=F(t,y)
\end{equation}
的一阶方程，其中\(F\)是给定的连续函数。若可导函数\(y=y(t)\)在区间\(I\)上满足此等式，就称它为方程在\(I\)上的一个{\heiti{}解}。例如，\(y=Ce^{kt}\)对任意常数\(C\)都满足\(y'=ky\)，所以同一方程可以有一族解。

实际问题通常还给出某一时刻的状态。方程连同条件\(y(t_0)=y_0\)，组成一个{\heiti{}初值问题}；\((t_0,y_0)\)称为初值。对于增长方程，初值确定\(C\)，得到
\begin{equation*}
 y(t)=y_0e^{k(t-t_0)}.
\end{equation*}
这确实给出了全部解：若\(y'=ky\)，则\((e^{-kt}y)'=0\)，由中值定理，\(e^{-kt}y\)在区间上为常数。

还可以从图像理解式\eqref{eq:calculus-ode-normal-form}。若解的图像经过\((t,y)\)，该处切线斜率就必须是\(F(t,y)\)。因此即使暂时求不出表达式，也可以由\(F\)的符号判断解在哪里上升、在哪里下降。若常数\(c\)满足\(F(t,c)=0\)，则\(y(t)=c\)是一个常值解，也称为{\heiti{}平衡解}。

\begin{example}
验证\(y=t-1+Ce^{-t}\)满足\(y'=t-y\)，再求其中满足\(y(0)=2\)的解。根据方程判断解的图像在直线\(y=t\)的上、下方分别怎样变化。
\end{example}

\begin{exercise}
\mbox{}\par\nobreak

（1）验证\(y=\tan(t+C)\)满足\(y'=1+y^2\)。求其中满足\(y(0)=0\)的解，并写出包含零点、使此解有定义的最大区间。

（2）某种群满足\(P'=kP\)、\(P(0)=P_0>0\)，经过\(3\)个时间单位数量翻倍。求\(k\)，并求数量达到\(3P_0\)所需的时间。

（3）对方程\(y'=y(2-y)\)，求平衡解，并按\(y\)的取值判断解的增减。若只知道\(y(0)=1\)，能否由当前斜率直接断定\(y(1)=2\)？
\end{exercise}

\subsection{可分离变量的方程}

设\(a,b\)连续。若方程具有形式
\begin{equation*}
 y'=a(t)b(y),
\end{equation*}
就可以尝试将两种变量分别放入一个原函数。在\(b(y)\ne0\)的范围内，取\(H'(y)=1/b(y)\)、\(A'(t)=a(t)\)。链式法则给出
\begin{equation*}
 \frac{d}{dt}H(y(t))=\frac{y'(t)}{b(y(t))}=a(t),
 \qquad H(y(t))=A(t)+C.
\end{equation*}
这样就把微分方程转化为\(y\)与\(t\)的关系。计算时常简记为
\begin{equation}\label{eq:calculus-ode-separation}
 \frac{dy}{b(y)}=a(t)\,dt,
 \qquad \int\frac{dy}{b(y)}=\int a(t)\,dt.
\end{equation}
这称为{\heiti{}分离变量法}，其依据正是上面的复合函数求导。若\(b(c)=0\)，还应保留在相除前已经得到的平衡解\(y=c\)。

例如，\(y'=ky\)在\(y\ne0\)时给出\(\ln|y|=kt+C\)，从而\(y=Ce^{kt}\)、\(C\ne0\)；补入平衡解\(y=0\)，就得到前面研究的整个解族。

\paragraph{有限容量下的增长} \mbox{}

当资源有限时，可以让单位数量的增长率随数量增加而降低。设\(K>0\)表示环境容量，取模型
\begin{equation}\label{eq:calculus-logistic-equation}
 y'=ky\left(1-\frac yK\right),\qquad k>0.
\end{equation}
当\(0<y<K\)时增长，接近\(K\)时增长率趋小；\(y=0\)与\(y=K\)是平衡解。

设初值为\(y(0)=y_0\)、\(0<y_0<K\)。在\(0<y<K\)内分离变量，利用
\begin{equation*}
 \frac1{y(1-y/K)}=\frac1y+\frac1{K-y},
\end{equation*}
得到\(\ln\frac{y}{K-y}=kt+C\)。代入初值并整理，便得
\begin{equation}\label{eq:calculus-logistic-solution}
 y(t)=\frac{K}{1+(K/y_0-1)e^{-kt}}.
\end{equation}
这个函数始终位于\(0\)与\(K\)之间，代回可验证它满足原方程。对\(t\ge0\)，它单调增加并趋于\(K\)。与指数增长相比，变化规律中的容量限制决定了长期行为。

\begin{example}
环境温度恒为\(20\)，物体温度\(T\)满足\(T'=-k(T-20)\)、\(k>0\)。已知\(T(0)=80\)、\(T(10)=50\)，求温度函数及降到\(35\)所需的时间，并解释它为什么持续下降却不会在有限时间内达到环境温度。
\end{example}

\begin{exercise}
\mbox{}\par\nobreak

（1）求初值问题\(y'=2t(1+y^2)\)、\(y(0)=0\)的解，并确定包含零点的最大解区间。

（2）求\(y'=y^2\)、\(y(0)=1\)的解及包含零点的最大解区间。再求平衡解，比较它们随时间的变化。

（3）设\(y\)由式\eqref{eq:calculus-logistic-solution}给出。求\(y''\)，判断增长速度在哪个数量水平最大，并讨论图像的凸性。

（4）对\(y'=ky(1-y/K)\)，设\(K/2<y_0<K\)。求达到\((y_0+K)/2\)所需的时间，再与从\(y_0\)增长到\(2y_0\)的要求比较，说明容量限制的作用。
\end{exercise}

\subsection{一阶线性微分方程}

若未知函数及其一阶导数以一次形式出现，整理为
\begin{equation}\label{eq:calculus-linear-ode}
 y'+p(t)y=q(t),
\end{equation}
就得到{\heiti{}一阶线性微分方程}。以下设\(p,q\)在所讨论的区间连续。

分离变量适合将变量拆开；这里可以利用另一条已有结构：乘积求导。若能找到函数\(\mu\)，使
\begin{equation*}
 (\mu y)'=\mu y'+\mu' y=\mu(y'+py),
\end{equation*}
原方程就可以化为一个乘积的导数。比较系数，只须令\(\mu'=p\mu\)。固定区间内的\(t_0\)，取
\begin{equation*}
 P(t)=\int_{t_0}^t p(s)\,ds,\qquad \mu(t)=e^{P(t)},
\end{equation*}
便有\(\mu'=p\mu\)、\(\mu(t_0)=1\)。\(\mu\)称为{\heiti{}积分因子}。

两边乘以\(\mu\)，方程化为\((\mu y)'=\mu q\)。从\(t_0\)积分到\(t\)，再用初值\(y(t_0)=y_0\)，得到
\begin{equation}\label{eq:calculus-linear-ode-solution}
 y(t)=e^{-P(t)}\left(y_0+\int_{t_0}^t e^{P(s)}q(s)\,ds\right).
\end{equation}
右端是已知函数，求导即可验证它满足方程与初值；反过来，任何解都必须满足刚才的积分等式。因此这个初值问题在该区间上恰有一个解。

例如，\(y'+ky=r\)、\(k>0\)，其中\(r\)为常数，取\(t_0=0\)，公式给出
\begin{equation*}
 y(t)=\frac rk+\left(y_0-\frac rk\right)e^{-kt}.
\end{equation*}
平衡值\(r/k\)由持续补充与按比例减少相抵消得到，初值决定偏离平衡值的幅度。令\(z=y-r/k\)，又回到\(z'=-kz\)。

将右端置零所得的方程称为{\heiti{}齐次方程}。若\(y_1,y_2\)满足原方程，则差\(z=y_1-y_2\)满足\(z'+pz=0\)，从而\(z=Ce^{-P}\)。因此其他解都由一个给定解加上齐次解得到；相同初值使差中的常数为零。

\begin{example}
容器内原有\(100\)升清水，每分钟流入\(2\)升浓度为\(3\)克/升的盐水，同时流出\(2\)升充分混合的溶液。设\(m(t)\)为\(t\)分钟后的盐量，建立初值问题，求\(m(t)\)，并判断盐量是否能在有限时间内达到\(300\)克。
\end{example}

\begin{exercise}
\mbox{}\par\nobreak

（1）求\(y'+2y=e^t\)、\(y(0)=0\)的解。

（2）求\(ty'+y=t^2\)、\(y(1)=0\)在\(t>0\)上的解，比较直接辨认乘积导数与使用积分因子的计算。

（3）求\(y'-2ty=2t\)、\(y(0)=1\)的解，并研究它的最小值及图像的对称性。

（4）设\(p,q\)在\([0,T]\)上连续，\(q\ge0\)，且\(y(0)=y_0\)。利用式\eqref{eq:calculus-linear-ode-solution}证明，\(y_0\ge0\)时，解在整个区间上非负；若\(y_0>0\)，则始终为正。

（5）设\(f\)连续且满足\(f(x)=1+\int_0^x(2t-f(t))\,dt\)。求\(f\)、它在\([0,2]\)上的最小值及\(\int_0^2 f(x)\,dx\)。从原式说明\(f\)为什么可导。
\end{exercise}

\begingroup
\subsection{二阶常系数线性微分方程}
\label{subsec:calculus-second-order-ode}
\setlength{\parskip}{4pt}
\setlength{\abovedisplayskip}{10pt}
\setlength{\belowdisplayskip}{10pt}

速度的变化率是加速度。若物体偏离平衡位置的位移为\(y\)，质量为\(m\)，回复力为\(-ky\)，则运动规律给出\(my''=-ky\)。这里除了初始位置，还要给出初始速度，才能描述物体从怎样的状态开始运动。

将这类关系推广，考虑
\begin{equation}\label{eq:calculus-second-order-ode}
 y''+ay'+by=f(t),
\end{equation}
其中\(a,b\)为实常数，\(f\)连续。这称为{\heiti{}二阶常系数线性微分方程}。它的解是满足等式的二阶连续可导函数；相应的初始条件为\(y(t_0)=y_0\)、\(y'(t_0)=v_0\)。先研究\(f=0\)的{\heiti{}齐次方程}。

\paragraph{特征方程} \mbox{}

指数函数求导后只改变常数因子，因此尝试\(y=e^{rt}\)。代入齐次方程，得到
\begin{equation}\label{eq:calculus-ode-characteristic}
 r^2+ar+b=0.
\end{equation}
它称为原方程的{\heiti{}特征方程}。根的类型将决定解的形式。记\(\Delta=a^2-4b\)，令\(y=e^{-at/2}u\)，直接求导可将齐次方程化为
\begin{equation*}
 u''-\frac\Delta4u=0.
\end{equation*}
下面分别求出三种情形的全部解，即{\heiti{}通解}。

1°\(\Delta>0\)。令\(\beta=\sqrt\Delta/2\)，则\(u''-\beta^2u=0\)。置\(z=u'-\beta u\)，有\(z'=-\beta z\)。先求\(z\)，再解\(u'-\beta u=z\)，得到\(u=C_1e^{\beta t}+C_2e^{-\beta t}\)。所以，若特征方程有不同实根\(r_1,r_2\)，通解为
\begin{equation*}
 y=C_1e^{r_1t}+C_2e^{r_2t}.
\end{equation*}

2°\(\Delta=0\)。此时\(u''=0\)，故\(u=C_1+C_2t\)。特征方程有重根\(r=-a/2\)，通解为
\begin{equation*}
 y=(C_1+C_2t)e^{rt}.
\end{equation*}
重根时，\(te^{rt}\)提供了另一个解。

3°\(\Delta<0\)。记\(\alpha=-a/2\)、\(\beta=\sqrt{-\Delta}/2\)，则\(u''+\beta^2u=0\)。正弦、余弦满足此式，给出
\begin{equation*}
 y=e^{\alpha t}(C_1\cos\beta t+C_2\sin\beta t).
\end{equation*}
这已经包含全部解。事实上，任给\(u''+\beta^2u=0\)的解，选取正弦、余弦的组合与它具有相同的初始值和初始导数。两者之差\(w\)满足\(w''+\beta^2w=0\)，且初始值与初始导数都为零。由于
\begin{equation*}
 \frac{d}{dt}\bigl((w')^2+\beta^2w^2\bigr)=2w'(w''+\beta^2w)=0,
\end{equation*}
括号内恒为零，故\(w=0\)。特征方程的共轭复根\(\alpha\pm i\beta\)中，实部决定指数因子，虚部的绝对值决定振荡部分的周期\(2\pi/\beta\)。

上述三种通解各含两个任意常数。代入\(y(t_0)=y_0\)、\(y'(t_0)=v_0\)，所得两个线性方程都能唯一确定这两个常数，因而齐次初值问题有唯一解。

例如，\(y''+\omega^2y=0\)、\(\omega>0\)，在\(y(0)=y_0\)、\(y'(0)=v_0\)下的解为
\begin{equation*}
 y(t)=y_0\cos\omega t+\frac{v_0}{\omega}\sin\omega t.
\end{equation*}
由三角函数的合成可知它作周期运动；初始位置与速度共同决定振幅和相位。加入\(ay'\)、\(a>0\)后，若仍有共轭复根，则指数因子\(e^{-at/2}\)使振幅随时间衰减。

\begin{example}
分别求下列初值问题的解，比较特征根与图像变化的关系：
\begin{equation*}
 \begin{aligned}
 &y''-3y'+2y=0,&&y(0)=1,\quad y'(0)=0;\\
 &y''+2y'+y=0,&&y(0)=0,\quad y'(0)=1;\\
 &y''+2y'+5y=0,&&y(0)=1,\quad y'(0)=-1.
 \end{aligned}
\end{equation*}
\end{example}

\paragraph{非齐次方程} \mbox{}

若已经找到式\eqref{eq:calculus-second-order-ode}的一个解\(y_p\)，则任意其他解与\(y_p\)之差都满足齐次方程。反过来，将\(y_p\)与任意齐次解相加，仍满足原方程。因此非齐次方程的通解为
\begin{equation*}
 y=y_p+y_h,
\end{equation*}
其中\(y_h\)为齐次通解，\(y_p\)称为一个{\heiti{}特解}。加上初始条件后，两个常数仍由函数值与导数唯一确定。

寻找特解时，可以利用求导后表达式的结构。例如，右端为\(Ae^{\lambda t}\)时，先尝试\(y_p=Be^{\lambda t}\)。若\(\lambda^2+a\lambda+b\ne0\)，代入便得\(B=A/(\lambda^2+a\lambda+b)\)。右端为多项式或正弦、余弦时，也可以先设同类表达式，再代入比较系数。

若试设的表达式已经是齐次解，代入后左端恒为零，就需要改变试设。例如，对
\begin{equation*}
 y''-3y'+2y=e^t,
\end{equation*}
\(Be^t\)不能产生右端，但代入\(Bte^t\)后左端为\(-Be^t\)，从而可取\(y_p=-te^t\)。特解提供方程要求的变化，齐次解则保留调节初始状态的自由。

\begin{example}
求\(y''+y=\sin t\)、\(y(0)=y'(0)=0\)的解。先尝试\(y_p=A\cos t+B\sin t\)，再尝试\(y_p=Ct\cos t\)，比较两次代入的结果，并判断解在\(t\ge0\)时是否有界。
\end{example}

\begin{exercise}
\mbox{}\par\nobreak

（1）分别求\(y''-4y=0\)、\(y''-4y'+4y=0\)、\(y''+4y=0\)的通解，说明它们的特征根怎样决定解的形式。

（2）求\(y''+4y=0\)、\(y(0)=1\)、\(y'(0)=2\)的解、振幅及最小正周期。证明\((y')^2+4y^2\)为常数，并由它核对振幅。

（3）求\(y''+3y'+2y=2\)、\(y(0)=y'(0)=0\)的解，研究它在\(t\ge0\)上的增减与极限。

（4）分别为\(y''-y=t\)与\(y''-y=e^t\)寻找特解，再写出通解。比较两题中选择试设表达式的依据。

（5）求\(y''-2y'+y=e^t\)的通解。比较试设\(Be^t\)、\(Bte^t\)、\(Bt^2e^t\)时的结果，说明重根怎样影响特解的选择。
\end{exercise}

\subsection{Euler方法}
\setlength{\abovedisplayskip}{10pt}
\setlength{\belowdisplayskip}{10pt}

解的表达式不易求出时，还可以从初值逐步计算。设\(y'=F(t,y)\)、\(y(t_0)=y_0\)。在\(t_0\)附近，局部一次近似为
\begin{equation*}
 y(t_0+h)\approx y_0+hF(t_0,y_0).
\end{equation*}
将这个近似值作为下一步的起点，再按该处的斜率前进，就得到递推
\begin{equation}\label{eq:calculus-euler-method}
 t_n=t_0+nh,\qquad y_{n+1}=y_n+hF(t_n,y_n),\qquad h>0.
\end{equation}
它称为{\heiti{}Euler方法}，\(h\)称为步长。将相邻点\((t_n,y_n)\)用线段相连，便得到近似解的折线。每段斜率由起点处的方程给出，计算本身只需反复代入已知函数\(F\)。

对二阶方程\(y''+ay'+by=f(t)\)，引入速度\(v=y'\)，就得到\(y'=v\)、\(v'=f(t)-av-by\)。以\(Y_n,V_n\)近似位置与速度，分别使用一次近似，得到同时更新的递推
\begin{equation*}
 Y_{n+1}=Y_n+hV_n,\qquad
 V_{n+1}=V_n+h\bigl(f(t_n)-aV_n-bY_n\bigr).
\end{equation*}
两个新值都由第\(n\)步的状态计算，初始位置与初始速度共同确定整个递推。

例如，对\(y'=y\)、\(y(0)=1\)，有\(y_{n+1}=(1+h)y_n\)，所以\(y_n=(1+h)^n\)。若把\([0,T]\)分成\(N\)段，终点近似为
\begin{equation*}
 y_N=\left(1+\frac TN\right)^N.
\end{equation*}
当\(T=1\)时，这正是第一章引入\(e\)的数列。这里它来自用递推近似连续增长：每一步按当前数量增长，再以更新后的数量继续计算。

\begin{figure}[!ht]
\centering
\includegraphics[width=.76\linewidth]{calculus-euler-method}
\caption{对\(y'=y\)、\(y(0)=1\)作Euler折线。蓝线为精确解\(e^t\)，橙线分别取步长\(1/2\)与\(1/4\)。}
\label{fig:calculus-euler-method}
\end{figure}

这个例子的误差可以直接估计。由
\begin{equation*}
 0\le u-\ln(1+u)=\int_0^u\frac{s}{1+s}\,ds\le\frac{u^2}2
 \qquad(u\ge0),
\end{equation*}
得到\(0\le T-N\ln(1+T/N)\le T^2/(2N)\)。再用\(1-e^{-v}\le v\)、\(v\ge0\)，可得
\begin{equation*}
 0\le e^T-\left(1+\frac TN\right)^N\le\frac{e^TT^2}{2N}.
\end{equation*}
因此固定终点、不断缩短步长时，近似值趋于精确值。对一般方程，也需要同时控制每步产生的误差和已有误差在后续计算中的传播。

\begin{example}
\label{ex:calculus-euler-values}
用Euler方法近似\(y'=t-y\)、\(y(0)=1\)在\(t=1\)处的解，分别取\(h=1/2\)与\(h=1/4\)。再求精确解，比较误差，并画出两条近似折线。
\end{example}

\begin{exercise}
\mbox{}\par\nobreak

（1）对\(y'=y\)、\(y(0)=1\)，计算\(h=1,1/2,1/4\)时在\(t=1\)处的近似值。利用\(e^t\)的凸性解释折线为什么位于精确解下方。

（2）设\(y'=-ky\)、\(y(0)=1\)、\(k>0\)。写出Euler递推的通项，求使所有近似值非负且单调不增的步长范围。再比较\(h=1/(2k)\)与\(h=3/(2k)\)时的变化。

（3）对\(y'=1+t^2\)、\(y(0)=0\)，把\([0,1]\)等分为\(N\)段。写出终点近似值，说明它与左端点矩形求积的关系，并估计误差。

（4）对\(y''+y=0\)、\(y(0)=1\)、\(y'(0)=0\)使用上述递推。证明\(Y_n^2+V_n^2=(1+h^2)^n\)，并与精确解的\(y^2+(y')^2=1\)比较。
\end{exercise}
\endgroup

\FloatBarrier
\newpage
\begingroup
\setlength{\parskip}{4pt}
\setlength{\abovedisplayskip}{10pt}
\setlength{\belowdisplayskip}{10pt}
\subsection{线性微分算子*}
\label{subsec:calculus-linear-differential-operators}

\paragraph{一阶方程的算子表示} \mbox{}

第\ref{sec:calculus-function-spaces}节把求导与积分看作作用在函数空间上的运算。现在取连续函数\(p\)，定义乘法算子\(M_p:C[a,b]\to C[a,b]\)、\((M_pu)(t)=p(t)u(t)\)。将\(D\)与限制在\(C^1[a,b]\)上的\(M_p\)相加，就得到线性算子
\begin{equation*}
 L=D+M_p,\qquad Lu=u'+pu.
\end{equation*}
一阶线性方程就是\(Ly=q\)，相应的齐次方程\(Ly=0\)研究\(L\)的核。由前面的解法，固定\(t_0\in[a,b]\)，记\(P(t)=\int_{t_0}^t p(s)\,ds\)，则
\begin{equation*}
 \ker L=\{Ce^{-P}:C\in\mathbb R\}.
\end{equation*}
因此只要找到一个解\(y_*\)，全部解便构成\(y_*+\ker L\)。初值\(y(t_0)=y_0\)从中确定唯一的一个函数。

积分因子也有相应的运算解释。乘积法则表明
\begin{equation*}
 Lu=e^{-P}D(e^Pu).
\end{equation*}
先乘以\(e^P\)，再求导，最后乘以\(e^{-P}\)，就得到\(L\)。求解\(Ly=q\)因而可以依次化为乘法、求原函数和确定初值。若限制输入为
\begin{equation*}
 C^1_{t_0}=\{u\in C^1[a,b]:u(t_0)=0\},
\end{equation*}
则\(L:C^1_{t_0}\to C[a,b]\)为双射，其逆由式\eqref{eq:calculus-linear-ode-solution}取\(y_0=0\)给出。第5.11节对求导算子的讨论，正对应这里\(p=0\)的情形。

\paragraph{二阶方程的算子分解} \mbox{}

记\(C^2[a,b]\)为具有二阶连续导数的函数组成的线性空间。对实常数\(a_1,a_0\)，定义
\begin{equation*}
 L_2:C^2[a,b]\to C[a,b],\qquad L_2u=u''+a_1u'+a_0u.
\end{equation*}
以\(\mathrm{Id}\)表示恒等运算，将各项的结果都看作连续函数，就可以记为\(L_2=D^2+a_1D+a_0\mathrm{Id}\)。若特征方程有实根\(r_1,r_2\)，则
\begin{equation*}
 L_2=(D-r_1\mathrm{Id})(D-r_2\mathrm{Id}).
\end{equation*}
右侧先将\(C^2\)中的函数送入\(C^1\)，再送入\(C\)。这正对应先令\(z=y'-r_2y\)，再依次求解两个一阶方程。重根时两个因子相同，第二次求解产生\(te^{rt}\)。

若根为\(\alpha\pm i\beta\)、\(\beta>0\)，则在实函数空间上仍可写成
\begin{equation*}
 L_2=(D-\alpha\mathrm{Id})^2+\beta^2\mathrm{Id},\qquad
 L_2(e^{\alpha t}u)=e^{\alpha t}(u''+\beta^2u).
\end{equation*}
乘以指数函数将方程化为振荡方程，这就是前面实变量换元的算子解释。常系数使这些一阶因子的顺序可以交换；若把常数\(r\)换成可导函数\(p(t)\)，则乘积求导会多出\(p'u\)，分解必须据实际运算重新核对。

第\ref{subsec:calculus-second-order-ode}小节的三种通解，分别给出\(\ker L_2\)的一组基：
\begin{equation*}
 \begin{aligned}
 &e^{r_1t},\ e^{r_2t}&&\text{（不同实根）},\\
 &e^{rt},\ te^{rt}&&\text{（重根）},\\
 &e^{\alpha t}\cos\beta t,\ e^{\alpha t}\sin\beta t
 &&\text{（共轭复根）}.
 \end{aligned}
\end{equation*}
三种情形中的核都是二维实线性空间；若有一个特解\(y_p\)，方程\(L_2y=f\)的全部解就是\(y_p+\ker L_2\)。取初值的映射
\begin{equation*}
 \ker L_2\longrightarrow\mathbb R^2,\qquad
 y\longmapsto\bigl(y(t_0),y'(t_0)\bigr)
\end{equation*}
是线性双射。两个初始数值因此与齐次解中的两个自由参数一一对应。

\paragraph{零初值下的逆运算} \mbox{}

\begingroup
\setlength{\abovedisplayskip}{7pt}
\setlength{\belowdisplayskip}{7pt}
\setlength{\jot}{1pt}
为使累积量一次求导无端点项、二次求导产生\(f(t)\)，取齐次解\(S\)满足\(S(0)=0\)、\(S'(0)=1\)。由三种根型，
\begin{equation*}
 S(t)=\begin{cases}
 \dfrac{e^{r_1t}-e^{r_2t}}{r_1-r_2},&r_1\ne r_2\text{为实根},\\[3pt]
 te^{rt},&r\text{为重根},\\[2pt]
 e^{\alpha t}\dfrac{\sin\beta t}{\beta},&\alpha\pm i\beta\text{为共轭复根}.
 \end{cases}
\end{equation*}
对连续函数\(f\)，按时间差累积：
\begin{equation*}
 (Gf)(t)=\int_{t_0}^t S(t-s)f(s)\,ds.
\end{equation*}
展开\(S(t-s)\)，由基本定理与乘积法则，得
\begin{equation*}
 \begin{aligned}
 (Gf)'(t)&=\int_{t_0}^t S'(t-s)f(s)\,ds,\\
 (Gf)''(t)&=f(t)+\int_{t_0}^t S''(t-s)f(s)\,ds.
 \end{aligned}
\end{equation*}
\endgroup
因此\(L_2Gf=f\)，且\((Gf)(t_0)=(Gf)'(t_0)=0\)。对振荡方程，\(S(t)=\sin(\omega t)/\omega\)，可用正弦的差角公式展开。

于是，限制在
\begin{equation*}
 C^2_{t_0,0}=\{u\in C^2[a,b]:u(t_0)=u'(t_0)=0\}
\end{equation*}
之后，\(L_2:C^2_{t_0,0}\to C[a,b]\)与\(G\)互为逆映射。若初值为一般的\(y_0,v_0\)，再加上具有这两个初值的齐次解即可。这也证明连续右端下二阶常系数初值问题的存在与唯一。

\begin{exercise}
\mbox{}\par\nobreak

（1）初始时刻取\(0\)。对\(L_2=(D-\mathrm{Id})^2\)，写出零初值下的逆运算\(G\)。利用分部积分证明\((Gf)(t)=e^tI^2(e^{-t}f(t))\)，其中\((Ig)(t)=\int_0^t g(s)\,ds\)。再取\(f(t)=e^t\)，核对重根情形的特解。

（2）设\(p,q\)一阶连续可导。分别计算\((D-M_p)(D-M_q)u\)与\((D-M_q)(D-M_p)u\)，求它们的差，并判断这两个运算何时可以交换。

（3）设\(\omega,\nu>0\)。用\(G\)求\(y''+\omega^2y=\sin\nu t\)、\(y(0)=y'(0)=0\)的解，比较\(\nu\ne\omega\)与\(\nu=\omega\)时解的有界性，解释相同频率的持续输入怎样改变运动。
\end{exercise}

\subsection{逐次逼近*}
\label{subsec:calculus-successive-approximations}
\begingroup
\setlength{\parskip}{3pt}
\setlength{\abovedisplayskip}{8pt}
\setlength{\belowdisplayskip}{8pt}

\paragraph{从微分关系到积分关系} \mbox{}

前面的解法与近似都在研究同一个初值问题。现在把它整体地写成累积关系：若\(y'=F(t,y)\)、\(y(t_0)=y_0\)，由基本定理，
\begin{equation}\label{eq:calculus-ode-integral-equation}
 y(t)=y_0+\int_{t_0}^t F(s,y(s))\,ds.
\end{equation}
反过来，若连续函数\(y\)满足此式，由\(F\)连续，右侧可导，求导并代入\(t=t_0\)便得到原方程与初值。因此这两种表达等价。

这里的被积函数仍然包含未知的\(y\)。可以先给出一个近似函数\(u\)，用它计算变化率，再积分得到新的函数：
\begin{equation*}
 (\mathcal Tu)(t)=y_0+\int_{t_0}^t F(s,u(s))\,ds.
\end{equation*}
在固定区间\([a,b]\)上，若\(F\)在\([a,b]\times\mathbb R\)连续，这就定义了从\(C[a,b]\)到自身的映射。在这一运算下保持不变的函数，即\(\mathcal Tu=u\)的{\heiti{}不动点}，正是初值问题的解。对于线性方程，\(F(t,u)=q(t)-p(t)u\)。以\(M_p\)记乘以\(p\)的运算，则
\begin{equation*}
 \mathcal Tu=y_0+I_{t_0}q-I_{t_0}M_pu,
 \qquad (I_{t_0}f)(t)=\int_{t_0}^t f(s)\,ds.
\end{equation*}
它由线性积分运算和固定函数的相加组成。Euler方法逐个更新离散的函数值；这里则从一个近似函数出发，更新整个函数。

\paragraph{从逐次积分到指数函数} \mbox{}

对\(y'=y\)、\(y(0)=1\)，在固定的\([0,T]\)上取\(u_0(t)=1\)，并递推定义
\begin{equation*}
 u_{n+1}(t)=1+\int_0^t u_n(s)\,ds.
\end{equation*}
于是\(u_1=1+t\)、\(u_2=1+t+t^2/2\)，由归纳法易证
\begin{equation*}
 u_n(t)=\sum_{j=0}^n\frac{t^j}{j!}.
\end{equation*}
这些函数恰是\(e^t\)在零点的Taylor多项式。由Lagrange余项，
\begin{equation*}
 \max_{0\le t\le T}|e^t-u_n(t)|\le\frac{e^TT^{n+1}}{(n+1)!}\longrightarrow0.
\end{equation*}
因此它们在整个区间上一致逼近\(e^t\)。积分误差满足
\begin{equation*}
 \left|\int_0^t(u_n(s)-e^s)\,ds\right|
 \le T\max_{0\le s\le T}|u_n(s)-e^s|,
\end{equation*}
可以在递推等式中取极限，得到\(e^t=1+\int_0^t e^s\,ds\)。有限次积分产生多项式，统一误差控制使它们的极限保留积分关系。

\paragraph{初值与唯一性} \mbox{}

若\(u,v\)都是上述积分方程的连续解，记\(\delta=\max_{[0,T]}|u-v|\)，则
\begin{equation*}
 |u(t)-v(t)|\le\int_0^t|u(s)-v(s)|\,ds.
\end{equation*}
先用\(|u-v|\le\delta\)得到\(|u(t)-v(t)|\le\delta t\)，再代回得到\(\delta t^2/2\)。由归纳法易证
\begin{equation*}
 |u(t)-v(t)|\le\frac{\delta t^n}{n!}\le\frac{\delta T^n}{n!}\qquad(n\ge1).
\end{equation*}
右端趋零，故\(u=v\)。这从积分关系证明了解的唯一性：初值相同，可能的差在反复累积估计下被压到零。

\begin{exercise}
\mbox{}\par\nobreak

（1）对\(y'=2ty\)、\(y(0)=1\)，从\(u_0=1\)出发作逐次积分。证明\(u_n(t)=\sum_{j=0}^nt^{2j}/j!\)，并在\([0,T]\)上给出它逼近精确解的统一误差界。

（2）若连续函数\(F\)满足\(|F(t,u)-F(t,v)|\le L|u-v|\)，仿照正文的反复积分估计，证明\([0,T]\)上具有相同初值的两个解相同。指出每次积分怎样产生新的阶乘因子。

\begin{samepage}
（3）对\(y'=2\sqrt{|y|}\)、\(y(0)=0\)，验证任意\(c\ge0\)都给出一个解：\(t\le c\)时取零，\(t>c\)时取\((t-c)^2\)。若从零函数开始逐次积分，会得到哪个解？比较本题的\(F\)与第（2）题的条件。
\end{samepage}
\end{exercise}
\endgroup

\subsection{Euler方法的误差估计*}
\label{subsec:calculus-euler-error}

Euler方法每一步产生的近似误差，会随递推继续传播。下面分别控制这两部分误差。设精确解在\([t_0,T]\)上二阶连续可导，\(|y''|\le M\)，且对该区间内的每个\(t\)以及任意实数\(u,v\)，有
\begin{equation*}
 |F(t,u)-F(t,v)|\le L|u-v|,\qquad L>0.
\end{equation*}
令\(h=(T-t_0)/N\)，用\(Y_n\)表示Euler近似值，\(E_n=|y(t_n)-Y_n|\)。Taylor公式表明每一步从精确值出发造成的误差不超过\(Mh^2/2\)，因而
\begin{equation*}
 E_{n+1}\le(1+Lh)E_n+\frac M2h^2,\qquad E_0=0.
\end{equation*}
反复代入并求有限等比和，得到
\begin{equation*}
 E_n\le\frac{Mh}{2L}\bigl((1+Lh)^n-1\bigr)
 \le\frac{Mh}{2L}\bigl(e^{L(T-t_0)}-1\bigr),\qquad 0\le n\le N.
\end{equation*}
固定区间后，误差至多为步长的一个常数倍。局部的二次误差经过许多步累积，成为整体的一次误差；控制增长因子，才能从局部近似得到整个区间上的结论。

\begin{example}
对\(y'=-y\)、\(y(0)=1\)，用上述估计确定一个步长，使\([0,1]\)上所有网格点的误差都小于\(10^{-2}\)。再由\(Y_n=(1-h)^n\)直接比较，判断一般误差界是否给出了必要的步长要求。
\end{example}

\begin{exercise}
对例\ref{ex:calculus-euler-values}，用Euler误差估计选择步长，再与实际误差比较。
\end{exercise}
\endgroup

\FloatBarrier
\newpage
\section{综合习题}
\label{sec:calculus-exercises}

以下题目覆盖全章，可根据问题选择已有的方法。星号题使用相同的函数与累积对象，进一步研究局部模型、统一误差及函数运算。

\begin{exercise}
\mbox{}\par\nobreak

（1）设\(f\)在\(a\)处可导。求\(\displaystyle\lim_{h\to0}\frac{f(a+3h)-f(a-2h)}h\)，并用微分解释各增量的作用。

（2）设\(f(0)=0\)，\(x\ne0\)时\(f(x)=x^2\sin(1/x)\)。判断它的连续性、可导性及导函数在零点的连续性。再与\(g(0)=0\)、\(g(x)=x\sin(1/x)\)比较，说明函数增量与差商分别需要怎样的控制。

（3）设\(f(x)=x^3+x\)，反函数为\(g\)。求\(g'(2)\)，再求\(\int_0^2g'(x)\,dx\)。第二问是否必须先求出\(g\)的显式表达式？

（4）讨论\(f_a(x)=e^x-ax\)在\(\mathbb R\)上的单调性、最小值及零点个数，并确定\(f_a(x)\ge0\)恒成立的参数范围。

（5）设\(f\)在\([0,1]\)上连续、内部可导，且\(f(0)=f(1)=0\)、\(|f'|\le2\)。估计\(|f(x)|\)以及\(\int_0^1|f(x)|\,dx\)，并构造达到积分上界的连续分段线性函数。比较这个函数与原题可导要求的关系。

（6）以2024年新课标I卷第10题的函数\(f(x)=(x-1)^2(x-4)\)为研究对象：求其单调区间与极值，讨论方程\(f(x)=k\)的不同实根个数，再求图像在\([1,4]\)上与横轴围成的面积。

（7）求\(y=e^x\)与\(y=-e^{-x}\)的公共切线。分别设两个切点，比较公共切线与公共点处相切的条件。

（8）设\(A(x)=\int_0^x(t^2-2t)\,dt\)，\(0\le x\le3\)。求其单调区间、最值与零点，并解释\(A\)的增减为何不由\(t^2-2t\)本身的单调性直接决定。

（9）设\(f\)在\([0,1]\)上连续，且\(f(x)+f(1-x)=1\)。求\(\int_0^1f(x)\,dx\)。\(\int_0^1xf(x)\,dx\)能否仅由这些条件唯一确定？举例或证明。

（10）计算\(\int_0^1 x/(1+x^2)\,dx\)、\(\int_0^{\pi/2}x\cos x\,dx\)和\(\int_{-1}^2|x^2-x|\,dx\)，说明各自选择换元、分部或区间拆分的依据。

（11）设\(f\)在\([0,1]\)上连续可导。证明\(\int_0^1f(x)\,dx=f(1)-\int_0^1xf'(x)\,dx\)。若\(f(0)=0\)、\(0\le f'\le1\)，估计\(\int_0^1 f\)。

（12）用四等分的中点与梯形公式估计\(\int_0^1dx/(1+x)\)。利用凸性判断上下界，再由\(\ln2\)核对。若只知道函数的单调性，能够得到哪些估计？

（13）细杆位于\([0,1]\)，线密度为\(\rho(x)=1+cx\)，其中\(c\ge0\)。求总质量和质心坐标，讨论质心怎样随\(c\)变化，并判断能否位于\(x=3/4\)。

（14）设\(f\)在\([0,1]\)上连续。证明存在\(\xi\in[0,1]\)，使\(\int_0^1xf(x)\,dx=f(\xi)/2\)。再说明把权函数\(x\)换成可能变号的函数时，结论为何未必成立。

（15）设\(f\)在\([0,1]\)上有界，且只有有限个间断点。它的累积函数是否一定连续？是否一定处处可导？分别说明依据并给出适当例子。

（16）分别构造：有无穷多个间断点但可积的函数；有界却不可积的函数；非零但绝对值积分为零的可积函数。比较这些例子与连续函数类的性质。

（17）（2025年北京卷第20题）已知函数\(f\)的定义域为\((-1,+\infty)\)，\(f(0)=0\)，且\(f'(x)=\ln(1+x)/(1+x)\)。设\(l_1\)是曲线在点\(A(a,f(a))\)、\(a\ne0\)处的切线。

\quad \exercisesubitem{1}求\(f'(x)\)的最大值；\exercisesubitem{2}当\(-1<a<0\)时，证明除切点外，曲线在\(l_1\)上方；\exercisesubitem{3}设过\(A\)且垂直于\(l_1\)的直线为\(l_2\)，\(l_1,l_2\)与横轴交点的横坐标分别为\(x_1,x_2\)。若\(a>0\)，求\((2a-x_1-x_2)/(x_2-x_1)\)的取值范围。

（18）（2022年新高考I卷第22题）已知函数\(f(x)=e^x-ax\)和\(g(x)=ax-\ln x\)有相同的最小值。

\quad \exercisesubitem{1}求\(a\)；\exercisesubitem{2}证明存在直线\(y=b\)，其与两条曲线\(y=f(x)\)、\(y=g(x)\)共有三个不同的交点，并且从左到右的三个交点的横坐标成等差数列。

（19）（2025年新课标II卷第18题）已知\(f(x)=\ln(1+x)-x+x^2/2-kx^3\)，其中\(0<k<1/3\)。

\quad \exercisesubitem{1}证明\(f\)在\((0,+\infty)\)内存在唯一的极值点和唯一的零点；\exercisesubitem{2}记它们分别为\(x_1,x_2\)，令\(g(t)=f(x_1+t)-f(x_1-t)\)。证明\(g\)在\((0,x_1)\)上单调递减，并比较\(2x_1\)与\(x_2\)的大小，证明你的结论。

（20）（2025年新课标I卷第19题）\exercisesubitem{1}求\(f(x)=5\cos x-\cos5x\)在\([0,\pi/4]\)上的最大值；\exercisesubitem{2}给定\(\theta\in(0,\pi)\)和\(a\in\mathbb R\)，证明存在\(y\in[a-\theta,a+\theta]\)，使\(\cos y\le\cos\theta\)；\exercisesubitem{3}设\(b\in\mathbb R\)，若存在\(\varphi\in\mathbb R\)，使\(5\cos x-\cos(5x+\varphi)\le b\)对所有实数\(x\)恒成立，求\(b\)的最小值。

\end{exercise}

\FloatBarrier
\newpage
\begin{exercise}[选做]
\mbox{}\par\nobreak

（1）设\(f(a)=a\)、\(f'(a)=\lambda\)，其中\(|\lambda|<1\)。由导数定义证明，存在\(q<1\)及\(a\)的一个邻域，使\(|f(x)-a|\le q|x-a|\)。据此研究从邻域内出发的迭代\(x_{n+1}=f(x_n)\)，给出收敛与误差估计，接回第一章的迭代数列。

（2）对\(f(x)=\ln(1+x)\)，用二阶Taylor公式估计\(\ln1.1\)，并给出误差上界。再用\(\int_1^{1.1}dt/t\)的中点公式估计同一个数，比较两种误差控制所需的信息。

（3）设\(f_n(x)=x^n\)，\(0\le x\le1\)。求逐点极限，证明不一致收敛，并计算\(\lim_n\int_0^1 f_n\)。这个例子能否说明“一致收敛是积分极限关系成立的必要条件”？

（4）设\(f_n\)是\([0,1]\)上的连续非负函数，\(\int_0^1f_n\to0\)。是否必有\(\|f_n\|_\infty\to0\)？尝试构造高度不变、宽度缩小的三角形图像，并解释面积误差与最大误差的区别。

（5）设\(F\in C^1[0,1]\)。证明\(F(x)=F(0)+\int_0^xF'(t)\,dt\)，并估计\(\|F\|_\infty\)。对\(F,G\)写出对应差值的估计，说明同时控制初值与导函数为什么能控制整个函数。

（6）用极坐标面积公式计算\(r=1+\cos\theta\)围成的面积，再把它写成平面参数曲线，用参数面积公式核对。说明两种局部贡献的表达为何不同。

（7）讨论\(\int_1^{+\infty}dx/[x(1+x^p)]\)的敛散性，其中\(p\in\mathbb R\)。收敛时给出尾部估计，再用积分判别法研究相应级数。

（8）设连续函数\(f_n,f\)定义在\([0,+\infty)\)，在每个有限区间上\(f_n\)一致收敛于\(f\)，且\(|f_n(x)|\le e^{-x}\)。证明\(f\)的反常积分绝对收敛，且\(\int_0^{+\infty}f_n\to\int_0^{+\infty}f\)。先固定截断位置控制有限范围，再统一估计尾部。

（9）设\(f\)在\([a,b]\)上有界、只有有限个间断点。用Darboux判据证明可积性，并说明“间断点附近缩小总长度”与“连续部分统一控制振幅”分别解决了什么问题。

（10）设\(f\)二阶连续可导，\(m\le f''\le M\)。用积分余项证明图像与任意切线之间的二次误差界，再由它推导中点数值积分在各小段上的误差界。
\end{exercise}

